% =====================================================================
%  Quantitative Methods for Urban and Regional Research
%  A Regional Policy Informatics Approach
%  Author: David S. Bieri
%  © 2011, 2014, 2026
% =====================================================================
\documentclass[12pt, a4paper, oneside, openany]{book}

%%%%%%%%%%%%%%%%%%%%%%%%%%%%%%%%%%%%%%%%%%%%%%%%%%%%%%%%%%%%%%%%%%%%%%%%%%
%%%%%  Put packages you want to use here and 'fancyhdr' is required   %%%%
%%%%%%%%%%%%%%%%%%%%%%%%%%%%%%%%%%%%%%%%%%%%%%%%%%%%%%%%%%%%%%%%%%%%%%%%%%
\usepackage[sectionbib]{natbib}
\usepackage{chapterbib}
\usepackage{thumbpdf}
\pdfcompresslevel=9
\usepackage{imakeidx}
\usepackage[justific=RaggedRight]{idxlayout}
\usepackage{subfigure,setspace,multirow,longtable,threeparttable,grffile,rotating}
\usepackage{pdflscape}
\usepackage{graphicx}
\usepackage[usenames,dvipsnames]{color}
\usepackage[bookmarks=true,citecolor=Blue,urlcolor=Blue,bookmarksopen=false,bookmarksnumbered=true,bookmarkstype=toc,colorlinks=true,linkcolor=blue]{hyperref}
\usepackage{amssymb} %% The amssymb package provides various useful mathematical symbols
%\usepackage{amsthm,amsmath} %% The amsthm package provides extended theorem environments
\usepackage{arydshln}
\usepackage{times}
\usepackage[usenames,dvipsnames]{pstricks}
\usepackage{pst-pdf}
\usepackage{epsfig,pifont}
\usepackage{pst-grad}
\usepackage{pst-plot}
\usepackage{geometry}
\usepackage{amsmath}
%\usepackage{multind}\ProvidesPackage{multind}
%% The lineno packages adds line numbers. Start line numbering with
%% \begin{linenumbers}, end it with \end{linenumbers}. Or switch it on
%% for the whole article with \linenumbers.
\usepackage{enumerate}
\usepackage{lineno}
\usepackage{fancyhdr}
\usepackage{arydshln}
\usepackage{epsfig}
\usepackage{cite}
\usepackage{amsmath}
\usepackage{theorem}
\usepackage{latexsym}
\usepackage{epic}
\usepackage{bbding}
\usepackage{appendix}
\usepackage{caption}
\captionsetup{labelfont=normalsize}
\usepackage{paralist}
\usepackage[mathcal]{euscript}
\usepackage{bm}
\usepackage{type1cm}
\usepackage{lettrine}
%\usepackage{librebaskerville}
\usepackage[latin1]{inputenc}
\usepackage{tikz}
\usetikzlibrary{shapes,arrows,matrix}
\usepackage{microtype}
%\usepackage[urw-garamond]{mathdesign}
\usepackage[T1]{fontenc}
\usepackage[scaled=0.85]{beramono}


%%%%%%%%%%%%%%%%%%%%%%%%%%%%%%%%%%%%%%%%%%%%%%%%%%%%%%%%%%%%%%%%%%%%%%%%%%
%%%%%                 Set line spacing = 1.5 here                   %%%%%%
%%%%%%%%%%%%%%%%%%%%%%%%%%%%%%%%%%%%%%%%%%%%%%%%%%%%%%%%%%%%%%%%%%%%%%%%%%
%\renewcommand{\baselinestretch}{1.5}
%%%%%%%%%%%%%%%%%%%%%%%%%%%%%%%%%%%%%%%%%%%%%%%%%%%%%%%%%%%%%%%%%%%%%%%%%%
%%%%%                      Your fancy heading                       %%%%%%
%%%%% For the final copy you need to remove '\bfseries\today' below %%%%%%
%%%%%%%%%%%%%%%%%%%%%%%%%%%%%%%%%%%%%%%%%%%%%%%%%%%%%%%%%%%%%%%%%%%%%%%%%%
%\pagestyle{fancy}
%\renewcommand{\chaptermark}[1]{\markright{\chaptername\ \thechapter.\ #1}}
%\renewcommand{\sectionmark}[1]{\markright{\thesection.\ #1}{}}
%\lhead[\fancyplain{}{}]%
%      {\fancyplain{}{\slshape\rightmark}}
%\chead[\fancyplain{}{}]%
%      {\fancyplain{}{}}
%\rhead[\fancyplain{}{}]%
%      {\fancyplain{}{\bfseries\thepage}}
%\lfoot[\fancyplain{}{}]%
%      {\fancyplain{}{}}
%\cfoot[\fancyplain{}{}]%
%      {\fancyplain{}{}}
%\rfoot[\fancyplain{}{}]%
%      {\fancyplain{}{}}
%%%%%%%%%%%%%%%%%%%%%%%%%%%%%%%%%%%%%%%%%%%%%%%%%%%%%%%%%%%%%%%%%%%%%%%%%%
%%%%%%%%%%%%Here you set the space between the main text%%%%%%%%%%%%%%%%%%
%%%%%%%%%%%%%%%%%%%and the start of the footnote%%%%%%%%%%%%%%%%%%%%%%%%%%
%%%%%%%%%%%%%%%%%%%%%%%%%%%%%%%%%%%%%%%%%%%%%%%%%%%%%%%%%%%%%%%%%%%%%%%%%%
\addtolength{\skip\footins}{5mm}
%%%%%%%%%%%%%%%%%%%%%%%%%%%%%%%%%%%%%%%%%%%%%%%%%%%%%%%%%%%%%%%%%%%%%%%%%%
%%%%%      Define new counter so you can have the equation           %%%%%
%%%%%    number 4.2.1a for example, this a gift from J.F.Blowey      %%%%%
%%%%%%%%%%%%%%%%%%%%%%%%%%%%%%%%%%%%%%%%%%%%%%%%%%%%%%%%%%%%%%%%%%%%%%%%%%
\newcounter{ind}
\def\eqlabon{
\setcounter{ind}{\value{equation}}\addtocounter{ind}{1}
\setcounter{equation}{0}
\renewcommand{\theequation}{\arabic{chapter}%
         .\arabic{ind}\alph{equation}}}
\def\eqlaboff{
\renewcommand{\theequation}{\arabic{chapter}%
         .\arabic{equation}}
\setcounter{equation}{\value{ind}}}

%%%%%%%%%%%%%%%%%%%%%%%%%%%%%%%%%%%%%%%%%%%%%%%%%%%%%%%%%%%%%%%%%%%%%%%%%%
%%%%%%%%%%%%           New theorem you want to use              %%%%%%%%%%
%%%%%%%%%%%%%%%%%%%%%%%%%%%%%%%%%%%%%%%%%%%%%%%%%%%%%%%%%%%%%%%%%%%%%%%%%%
%{\theorembodyfont{\rmfamily}\newtheorem{Pro}{{\textbf Proposition}}[section]}
%{\theorembodyfont{\rmfamily}\newtheorem{The}{{\textbf Theorem}}[section]}
%{\theorembodyfont{\rmfamily}\newtheorem{Def}[The]{{\textbf Definition}}}
%{\theorembodyfont{\rmfamily}\newtheorem{Cor}[The]{{\textbf Corollary}}}
%{\theorembodyfont{\rmfamily}\newtheorem{Lem}[The]{{\textbf Lemma}}}
%{\theorembodyfont{\rmfamily}\newtheorem{Exp}{{\textbf Example}}[section]}
%\def\remark{\textbf{Remark}:}
%\def\remarks{\textbf{Remarks}:}
%\def\bproof{\textbf{Proof}: }
%\def\eproof{\hfill$\Box$}
%%%%%%%%%%%%%%%%%%%%%%%%%%%%%%%%%%%%%%%%%%%%%%%%%%%%%%%%%%%%%%%%%%%%%%%%%%
%%%%%%%    Bold font in math mode, a gift from J.F.Blowey       %%%%%%%%%%
%%%%%%%%%%%%%%%%%%%%%%%%%%%%%%%%%%%%%%%%%%%%%%%%%%%%%%%%%%%%%%%%%%%%%%%%%%
\def\bv#1{\mbox{\boldmath$#1$}}
%%%%%%%%%%%%%%%%%%%%%%%%%%%%%%%%%%%%%%%%%%%%%%%%%%%%%%%%%%%%%%%%%%%%%%%%%%
%%%%%%%        New symbol which is not defined in Latex         %%%%%%%%%%
%%%%%%%                 a gift from J.F.Blowey                  %%%%%%%%%%
%%%%%%%%%%%%%%%%%%%%%%%%%%%%%%%%%%%%%%%%%%%%%%%%%%%%%%%%%%%%%%%%%%%%%%%%%%
% The Mean INTegral
% to be used in displaystyle
\def\mint{\textstyle\mints\displaystyle}
% to be used in textstyle
\def\mints{\int\!\!\!\!\!\!{\rm-}\ }
%
% The Mean SUM
% to be used in displaystyle
\def\msum{\textstyle\msums\displaystyle}
% to be used in textstyle
\def\msums{\sum\!\!\!\!\!\!\!{\rm-}\ }
%%%%%%%%%%%%%%%%%%%%%%%%%%%%%%%%%%%%%%%%%%%%%%%%%%%%%%%%%%%%%%%%%%%%%%%%%%
%%%%%%%%%%            Define your short cut here              %%%%%%%%%%%%
%%%%%%%%%%%%%%%%%%%%%%%%%%%%%%%%%%%%%%%%%%%%%%%%%%%%%%%%%%%%%%%%%%%%%%%%%%
\def\poincare{Poincar\'e }
\def\holder{H\"older }


\def\citeapos#1{\citeauthor{#1}'s (\citeyear{#1})}

\def\newblock{\hskip .11em plus .33em minus .07em}


\newcommand{\HRule}{\rule{\linewidth}{0.2mm}}

\newcommand\independent{\protect\mathpalette{\protect\independenT}{\perp}}
\def\independenT#1#2{\mathrel{\rlap{$#1#2$}\mkern2mu{#1#2}}} %\newcommand\indepent is used to make the double perp. "indpendent given" symbol


\long\def\symbolfootnote[#1]#2{\begingroup%
\def\thefootnote{\fnsymbol{footnote}}\footnote[#1]{#2}\endgroup}

\def\blfootnote{\xdef\@thefnmark{}\@footnotetext} 

\bibliographystyle{econometrica}

\usepackage{listings}

\definecolor{dkgreen}{rgb}{0,0.6,0}
\definecolor{gray}{rgb}{0.4,0.4,0.4}
\definecolor{lightgray}{rgb}{0.9,0.9,0.9}
\definecolor{mauve}{rgb}{0.58,0,0.82}


\newcommand\indexkeywords[1]{\index{keywords, #1}}
%\newcommand\lstindexmacro[1]{\index{{\ttfamily#1}}}
\newcommand{\indexfonction}[1]{\index{#1@\texttt{#1}}}

\lstloadlanguages{R}
\lstdefinelanguage{Renhanced}%
  {keywords={abbreviate,abline,abs,acos,acosh,action,add1,add,%
      aggregate,alias,Alias,alist,all,anova,any,aov,aperm,append,apply,%
      approx,approxfun,apropos,Arg,args,array,arrows,as,asin,asinh,%
      atan,atan2,atanh,attach,attr,attributes,autoload,autoloader,ave,%
      axis,backsolve,barplot,basename,besselI,besselJ,besselK,besselY,%
      beta,binomial,body,box,boxplot,break,browser,bug,builtins,bxp,by,%
      c,C,call,Call,case,cat,category,cbind,ceiling,character,char,%
      charmatch,check,chol,chol2inv,choose,chull,class,close,cm,codes,%
      coef,coefficients,co,col,colnames,colors,colours,commandArgs,%
      comment,complete,complex,conflicts,Conj,contents,contour,%
      contrasts,contr,control,helmert,contrib,convolve,cooks,coords,%
      distance,coplot,cor,cos,cosh,count,fields,cov,covratio,wt,CRAN,%
      create,crossprod,cummax,cummin,cumprod,cumsum,curve,cut,cycle,D,%
      data,dataentry,date,dbeta,dbinom,dcauchy,dchisq,de,debug,%
      debugger,Defunct,default,delay,delete,deltat,demo,de,density,%
      deparse,dependencies,Deprecated,deriv,description,detach,%
      dev2bitmap,dev,cur,deviance,off,prev,,dexp,df,dfbetas,dffits,%
      dgamma,dgeom,dget,dhyper,diag,diff,digamma,dim,dimnames,dir,%
      dirname,dlnorm,dlogis,dnbinom,dnchisq,dnorm,do,dotplot,double,%
      download,dpois,dput,drop,drop1,dsignrank,dt,dummy,dump,dunif,%
      duplicated,dweibull,dwilcox,dyn,edit,eff,effects,eigen,else,%
      emacs,end,environment,env,erase,eval,equal,evalq,example,exists,%
      exit,exp,expand,expression,External,extract,extractAIC,factor,%
      fail,family,fft,file,filled,find,fitted,fivenum,fix,floor,for,%
      For,formals,format,formatC,formula,Fortran,forwardsolve,frame,%
      frequency,ftable,ftable2table,function,gamma,Gamma,gammaCody,%
      gaussian,gc,gcinfo,gctorture,get,getenv,geterrmessage,getOption,%
      getwd,gl,glm,globalenv,gnome,GNOME,graphics,gray,grep,grey,grid,%
      gsub,hasTsp,hat,heat,help,hist,home,hsv,httpclient,I,identify,if,%
      ifelse,Im,image,\%in\%,index,influence,measures,inherits,install,%
      installed,integer,interaction,interactive,Internal,intersect,%
      inverse,invisible,IQR,is,jitter,kappa,kronecker,labels,lapply,%
      layout,lbeta,lchoose,lcm,legend,length,levels,lgamma,library,%
      licence,license,lines,list,lm,load,local,locator,log,log10,log1p,%
      log2,logical,loglin,lower,lowess,ls,lsfit,lsf,ls,machine,Machine,%
      mad,mahalanobis,main,make,link,margin,match,Math,matlines,mat,matplot,%
      matpoints,matrix,max,mean,median,memory,menu,merge,methods,min,%
      missing,Mod,mode,model,response,mosaicplot,mtext,mvfft,na,nan,%
      names,omit,nargs,nchar,ncol,NCOL,new,next,NextMethod,nextn,%
      nlevels,nlm,noquote,NotYetImplemented,NotYetUsed,nrow,NROW,null,%
      numeric,\%o\%,objects,offset,old,on,Ops,optim,optimise,optimize,%
      options,or,order,ordered,outer,package,packages,page,pairlist,%
      pairs,palette,panel,par,parent,parse,paste,path,pbeta,pbinom,%
      pcauchy,pchisq,pentagamma,persp,pexp,pf,pgamma,pgeom,phyper,pico,%
      pictex,piechart,Platform,plnorm,plogis,plot,pmatch,pmax,pmin,%
      pnbinom,pnchisq,pnorm,points,poisson,poly,polygon,polyroot,pos,%
      postscript,power,ppoints,ppois,predict,preplot,pretty,Primitive,%
      print,prmatrix,proc,prod,profile,proj,prompt,prop,provide,%
      psignrank,ps,pt,ptukey,punif,pweibull,pwilcox,q,qbeta,qbinom,%
      qcauchy,qchisq,qexp,qf,qgamma,qgeom,qhyper,qlnorm,qlogis,qnbinom,%
      qnchisq,qnorm,qpois,qqline,qqnorm,qqplot,qr,Q,qty,qy,qsignrank,%
      qt,qtukey,quantile,quasi,quit,qunif,quote,qweibull,qwilcox,%
      rainbow,range,rank,rbeta,rbind,rbinom,rcauchy,rchisq,Re,read,csv,%
      csv2,fwf,readline,socket,real,Recall,rect,reformulate,regexpr,%
      relevel,remove,rep,repeat,replace,replications,report,require,%
      resid,residuals,restart,return,rev,rexp,rf,rgamma,rgb,rgeom,R,%
      rhyper,rle,rlnorm,rlogis,rm,rnbinom,RNGkind,rnorm,round,row,%
      rownames,rowsum,rpois,rsignrank,rstandard,rstudent,rt,rug,runif,%
      rweibull,rwilcox,sample,sapply,save,scale,scan,scan,screen,sd,se,%
      search,searchpaths,segments,seq,sequence,setdiff,setequal,set,%
      setwd,show,sign,signif,sin,single,sinh,sink,solve,sort,source,%
      spline,splinefun,split,sqrt,stars,start,stat,stem,step,stop,%
      storage,strstrheight,stripplot,strsplit,structure,strwidth,sub,%
      subset,substitute,substr,substring,sum,summary,sunflowerplot,svd,%
      sweep,switch,symbol,symbols,symnum,sys,status,system,t,table,%
      tabulate,tan,tanh,tapply,tempfile,terms,terrain,tetragamma,text,%
      time,title,topo,trace,traceback,transform,tri,trigamma,trunc,try,%
      ts,tsp,typeof,unclass,undebug,undoc,union,unique,uniroot,unix,%
      unlink,unlist,unname,untrace,update,upper,url,UseMethod,var,%
      variable,vector,Version,vi,warning,warnings,weighted,weights,%
      which,while,window,write,\%x\%,x11,X11,xedit,xemacs,xinch,xlim,xor,%
      xpdrows,xy,xyinch,yinch,ylim,zapsmall,zip},%
   otherkeywords={!,!=,~,$,*,\%,\&,\%/\%,\%*\%,\%\%,<-,<<-,_,/},%
   alsoother={._$},%
   sensitive,%
   morecomment=[l]\#,%
   morestring=[d]",%
   morestring=[d]'% 2001 Robert Denham
  }


\lstset{ %
  backgroundcolor=\color{white},   % choose the background color; you must add \usepackage{color} or \usepackage{xcolor}
  basicstyle=\scriptsize\ttfamily,        % the size of the fonts that are used for the code
  breakatwhitespace=false,         % sets if automatic breaks should only happen at whitespace
  breakindent=25pt,
  breaklines=true,                 % sets automatic line breaking
  captionpos=b,                    % sets the caption
  commentstyle=\color{dkgreen}\textsl,    % comment style
  deletekeywords={...},            % if you want to delete keywords from the given language
  escapeinside={\%@}{@)},          % if you want to add LaTeX within your code
  extendedchars=true,              % lets you use non
  frame=none,                      % adds a frame around the code
  index=[1][keywords],
%indexstyle=[1]\lstindexmacro}
  indexstyle=\indexfonction,
  keepspaces=true,                 % keeps spaces in text, useful for keeping indentation of code (possibly needs columns=flexible)
  keywordstyle=\color{blue},       % keyword style
  language=Renhanced,              % the language of the code
  morekeywords={*,...},            % if you want to add more keywords to the set
  numbers=left,                    % where to put the line
  numbersep=5pt,                   % how far the line
  numberstyle=\tiny\color{gray},   % the style that is used for the line
  rulecolor=\color{black},         % if not set, the frame
  showspaces=false,                % show spaces everywhere adding particular underscores; it overrides 'showstringspaces'
  showstringspaces=false,          % underline spaces within strings only
  showtabs=false,                  % show tabs within strings adding particular underscores
  stepnumber=1,                    % the step between two line
  stringstyle=\color{mauve},       % string literal style
  tabsize=2,                       % sets default tabsize to 2 spaces
  title=\lstname                   % show the filename of files included with \lstinputlisting; also try caption instead of title
}


\title{\huge\textbf{Quantitative Methods for\\Urban and Regional Research}\\[10pt]
       \Large\textit{A Regional Policy Informatics Approach}}
\author{\Large David S.\ Bieri\\[4pt]
        \normalsize University of Michigan\\[6pt]
        \footnotesize\copyright~2011, 2014, 2026\\[2pt]
        \scriptsize This manuscript may be printed and reproduced for individual
        or instructional use, but may not be printed for commercial purposes.
        Licensed under a Creative Commons
        \href{http://creativecommons.org/licenses/by-nc-sa/3.0/deed.en_US}{%
        Attribution-NonCommercial-ShareAlike 3.0 Unported License}.}
\date{This version: \today}

\makeindex

% =====================================================================
\begin{document}

\maketitle
\frontmatter
\setcounter{tocdepth}{2}
\setcounter{page}{3}

% ── Acknowledgements ──────────────────────────────────────────────────
\chapter*{Acknowledgements}
This manuscript has evolved from a first-year graduate course in quantitative
methods for urban planning students. Several cohorts of students at the
University of Michigan and at Virginia Tech have provided helpful feedback and
fairly patiently endured excursions into matrix algebra, occasionally agreeing
with me that it was important to derive certain concepts from first principles.
James Carpenter, Miles Hamby, Arthur Prokosch, and Cheng Wang provided
excellent research assistance and diligently transcribed a large part of my
iteratively amended lecture notes into \LaTeX{}. This version of the manuscript
has been substantially expanded and restructured to reflect the emergence of a
regional policy informatics paradigm and the maturation of open-source spatial
analytics workflows.

% ── Table of Contents, Figures, Tables ───────────────────────────────
\tableofcontents
\listoffigures
\listoftables

% ── Main matter ───────────────────────────────────────────────────────
\mainmatter

% ── Preface ───────────────────────────────────────────────────────────
% =====================================================================
%  Chapter 0: Preface — Caveat Computor
%  Source: Intro.tex (Overleaf) — adapted and expanded
% =====================================================================

\chapter*{Preface: \textit{Caveat Computor}}\label{chap_Preface}
\addcontentsline{toc}{chapter}{Preface: \textit{Caveat Computor}}

\section*{What this book is about}

The primary objective of this manuscript is to introduce students to the
quantitative methods and techniques used in urban and regional research. The
emphasis is on using methods in the context of planning, urban policy, and
regional science research: matching the method to the problem, developing an
interactive engagement with data, promoting creative and accessible styles of
data presentation, and cultivating a critical literacy of methods and their
limitations.

In its intellectual lineage, this text is unambiguously in the spirit of
Walter \citeapos{Isard1960a} classic \emph{Methods of Regional Analysis}---a
methodological manifesto for regional science as a scientific project, if ever
there was one.\footnote{Now in its seventh edition, \emph{Methods of
Interregional and Regional Analysis} \citep{Isard1998} remains the
methodological benchmark for regional science, covering regional and
interregional input-output analysis, econometrics, programming and industrial
and urban complex analysis, gravity and spatial interaction models, and applied
general interregional equilibrium models. In contrast to the largely
theoretical treatment in that volume, this manual focuses on the empirical
aspects of application, with special emphasis on the connections between theory
and evidence.} Of all the subdisciplines subsumed under the heading of
regional science, urban and regional planning is perhaps the archetypical
interdisciplinary venture---yet it still relies on solid disciplinary
skills. Planning educators regularly define an accomplished planner as
simultaneously ``the best geographer, the best economist, the best
environmental scientist, the best political scientist, and the best
architect'' \citep{Fainstein2011}.

\section*{The computational abundance paradox}

As computational costs have fallen precipitously and data availability has
exploded, one might expect the quality of empirical analysis to have risen in
proportion. The opposite tendency is equally plausible---and, this author
believes, equally common. Software and analytical tools are becoming ever more
user-friendly, and hitherto prohibitively complex analyses now require only a
few clicks deployed via an attractive graphical user interface. Yet perhaps
paradoxically, the more we rely on technology and the more accessible data
tools become, the less time is spent on careful and critical analysis. This
leads to a latent tendency to produce more empirical output rather than better
analysis. Just because we can (cheaply and conveniently) compute something does
not mean that we should. In other words: \emph{caveat computor}.

This paradox has only deepened with the emergence of large language model
(LLM) coding assistants capable of generating spatial regression code from a
single natural-language prompt. When generating a spatial lag model requires
less effort than reading this sentence, the limiting constraint on analysis
quality shifts entirely to whether the analyst understands what the model is
doing---whether the weight matrix choice is defensible, whether the
identification strategy is credible, whether the coefficient can be given a
causal interpretation. The theoretical rigour of Parts~I and~II of this
book is \emph{justified by}, not in tension with, the ease of modern code
generation. The most dangerous analyst is one who can produce polished output
without understanding its assumptions.

\section*{A new paradigm: regional policy informatics}

This book proposes a synthesis that takes the credibility revolution in
econometrics seriously, treats spatial structure as an inferential challenge
rather than a nuisance parameter, and integrates the quantitative traditions
of urban economics, regional science, economic geography, and real estate
around the unifying framework of spatial equilibrium.

\noindent\textbf{Central claim.} Spatial analysis in the social sciences has
systematically produced \emph{geographically-correlated} answers rather than
\emph{causally-identified} ones. Geographic confounding---the systematic
confusion of spatial correlation with causal identification---is to regional
analytics what omitted variable bias is to cross-sectional econometrics. The
response is a credibility-first, spatially-literate workflow built entirely on
open-source tools and public data infrastructure.

We define \textbf{regional policy informatics} as the systematic application
of computational, statistical, and causal inference methods to the analysis of
spatially-structured policy problems, using reproducible workflows and open
data infrastructure. What distinguishes this from ``spatial data science''
(the term used by the Python-oriented GIS community) is the explicit
commitment to policy-relevant causal identification. Spatial data science
optimises for pattern recognition and visualisation---it is largely descriptive
and predictive. Regional policy informatics is evaluative: it asks whether
interventions work, for whom, and why.

\section*{Four questions that frame good quantitative research}

The analytical philosophy underlying this text is organised around four
questions that \citet{Angrist2009a} argue should frame any credible empirical
project:

\begin{enumerate}
  \item What is the \emph{causal relationship of interest}? This question
    gives rise to the importance of the counterfactual.
  \item What \emph{ideal experiment} would capture this causal effect?
  \item What is the \emph{identification strategy}? How is the observational
    data used to approximate the ideal experiment?
  \item What is the \emph{mode of inference}? (microdata, clustered groups,
    spatial units, etc.)
\end{enumerate}

These four questions are not only the organising principle of Part~II
(The Linear Model and Causal Inference); they reappear throughout the book
as a touchstone for evaluating any spatial or regional analysis.

\section*{The open-source commitment}

This text is built around three principles:

\begin{enumerate}
  \item \textbf{Open source software} achieves interoperability and
    reproducibility without licensing barriers. The full analytical toolkit---
    from spatial regression to input-output modelling---is implemented in R
    and Python using freely available packages.

  \item \textbf{Public data infrastructure} supports a full regional
    analytics workflow without proprietary data. The Census, ACS, FRED,
    LEHD, BEA Regional Accounts, HUD, and OpenStreetMap collectively provide
    richer inputs than most proprietary alternatives.

  \item \textbf{Replicable research} is a scientific norm, not a technical
    nicety. The \emph{American Economic Review} requires that all data and
    code sufficient for replication be deposited prior to publication
    \citep{AER2024}. Urban planning and geography journals have been slower to
    adopt such standards; this text argues that they should not be.
\end{enumerate}

\section*{Prerequisites and how to use this book}

The material is aimed at first-year graduate students (or advanced
undergraduates) who have taken at minimum an introductory statistics course
covering the basics of probability, confidence intervals, hypothesis testing,
and bivariate regression. Exposure to linear algebra is helpful but not
required---Part~I provides everything needed from scratch.

Chapter~1 reviews probability and statistics. Chapter~2 introduces the matrix
algebra toolkit that underlies all subsequent chapters. Chapter~3 surveys data
sources and measurement. These three chapters constitute Part~I and can be
read selectively by students with stronger preparation. The heart of the book
is Parts~II and~III: Chapters~4--6 develop the linear model and causal
inference toolkit; Chapters~7--9 apply that toolkit to the spatial economy.
Part~IV covers time series, panel data, and the frontiers of computational
regional analysis.

\renewcommand\bibname{Further reading}
\begin{subappendices}
\bibliographystyle{econometrica}
\bibliography{Literature/OverLord}
\IfFileExists{Literature/Preface.bbl}{\input{Literature/Preface.bbl}}{\typeout{Need to run bibtex on Preface}}
\end{subappendices}


% ── Part I: Foundations ───────────────────────────────────────────────
\part{Foundations}\label{part1}
% =====================================================================
%  Chapter 1: Probability, Statistics, and the Language of Uncertainty
%  Source: Review.tex (Overleaf) + P1 transcription
% =====================================================================

\chapter{Probability, Statistics, and the Language of Uncertainty}
\label{chap_Review}

\section{A Paradigm Shift: From Classical to Frequentist to Bayesian}
\label{sec_paradigm}

``Probability approach in econometrics'' by Trygve \citet{Haavelmo1944} stands
at the beginning of a paradigm shift in statistical thinking; the
\emph{classical} interpretation of statistics is giving way to the
\emph{frequentist} interpretation. Under the classical interpretation,
probabilities are linked to the principle of indifference and statistical
events are governed by natural symmetry (e.g.\ coin tosses, die rolls).

Under the frequentist interpretation, probabilities of events are linked to
the frequency with which a given event is observed:
\begin{equation}
  P(x) = \lim_{n_t \to \infty} \frac{n_x}{n_t},
\end{equation}
where $n_x$ is the number of times event $x$ is observed and $n_t$ is the
total number of trials.

This probabilistic approach to statistics, which encompasses Bayesian
statistics\index{Bayesian statistics}, represents the world in terms of
``degrees of belief'' (Bayesian priors). More formally, \textbf{Bayes'
Theorem}\index{Bayes' Theorem} is expressed as the conditional probability of
event $A$ given $B$:

\begin{tcolorbox}[keyresult]
\begin{equation}
  P(A \mid B) = \frac{P(B \mid A)\,P(A)}{P(B)}
\end{equation}
\end{tcolorbox}


\section{Random Variables: Types and Typology}
\label{sec_RVtypes}

A random variable, usually written as $X$, is a variable whose possible values
are numerical outcomes of a random phenomenon. A random variable $X$ is a
(real) number with the following characteristics:

\begin{enumerate}[label=\alph*)]
  \item The value of $X$ is unknown.
  \item The possible values $X$ can take are known; the set of such values is
    called the \emph{support}, denoted $\boldsymbol{\mathcal{S}}_X$.
  \item If $\boldsymbol{\mathcal{A}} \subseteq \mathbb{R}$, we know the
    probability that $X$ belongs to $\boldsymbol{\mathcal{A}}$, denoted
    $P(X \in \boldsymbol{\mathcal{A}})$.
  \item If the value $x$ of $X$ becomes known such that $x=X$, then $x$ is
    called a \emph{realisation} of $X$.
\end{enumerate}

Random variables fall into two categories:

\begin{itemize}
  \item \textbf{Discrete random variables}\index{random variable!discrete}:
    Take only a countable number of distinct values (0, 1, 2, \ldots).
    Usually (but not necessarily) counts.
  \item \textbf{Continuous random variables}\index{random variable!continuous}:
    Take an infinite number of possible values. Usually measurements (height,
    weight, time).
\end{itemize}

We can also distinguish \emph{qualitative} (categorical) from
\emph{quantitative} variables:

\begin{itemize}
  \item \textbf{Nominal variables}\index{random variable!nominal}: Mutually
    exclusive, unordered categories (e.g.\ suburb type).
  \item \textbf{Ordinal variables}\index{random variable!ordinal}: Ordered
    categories where the differences between values are not meaningful (e.g.\
    quality-of-life rankings on a 1--10 scale).
  \item \textbf{Interval variables}\index{random variable!interval}:
    Differences are meaningful but there is no true zero (e.g.\ temperature
    in Celsius).
  \item \textbf{Ratio variables}\index{random variable!ratio}: A true zero
    exists and ratios are meaningful (e.g.\ height, weight, income).
\end{itemize}

\begin{table}[htp]
\centering
\caption{Basic statistical operations by variable type}
\label{tbl_Vars}
\renewcommand{\arraystretch}{1.3}
\begin{tabular}{@{}lcccc@{}}
\toprule
& \multicolumn{4}{c}{\emph{Variable type}} \\
\cmidrule{2-5}
\emph{Operation} & Nominal & Ordinal & Interval & Ratio \\
\midrule
Frequency distribution    & Yes & Yes & Yes & Yes \\
Median and percentiles    & No  & Yes & Yes & Yes \\
Addition / subtraction    & No  & No  & Yes & Yes \\
Mean, s.d., s.e.          & No  & No  & Yes & Yes \\
Ratio / coefficient of variation & No & No & No & Yes \\
\bottomrule
\end{tabular}
\renewcommand{\arraystretch}{1}
\end{table}


\section{Expectation and Variance}
\label{sec_Expectation}

\subsection{Expectation of a Discrete Random Variable}

\begin{equation}
  \E(X) = x_1 p_1 + x_2 p_2 + \cdots + x_n p_n = \sum_{i=1}^{n} x_i p_i
\end{equation}

In a finite sample, $\E(X)$ denotes the sample mean, and the sum of all
probabilities must equal one: $\sum_i p_i = 1$.

\paragraph{Example: The St.\ Petersburg Paradox}\index{St. Petersburg paradox}

Consider a gamble where an unbiased coin is tossed until tails ($T$) appears.
If $T$ shows up on the $n$-th toss, the payoff is $\$2^{n-1}$. How much would
you be willing to wager to participate?

\begin{align*}
  \E(X) &= \$2^0 \cdot \tfrac{1}{2} + \$2^1 \cdot \tfrac{1}{4}
           + \$2^2 \cdot \tfrac{1}{8} + \cdots \\
        &= \$1\cdot\tfrac{1}{2} + \$2\cdot\tfrac{1}{4}
           + \$4\cdot\tfrac{1}{8} + \cdots
         = \sum_{i=1}^{\infty} \$0.50 = \$\infty
\end{align*}

This is the St.~Petersburg Paradox, famously described by Daniel
Bernoulli\index{Bernoulli, Daniel} (1700--1782) to illustrate the failure of
the expected value principle in decision theory under uncertainty. The expected
payout is infinite, yet willingness to pay (WTP)\index{willingness to pay}
must be much less. The classical resolution introduces an explicit utility
function and the presumption of diminishing marginal utility, implying
$\E[U] < \E[X]$.

\paragraph{Example: Expected value of a pair of dice}

\begin{equation*}
  \E(X) = \tfrac{1}{6}\cdot1 + \tfrac{1}{6}\cdot2 + \cdots
          + \tfrac{1}{6}\cdot6 = 3.5
\end{equation*}

For a pair of dice, $\mu_X = 7$. Bivariate probabilities are shown in
Table~\ref{tbl_Dice}.

\begin{table}[htp]
\centering
\caption{Bivariate probabilities for a pair of dice}
\label{tbl_Dice}
\renewcommand{\arraystretch}{1.3}
\begin{tabular}{@{}lll|lll@{}}
\toprule
$X$ & $p$ & $Xp$ & $X$ & $p$ & $Xp$ \\
\midrule
$x_1$ & $p_1$ & $x_1p_1$ & 2  & 1/36 & 2/36  \\
$x_2$ & $p_2$ & $x_2p_2$ & 3  & 2/36 & 6/36  \\
      &       &           & 4  & 3/36 & 12/36 \\
      &       &           & 5  & 4/36 & 20/36 \\
      &       &           & 6  & 5/36 & 30/36 \\
      &       &           & 7  & 6/36 & 42/36 \\
      &       &           & 8  & 5/36 & 40/36 \\
      &       &           & 9  & 4/36 & 36/36 \\
      &       &           & 10 & 3/36 & 30/36 \\
      &       &           & 11 & 2/36 & 22/36 \\
$x_n$ & $p_n$ & $x_np_n$ & 12 & 1/36 & 6/36  \\
\midrule
\multicolumn{3}{l}{$\E(X)=\sum x_i p_i$} &&
\multicolumn{2}{r}{252/36 = 7} \\
\bottomrule
\end{tabular}
\renewcommand{\arraystretch}{1}
\end{table}

\subsection{Expectation Rules}

\begin{enumerate}
  \item $\E[X + Y + Z] = \E[X] + \E[Y] + \E[Z]$
  \item $\E[bX] = b\E[X]$
  \item $\E[b] = b$
  \item $\E[a + bX] = a + b\E[X]$
\end{enumerate}

\paragraph{Example} Let $V \sim \mathcal{N}(0,1)$ and $Z \equiv \mu + \sigma V$.
Then $\E[Z] = \mu$ and $\Var[Z] = \sigma^2$, since $\E[V]=0$ and
$\E[V^2]=1$.

\subsection{Variance}

\begin{equation}
  \Var(X) = \sigma^2_X = \E\bigl[(X-\mu_X)^2\bigr]
  = \sum_{i=1}^{n}(x_i - \mu_X)^2 p_i
\end{equation}

Alternative expression (computing formula):
\begin{align}
  \Var(X) &= \E\bigl[X^2 - 2\mu_X X + \mu_X^2\bigr] \nonumber \\
          &= \E[X^2] - 2\mu_X \E[X] + \mu_X^2 \nonumber \\
          &= \E[X^2] - \mu_X^2
\end{align}

\subsection{Fixed and Random Decomposition of Variables}

Let the random variable $X$ be decomposed into a constant mean $\mu_X$ and a
stochastic disturbance term $u$:
\begin{equation}
  X = \mu_X + u, \qquad \E[u] = \E[X - \mu_X] = 0
\end{equation}
Since all variation in $X$ is due to $u$, the variance of $X$ equals the
variance of the disturbance:
\begin{equation}
  \sigma^2_X = \E[(X-\mu_X)^2] = \E[u^2] = \sigma^2_u
\end{equation}
This decomposition into a \emph{systematic} (fixed) component and a
\emph{stochastic} (random) component is the conceptual foundation of every
regression model in this book.


\section{Continuous Random Variables}
\label{sec_ContRV}

All random variables have a \textbf{probability density function}
(PDF)\index{density function!probability} and a \textbf{cumulative density
function} (CDF)\index{density function!cumulative}:
\begin{equation}
  \Pr[A \leq x \leq B] = \int_A^B f(x)\,dx, \qquad
  F(x) = \int_{-\infty}^{x} f(t)\,dt = F(B) - F(A)
\end{equation}

The expectation and variance of a continuous random variable:
\begin{equation}
  \E[X] = \int_{-\infty}^{\infty} x\,f(x)\,dx = \mu_x,
  \qquad
  \Var(X) = \int_{-\infty}^{\infty}(x - \mu_x)^2 f(x)\,dx = \sigma_x^2
\end{equation}

\begin{figure}[htp]
\centering
\begin{tikzpicture}[>=Stealth, scale=1.1]
  % Left panel: PDF
  \begin{scope}[xshift=0cm]
    \draw[->] (-2.5,0) -- (2.7,0) node[right] {$x$};
    \draw[->] (0,-0.1) -- (0,1.6) node[above] {$f(x)$};
    \draw[very thick, domain=-2.4:2.4, samples=80]
      plot (\x, {1.3*exp(-0.8*\x*\x)});
    % Shaded region between -1.5 and -0.5
    \fill[gray!30, domain=-1.5:-0.5, samples=40]
      (-1.5,0) -- plot (\x, {1.3*exp(-0.8*\x*\x)}) -- (-0.5,0) -- cycle;
    \draw[dashed] (-1.5,0) node[below]{\small$A$}
                         -- (-1.5, {1.3*exp(-0.8*1.5*1.5)});
    \draw[dashed] (-0.5,0) node[below]{\small$B$}
                         -- (-0.5, {1.3*exp(-0.8*0.5*0.5)});
    \node[below=18pt] at (0,0) {PDF:\; $f(x) \equiv F'(x)$};
  \end{scope}
  % Right panel: CDF
  \begin{scope}[xshift=6cm]
    \draw[->] (-2.5,0) -- (2.7,0) node[right] {$x$};
    \draw[->] (0,-0.1) -- (0,1.6) node[above] {$F(x)$};
    \draw[very thick, domain=-2.4:2.4, samples=80]
      plot (\x, {1.3/(1+exp(-2.2*\x))});
    \draw[dashed] (-1.5,0) node[below]{\small$A$}
                         -- (-1.5, {1.3/(1+exp(2.2*1.5))});
    \draw[dashed] (-0.5,0) node[below]{\small$B$}
                         -- (-0.5, {1.3/(1+exp(2.2*0.5))});
    \node[below=18pt] at (0,0) {CDF:\; $F(x) = \int_{-\infty}^x f(t)\,dt$};
  \end{scope}
\end{tikzpicture}
\caption{Probability density function (PDF) and cumulative density function
(CDF) for a continuous random variable. The shaded area in the PDF equals
$\Pr[A \leq x \leq B]$.}
\label{fig_PDFCDF}
\end{figure}


\section{Moments: Skewness and Kurtosis}
\label{sec_Moments}

Higher-order moments describe the shape of a distribution beyond its mean and
variance.

\noindent\textbf{Skewness} ($3$rd central moment / $\sigma^3$):
\begin{equation}
  \gamma_1 = \frac{\E[(X-\mu)^3]}{\sigma^3}
  \qquad
  \begin{cases} > 0 & \text{right-skewed (positive tail)} \\
                = 0 & \text{symmetric} \\
                < 0 & \text{left-skewed (negative tail)}
  \end{cases}
\end{equation}

\noindent\textbf{Kurtosis} ($4$th central moment / $\sigma^4$):
\begin{equation}
  \gamma_2 = \frac{\E[(X-\mu)^4]}{\sigma^4}
  \qquad
  \begin{cases} > 3 & \text{leptokurtic (fat tails)} \\
                = 3 & \text{mesokurtic (normal)} \\
                < 3 & \text{platykurtic (thin tails)}
  \end{cases}
\end{equation}


\section{Covariance, Correlation, and Joint Distributions}
\label{sec_CovCorr}

\begin{equation}
  \Cov(X,Y) = \sigma_{XY} = \E[(X-\mu_X)(Y-\mu_Y)]
             = \E[XY] - \mu_X\mu_Y
\end{equation}

\begin{equation}
  \rho_{XY} = \frac{\Cov(X,Y)}{\sigma_X \sigma_Y} \;\in\; [-1,+1]
\end{equation}

\noindent\textbf{Covariance and correlation rules:}
\begin{enumerate}[label=(\roman*)]
  \item $\Cov(X,X) = \Var(X)$
  \item $\Cov(aX,bY) = ab\,\Cov(X,Y)$
  \item $\Var(X+Y) = \Var(X) + \Var(Y) + 2\Cov(X,Y)$
  \item If $X \ind Y$: $\Cov(X,Y) = 0$ and $\Var(X+Y) = \Var(X)+\Var(Y)$
\end{enumerate}

\noindent\textbf{Sample estimators:}
\begin{equation}
  \hat{\sigma}_{XY} = \frac{1}{n-1}\sum_{i=1}^{n}(x_i-\bar{x})(y_i-\bar{y}),
  \qquad
  \hat{\rho}_{XY} = \frac{\hat{\sigma}_{XY}}{\hat{\sigma}_X\hat{\sigma}_Y}
\end{equation}


\section{Estimator Properties}
\label{sec_EstimatorProps}

Let $\hat{\theta}$ be an estimator of population parameter $\theta$.

\subsection{Unbiasedness}

$\hat{\theta}$ is \textbf{unbiased} if $\E[\hat{\theta}] = \theta$.
The \textbf{bias} is $\text{Bias}(\hat{\theta}) = \E[\hat{\theta}] - \theta$.

\subsection{Efficiency}

Among all unbiased estimators, $\hat{\theta}$ is \textbf{efficient} (best) if
it has the smallest variance. The Cramér-Rao lower bound defines the minimum
variance achievable by an unbiased estimator.

\begin{figure}[htp]
\centering
\begin{tikzpicture}[>=Stealth, scale=1.0]
  % Efficient estimator (narrow)
  \draw[very thick, domain=-2:2, samples=60]
    plot (\x, {1.8*exp(-3*\x*\x)}) node[right]{\small Efficient $\hat\theta_1$};
  % Inefficient estimator (wide)
  \draw[thick, dashed, domain=-3.5:3.5, samples=80]
    plot (\x, {0.9*exp(-0.6*\x*\x)}) node[right]{\small Inefficient $\hat\theta_2$};
  % True value
  \draw[gray, thin] (0,-0.1) -- (0,1.95);
  \node[below] at (0,-0.1) {$\theta$};
  \draw[->] (-3.8,0) -- (4.2,0) node[right] {$\hat\theta$};
  \draw[->] (0,-0.15) -- (0,2.2);
\end{tikzpicture}
\caption{Efficiency: $\hat\theta_1$ is more efficient than $\hat\theta_2$
because it has smaller variance around the true parameter value $\theta$.}
\label{fig_Efficiency}
\end{figure}

\subsection{Mean Squared Error}

\begin{tcolorbox}[keyresult]
\begin{equation}
  \MSE(\hat{\theta}) = \E[(\hat{\theta} - \theta)^2]
  = \Var(\hat{\theta}) + \bigl[\text{Bias}(\hat{\theta})\bigr]^2
\end{equation}
\end{tcolorbox}

This decomposition shows the fundamental bias--variance trade-off: a slightly
biased estimator may have lower MSE than an unbiased one if its variance is
sufficiently smaller.

\subsection{Consistency}

$\hat{\theta}_n$ is \textbf{consistent} if $\hat{\theta}_n
\xrightarrow{p} \theta$ as $n \to \infty$:
\begin{equation}
  \lim_{n \to \infty} P(|\hat{\theta}_n - \theta| > \varepsilon) = 0
  \quad \forall\,\varepsilon > 0
\end{equation}

\begin{figure}[htp]
\centering
\begin{tikzpicture}[>=Stealth, scale=0.95]
  \def\tp{0}  % true parameter
  \foreach \s/\lbl in {1.2/{$n=10$}, 0.7/{$n=50$}, 0.3/{$n=200$}}{
    \draw[domain={\tp-3.5*\s}:{\tp+3.5*\s}, samples=60, thick]
      plot (\x, {1/(\s*sqrt(2*3.14159))*exp(-0.5*((\x-\tp)/\s)^2)});
  }
  \draw[->] (-4.2,0) -- (4.2,0) node[right]{$\hat\theta$};
  \draw[->] (0,-0.1) -- (0,1.5);
  \draw[gray, thin] (\tp,-0.05) -- (\tp,1.45);
  \node[below] at (\tp,-0.05) {$\theta$};
  \node[right] at (1.5,1.3) {\small $n=200$};
  \node[right] at (1.8,0.55) {\small $n=50$};
  \node[right] at (2.5,0.28) {\small $n=10$};
\end{tikzpicture}
\caption{Consistency: as $n$ increases, the sampling distribution of
$\hat\theta$ shrinks around the true value $\theta$.}
\label{fig_Consistency}
\end{figure}


\section{Statistical Inference}
\label{sec_Inference}

\subsection{The Hypothesis Testing Framework}

\begin{enumerate}
  \item Formulate $H_0$ (null) and $H_1$ (alternative)
  \item Choose significance level $\alpha$ (typically 0.01, 0.05, 0.10)
  \item Compute test statistic and $p$-value
  \item Reject $H_0$ if $p < \alpha$
\end{enumerate}

\begin{table}[htp]
\centering
\caption{Decision outcomes in hypothesis testing}
\label{tbl_TypeErrors}
\renewcommand{\arraystretch}{1.3}
\begin{tabular}{@{}lcc@{}}
\toprule
& $H_0$ True & $H_0$ False \\
\midrule
Fail to reject $H_0$ & Correct ($1-\alpha$) & Type II error ($\beta$) \\
Reject $H_0$         & Type I error ($\alpha$) & Correct ($1-\beta$) \\
\bottomrule
\end{tabular}
\renewcommand{\arraystretch}{1}
\end{table}

\subsection{The $z$-Statistic and $p$-Values}

For a test of $H_0: \mu = \mu_0$:
\begin{equation}
  z = \frac{\bar{x} - \mu_0}{\sigma/\sqrt{n}}
  \sim \mathcal{N}(0,1) \quad \text{under } H_0
\end{equation}

\begin{tcolorbox}[keyresult]
Two-sided $p$-value:\quad $p(\hat{z}) = 2\bigl(1 - \Phi(|\hat{z}|)\bigr)$
\end{tcolorbox}

\subsection{Statistical Power}

The \textbf{power} of a test is $1 - \beta$, the probability of correctly
rejecting a false null hypothesis.

\begin{figure}[htp]
\centering
\begin{tikzpicture}[>=Stealth, scale=1.0]
  % H0 distribution (centred at 0)
  \draw[very thick, domain=-3:3, samples=80]
    plot (\x, {1.2*exp(-0.8*\x*\x)});
  \node at (-1.8, 0.9) {$H_0$};
  % H1 distribution (centred at 2)
  \draw[very thick, dashed, domain=-1:5, samples=80]
    plot (\x, {1.2*exp(-0.8*(\x-2)*(\x-2))});
  \node at (3.8, 0.9) {$H_1$};
  % Critical value
  \draw[gray, thin] (1.4,-0.05) -- (1.4,1.3);
  \node[below] at (1.4,-0.05) {\small $z_\alpha$};
  % Type II error region (shaded)
  \fill[gray!30, domain=-1:1.4, samples=40]
    (-1,0) -- plot (\x, {1.2*exp(-0.8*(\x-2)*(\x-2))}) -- (1.4,0) -- cycle;
  \node at (0.2, 0.18) {\small $\beta$};
  % Rejection region label
  \node at (2.5, 0.2) {\small Power $= 1-\beta$};
  \draw[->] (-3.5,0) -- (5.5,0) node[right]{$z$};
  \draw[->] (0,-0.1) -- (0,1.5);
\end{tikzpicture}
\caption{Statistical power. The shaded region represents the Type II error
probability $\beta$; power is $1-\beta$.}
\label{fig_Power}
\end{figure}

\subsection{Confidence Intervals}

A $(1-\alpha)\times100\%$ confidence interval for $\mu$:
\begin{equation}
  \bar{x} \pm z_{\alpha/2}\,\frac{\sigma}{\sqrt{n}}
\end{equation}

With unknown $\sigma$, replace $z_{\alpha/2}$ with $t_{\alpha/2, n-1}$.
The \textbf{$t$-distribution}\index{t-distribution} with $\nu$ degrees of
freedom has heavier tails than the standard normal and converges to
$\mathcal{N}(0,1)$ as $\nu \to \infty$.

\subsection{Probability Limits and Asymptotic Theory}

The \textbf{probability limit}\index{probability limit} ($\plim$) satisfies
the following rules:

\begin{enumerate}[label=(\roman*)]
  \item $\plim(X + Y) = \plim(X) + \plim(Y)$
  \item $\plim(XY) = \plim(X)\cdot\plim(Y)$
  \item $\plim(X/Y) = \plim(X)/\plim(Y)$ \;(if $\plim(Y) \neq 0$)
  \item $\plim[g(X)] = g(\plim X)$ \;(continuous $g$; Slutsky's theorem)
  \item $\plim(\bar{X}) = \mu$ \;(law of large numbers)
  \item $\plim(\hat{\sigma}^2) = \sigma^2$ \;(consistency of sample variance)
\end{enumerate}


\section{The Central Limit Theorem}
\label{sec_CLT}

\subsection{Lindberg-Lévy CLT}

For i.i.d.\ random variables $x_t$ with mean $\mu$ and variance $\sigma^2$,
the \textbf{Central Limit Theorem}\index{central limit theorem} (CLT) states
that:

\begin{tcolorbox}[keyresult]
\begin{equation}
  z_t \equiv \frac{1}{\sqrt{n}}\sum_{t=1}^{n}\frac{x_t - \mu}{\sigma}
  \xrightarrow{d} \mathcal{N}(0,1) \quad \text{as } n \to \infty
  \label{eqn_CLT}
\end{equation}
\end{tcolorbox}

\subsection{CLT in Practice: How Large Does $n$ Need to Be?}

The CLT becomes effective at surprisingly small $n$ for many distributions:

\begin{itemize}
  \item \textbf{Binomial:} Approximation improves dramatically at $n=50$.
  \item \textbf{Exponential:} CLT takes effect visibly at $n=50$.
  \item \textbf{$t$-distribution:} Even at $n=10$, centered random variables
    approximate $\mathcal{N}(0,1)$ well (the two distributions are similar
    except for heavier tails in the $t$).
  \item \textbf{Gamma/Beta:} Reasonable approximation at smallest $n$,
    reliable at $n=50$.
  \item \textbf{Log-normal:} CLT kicks in only at $n > 100$ due to high skew.
\end{itemize}

\noindent\textbf{Variance and sample size as substitutes.} The $\chi^2(m)$
distribution approaches $\mathcal{N}(m, 2m)$ as $m$ becomes large. A higher
variance implies a larger denominator in the centering process, which in turn
requires \emph{fewer} observations for the distribution to be approximately
standard normal. Variance and sample size are substitutes in the centering
process.

\renewcommand\bibname{Further reading}
\begin{subappendices}
\bibliographystyle{econometrica}
\bibliography{Literature/OverLord}
\IfFileExists{Literature/Review.bbl}{\chapter{Review}\label{chap_Review}
\section{Preliminaries}\label{chap_Review:Prelim}
``Probability approach in econometrics'' by Trygve \citet{Haavelmo1944} stands at the beginning of a paradigm shift in statistical thinking; the \emph{classical} interpretation of statistics is giving way to the \emph{frequentist} interpretation. Under the classical interpretation, probabilities are linked to the principle of indifference and statistical events are governed by natural symmetry (e.g. coin tosses, die etc.) This constraining set-up of what gives rise to probabilities is one of the key problems of the classical view. 

Under the frequentist interpretation on the other hand, probabilities of events are linked to the frequency with which a given event is observed. This implies that

$$P(x)= \lim\limits_{n_{t} \rightarrow \infty} \frac{n_{x}}{n_{t}},$$ 

where $n_{x}$ is the number of observed trials and $n_{t}$ is the total number of trials. This probabilistic approach to statistics, which includes Bayesian statistics\index{Bayesian statistics}, represents the world in terms of ``degrees of belief'' (Bayesian priors).

More formally, Bayes' Theorem is expressed as the conditional probability of occurence A given B such that $$P(A|B) = \frac{P(B|A)P(A)}{P(B)}.$$\index{Bayes' Theorem|see{Bayesian statistics}}


\section{Statistics and probability theory}\label{chap_Review:Stats}
\subsection{Types of random variables}
A random variable, usually written as $X$, is a variable whose possible values are numerical outcomes of a random phenomenon. A random variable $X$  is a (real) number with the following characteristics:

\begin{inparaenum}[\itshape a\upshape)]
\item the value of $X$ is unknown;
\item the possible values $X$ can take are known; the set of such values is called support and denoted by $\bm{\mathcal{S}}_{X}$;
\item if $\bm{\mathcal{A}}$  is a subset of the set of real numbers $\mathbb{R}$  (i.e. $\bm{\mathcal{A}}\subseteq \mathbb{R}$), we know the probability that $X$ belongs to $\bm{\mathcal{A}}$ which is denoted by $P(X\in\bm{\mathcal{A}})$.
\item If (and when) the value $x$ of $X$  becomes known, such that $x=X$, $x$ is then referred to as the realization of $X$. Hence, the support $\bm{\mathcal{S}}_{X}$  is the set of all possible realizations of $X$.
\end{inparaenum}

Random variables generally fall into two categories, discrete and continuous:

\begin{itemize}
\item Discrete random variables\index{random variable!discrete}: Takes on only a countable number of distinct values such as $0,1,2,3,4 \ldots$ Discrete random variables are usually (but not necessarily) counts. If a random variable can take only a finite number of distinct values, then it must be discrete.
\item Continuous random variables\index{random variable!continuous}: Takes an infinite number of possible values. Continuous random variables are usually measurements. Examples include height, weight, the amount of sugar in an orange, the time required to run a mile.
\end{itemize}

In addition to distinguishing between discrete and continuous random variables, we can also categorise random variables into \emph{qualitative} variables or \emph{quantitative} variables. Qualitative\index{qualitative variable|see{random variable}} or categorical variables are either nominal, ordinal or dichotomous. 

\begin{itemize}
\item Nominal variables: \index{random variable!nominal}\index{categorical variable|see{random variable}} A variable for mutually exclusive, but not ordered categories such as, for example, your study might compare five different types of suburbs (inner-ring suburbs, commuter suburb, sitcom suburb, high-income suburb, minority suburb). 
\item An ordinal variable\index{random variable!ordinal} is a random variable where the order matters but not the difference between values. For example, quality of life rankings might express how nice a place is on a scale of 1 to 10. A score of 7 means a nicer place that a score of 5, and that is more than a score of 3. But the difference between the 7 and the 5 may not be the same as that between 5 and 3. The values simply express an order. 
\end{itemize}

Quantitative variables\index{quantitative variable|see{random variable}} are either interval variables where differentiation or transformation preserves the units (such as in the case of temperature measurements in C, F) or ratio variables where a value of 0 means that nothing of that variable is present. Thus temperature measured in Kelvin or height, for example, are ratio variables.

\begin{itemize}
\item An interval variable\index{random variable!interval} is a measurement where the difference between two values is meaningful. The difference between a temperature of 100 degrees and 90 degrees is the same difference as between 90 degrees and 80 degrees. 
\item When working with ratio variables\index{random variable!ratio}, you can look at the ratio of two measurements. A weight of 4 grams is twice a weight of 2 grams, because weight is a ratio variable. A temperature of 100 degrees Celsius is not twice as hot as 50 degrees C, because temperature C is not a ratio variable. Similarly, a pH of 3 is not twice as acidic as a pH of 6, because pH is an interval variable and not a ratio variable. 
\end{itemize}

\begin{table}[htp] 
\centering
\caption{Basic statistical operations by variable type}
\footnotesize
\begin{tabular}{lcccc}\label{tbl_Vars}
\\ \hline
\hline
&\multicolumn{4}{c}{\emph{Random variable type}}\\
\emph{Operation}&Nominal&Ordinal&Ratio&Interval\\
\cline{2-5}
Frequency distribution&Yes&Yes&Yes&Yes\\
Median and percentiles&No&Yes&Yes&Yes\\
Addition/subtraction&No&No&Yes&Yes\\
Mean, standard deviation, standard error of the mean&No&No&Yes&Yes\\
Ratio/coefficient of variation&No&No&No&Yes\\
\hline
\hline \\
\normalsize
\end{tabular}
\end{table}


\subsection{Expectations}
\subsubsection{Expectation of a discrete random variable}
\begin{eqnarray}
\mathbb{E}(X) &=& x_{1}p_{1} + x_{2}p_{2} + x_{3}p_{3} +\ldots + x_{n}p_{n}\nonumber \\
&=& \sum^{n}_{i=1}x_{i}p_{i}
\end{eqnarray}

Recall that in a finite sample, $\mathbb{E}(X)$ denotes the sample mean and the sum of all probabilities must be equal to one, i.e. $\sum_{i}p_{i}=1$.

\paragraph{Example} Consider the following gamble where an unbiased coin with heads ($H$) and tails ($T$) is tossed $n$ times until tails ($T$) appears. If $T$ shows up at the $n$-th toss, the payoff is $\$2^{n-1}$. How much would you be willing to put up as a wager to participate in this game?

\begin{eqnarray*}
\mathbb{E}(X) &=& \$2^{0}\cdot\frac{1}{2} + \$2^{1}\cdot\frac{1}{2}\cdot\frac{1}{2} + \$2^{2}\cdot\frac{1}{2}\cdot\frac{1}{2}\cdot\frac{1}{2} + \ldots \\
&=& \$1\cdot\frac{1}{2} + \$2\cdot\frac{1}{4} + \$4\cdot\frac{1}{8} + \ldots\\
&=& \sum^{\infty}_{i=1}\$ 0.5 = \$\infty
\end{eqnarray*}

This is the \emph{St. Petersburg paradox}\index{St. Petersburg paradox} famously defined and described by Daniel Bernoulli\index{Bernoulli, Daniel} (Swiss mathematician (1700--1782), father of modern fluid dynamics) to illustrate the failure of expected value principle (as opposed to expected utility) in decision theory under uncertainty. Thus the expected payout value of a St. Petersburg game is an infinite amount, yet clearly the willingness to pay (WTP)\index{WTP|see{willingness to pay}}\index{willingness to pay} to participate in such a game must be (much) less than its expected value. The classical resolution of the paradox thus involves the explicit introduction of a utility function\index{utility}, an expected utility hypothesis, and the presumption of diminishing marginal utility of money. This implies $\mathbb{E}(U) < \mathbb{E}(X)$.

\paragraph{Example} What is the expected value of throwing a single die (or a pair of dice)?

\begin{eqnarray*}
\mathbb{E}(X) = \frac{1}{6}\cdot 1 + \frac{1}{6}\cdot 2 + \frac{1}{6}\cdot 3 + \ldots + \frac{1}{6}\cdot 6 = 3.5
\end{eqnarray*}

Thus the sample mean $\mu_{X} = 3.5$ and analogously, for a pair of dice, $\mu_{X} = 7$. A table with the bivariate probabilities is shown in table~\ref{tbl_Dice} below:


\begin{table}[htp] 
\centering
\caption{Bivariate probabilities for a pair of dice and expected value of $X$}
\footnotesize
\bigskip
\begin{tabular}{lll|lll}\label{tbl_Dice}
\\ \hline
\hline
$X$&$p$&$Xp$&$X$&$p$&$Xp$\\
\hline
$x_{1}$ & $p_{1}$ & $x_{1}p_{1}$ & 2 & 1/36 & 2/36\\
$x_{2}$ & $p_{2}$ & $x_{2}p_{2}$ & 3 & 2/36 & 6/36\\
$x_{3}$ & $p_{3}$ & $x_{3}p_{3}$ & 4 & 3/36 & 12/36\\
$\ldots$ & $\ldots$ & $\ldots$ & 5 & 4/36 & 20/36\\
$\ldots$ & $\ldots$ & $\ldots$ & 6 & 5/36 & 30/36\\
$\ldots$ & $\ldots$ & $\ldots$ & 7 & 6/36 & 42/36\\
$\ldots$ & $\ldots$ & $\ldots$ & 8 & 5/36 & 40/36\\
$\ldots$ & $\ldots$ & $\ldots$ & 9 & 4/36 & 36/36\\
$\ldots$ & $\ldots$ & $\ldots$ & 10 & 3/36 & 30/36\\
$\ldots$ & $\ldots$ & $\ldots$ & 11 & 2/36 & 22/36\\
$x_{n}$ & $p_{n}$ & $x_{n}p_{n}$ & 12 & 1/36 & 6/36\\

\hline
\multicolumn{3}{l}{Total: $\mathbb{E}(X)=\sum^{n}_{i=1}x_{i}p_{i}$}&&\multicolumn{2}{r}{252/36 = 7}\\
\hline \\
\normalsize
\end{tabular}
\end{table}


\subsection{Expectations rules}
\begin{enumerate}
\item $\mathbb{E}[X + Y + Z] = \mathbb{E}(X) + \mathbb{E}(Y) + \mathbb{E}(Z);$
\item $\mathbb{E}[bX] = b\mathbb{E}(X);$
\item $\mathbb{E}[b] = b$
\end{enumerate}

\paragraph{Example} Different rules can be combined, e.g. $\mathbb{E}[a + bX] = a + b\mathbb{E}(X)$. Let a normally distributed random variable $V$ can have any mean $\mu$ and any positive variance $\sigma^{2}$. Such a random variable is said to follow the $\mathcal{N}(\mu, \sigma^{2})$ distribution. A standard normal variable therefore has the $\mathcal{N}(0, 1)$ distribution.

Suppose that $V$ has the standard normal distribution. What are the corresponding mean $\mu$ and variance $\sigma^{2}$ of a new random variable $Z \equiv \mu + \sigma V$? For the mean, we use the fact that $\mathbb{E}\left[V\right]=0$ which yields

\begin{eqnarray*}
\mathbb{E}\left[Z\right] &=& \mathbb{E}\left[\mu + \sigma V\right] = \mathbb{E}\left[\mu\right] + \mathbb{E}\left[\sigma V\right] = \mu + \sigma\mathbb{E}\left[V\right]\\
&=& \mu.
\end{eqnarray*}

Similarly, when deriving the variance, recall that $\text{var}(V)=\mathbb{E}\left[V^{2}\right]=1$ such that

\begin{eqnarray*}
    \mathbb{E}\left[(Z-E\left[Z\right])^{2}\right] &=& \mathbb{E}\left[Z^{2}-2Z\mathbb{E}\left[Z\right]+ (\mathbb{E}\left[Z\right])^{2}\right]= \mathbb{E}\left[Z^{2} -2Z\mu + \mu^{2}\right]\\
    &=& \mathbb{E}\left[Z^{2}\right] -2\mu\mathbb{E}\left[Z\right] + \mathbb{E}\left[\mu^{2}\right]= \mathbb{E}\left[Z^{2}\right] -2\mu^{2} + \mu^{2}\\
    &=& \mathbb{E}\left[(\mu^{2}+ 2\mu\sigma V + \sigma^{2}V^{2})\right] -\mu^{2} \\
    &=& \mu^{2} + 2\mu\sigma \mathbb{E}\left[V\right] +\sigma^{2}\mathbb{E}\left[V^{2}\right] -\mu^{2} =  \\
    &=& \sigma^{2}.
    \end{eqnarray*}


\subsubsection{Variance of a discrete random variable}
\begin{eqnarray}
\textrm{Var}(X) = \sigma^{2}_{X} &=& \mathbb{E}[(X-\mu_{X})^{2}]\nonumber \\
&=& (x_{1} - \mu_{X})^{2}p_{1} + (x_{2} - \mu_{X})^{2}p_{2} + \ldots + (x_{n} - \mu_{X})^{2}p_{n} \nonumber\\
&=& \sum^{n}_{i=1}(x_{i} - \mu_{X})^{2}p_{i}
\end{eqnarray}

Derivation of alternative expression for the variance:

\begin{eqnarray}
\textrm{Var}(X) = \sigma^{2}_{X} &=& \mathbb{E}[(X-\mu_{X})^{2}]\nonumber \\
&=& \mathbb{E}[(X-\mu_{X})((X-\mu_{X})] \nonumber\\
&=& \mathbb{E}[X^{2}-2\mu_{X}X + \mu_{X}^{2}] \nonumber\\
&=& \mathbb{E}[X^{2}] -\mathbb{E}[2\mu_{X}X] + \mathbb{E}[\mu_{X}^{2}] \nonumber\\
&=& \mathbb{E}[X^{2}] -2\mu_{X}\mathbb{E}[X] + \mu_{X}^{2} \nonumber\\
&=& \mathbb{E}[X^{2}] -2\mu_{X}^{2} + \mu_{X}^{2} = \mathbb{E}[X^{2}] - \mu_{X}^{2}
\end{eqnarray}

\subsubsection{Fixed and random components of variables}
Let the random variable $X$ be decomposed into two parts, namely a constant mean $\mu_{X}$ and a variable component $u$. Think of $u$ as measurement error or a stochastic disturbance term with $\mathbb{E}(u)= \mathbb{E}(X-\mu_{X}) = \mathbb{E}(X) -\mu_{X}= 0$.

Because all variation in $X$ is due to $u$, the variance of $X$, $\sigma^{2}_{X}$, is identical to the variance of the disturbance term, $\sigma_{u}$. More formally, we can write

\begin{eqnarray*}
\sigma^{2}_{X}&=& \mathbb{E}[(X-\mu_{X})^{2}] = \mathbb{E}(u^{2})\\
\sigma^{2}_{u}&=& \mathbb{E}[(u-\mu_{u})^{2}] = \mathbb{E}(u^{2}).
\end{eqnarray*}

\section{Probability distributions}
All random variables (discrete and continuous) have a probability density function (PDF)\index{PDF|see{density function}}\index{density function!probability} and a cumulative density function (CDF)\index{CDF|see{density function}}\index{density function!cumulative}. These are functions that give the probability that the random variable $X$ is less than or equal to $x$, for every value $x$. For a discrete random variable, the cumulative distribution function is found by summing up the probabilities.


\begin{figure}[htp]
\caption{Probability density function with $f(x)\equiv F'(x)$ }\label{fig_PDF}
\begin{center}
\includegraphics[width=0.7\textwidth]{"Figures/PDFjc"}
\end{center}
\footnotesize
\end{figure}


\begin{figure}[htp]
\caption{Cumulative density function with $\int f(x) \equiv F(x)$}\label{fig_CDF}
\begin{center}
\includegraphics[width=0.7\textwidth]{"Figures/CDFjc"}
\end{center}
\footnotesize
\end{figure}

Using R, it is fairly straightforward to produce the PDF and CDF of a standard normal distribution:

\begin{lstlisting}
#Probability density function with shaded area
cord.x <- c(-2,seq(-2,-1,0.01),-1) 
cord.y <- c(0,dnorm(seq(-2,-1,0.01)),0) 
curve(dnorm(x,0,1),xlim=c(-3,3),main='PDF', xlab='Marginal Density') 
polygon(cord.x,cord.y,col='blue')


#Cumulative density function
x <- rnorm(10000)
plot(ecdf(x), ylab="Probability")
\end{lstlisting}


\begin{eqnarray*}
\text{Pr.}[A\leq x\leq B]=\int_{A}^{B} f(x)dx
\end{eqnarray*}

$\therefore$ we can write the expectation of a continuous random variable as:

\begin{eqnarray*}
\text{CDF}=\int_{-\infty}^{x} f(x)dx = F(B)-F(A)=\int_{-\infty}^{x}F'(x)dx
\end{eqnarray*}

Equivalently, in the case of a continuous random variable the variance can be expressed as:

\begin{eqnarray}
\mathbb{E}\left[x\right]=\int x f(x)dx = \mu_{x}
\end{eqnarray}
\text{similarly}
\begin{eqnarray}
\mathbb{E}\left[(x-\mu_{x})^{2}\right]=\sigma_{x}^{2}=\int(x-\mu_{x})^{2}f(x)dx
\end{eqnarray}



\subsection{Expectation of a continuous random variable}
In the case of a continuous random variable, we can express the expectation in terms of a function of the random variable such that

\begin{eqnarray}
\mathbb{E}(X) &=& \mathbb{E}[g(X)] = g(x_{1})p_{1} + g(x_{2})p_{2} + \ldots + g(x_{n})p_{n} \nonumber\\
&=& \sum^{n}_{i=1}g(x_{i})p_{i}
\end{eqnarray}


\subsection{Lindberg-L\'{e}vy Central Limit Theorem}
For the random variable $x_{t}$ which is \emph{i.i.d.} with mean $\mu$ and variance $\sigma^{2}$, the Lindberg-L\'{e}vy central limit theorem (CLT) states that the quantity

\begin{eqnarray}
z_{t} \equiv \frac{1}{\sqrt{n}}\sum^{n}_{t=1}\frac{x_{t}-\mu}{\sigma}\label{eqn.CLT}
\end{eqnarray}

is asymptotically distributed as N(0,1). The impact of the central limit theorem on different distributions across three sample sizes, $n=10, 50, 100$, is shown below.


\begin{figure}[htp]
\centering
\includegraphics[width=0.49\textwidth]{Figures/binomial_PDF.pdf} \hfill
\includegraphics[width=0.49\textwidth]{Figures/binomial_CDF.pdf}\\
\includegraphics[width=0.49\textwidth]{Figures/expon_PDF.pdf} \hfill
\includegraphics[width=0.49\textwidth]{Figures/expon_CDF.pdf}\\
\caption{Normal approximation for binomial and exponential distribution}\label{fig.CLT1}
\end{figure}

\begin{figure}[htp]
\centering
\includegraphics[width=0.49\textwidth]{Figures/t_PDF.pdf} \hfill
\includegraphics[width=0.49\textwidth]{Figures/t_CDF.pdf}\\
\includegraphics[width=0.49\textwidth]{Figures/f_PDF.pdf} \hfill
\includegraphics[width=0.49\textwidth]{Figures/f_CDF.pdf}\\
\caption{Normal approximation for $t$- and $F$-distribution}\label{fig.CLT2}
\end{figure}

\begin{figure}[htp]
\centering
\includegraphics[width=0.49\textwidth]{Figures/gamma_PDF.pdf} \hfill
\includegraphics[width=0.49\textwidth]{Figures/gamma_CDF.pdf}\\
\includegraphics[width=0.49\textwidth]{Figures/beta_PDF.pdf} \hfill
\includegraphics[width=0.49\textwidth]{Figures/beta_CDF.pdf}\\
\caption{Normal approximation for $\Gamma$- and Beta distribution}\label{fig.CLT3}
\end{figure}

Overall, this exercise demonstrates how powerful the CLT really is; for all six different distributions, $1/\sqrt(n)$ times the sum of $n$ centered random variables approximately follows a normal distribution even at $n$ as small as 10. Figures~\ref{fig.CLT1} to figure~\ref{fig.CLT3} show the empirical distribution functions for the six distributions.

\begin{itemize}
\item[\emph{Binomial}] The approximation to the standard normal distribution increases dramatically at $n=50$ as the step function is still visible at $n=10$.
\item[\emph{Exponential}] CLT takes effect visibly at $n=50$
\item[\emph{t}] Even at $n=10$, the centered random variables from the $t$-distribution approximate N(0,1) well. It is hardly surprising since the two distributions are very similar except that the t-distribution displays wider tails.
\item[\emph{F}] Centering for random variables from the $F$-distribution approximates the normal distribution only at comparatively large $n$, with the CLT fully taking effect at $n=100$
\item[\emph{Gamma}] Centered random variables from the $\Gamma$-distribution approximate N(0,1) reasonably well at the smallest $n$ with the CLT kicking more reliably at $n=50$.
\item[\emph{Beta}] Even at $n=10$, the centered random variables from the Beta-distribution approximate N(0,1) well. However for this distribution, the CLT-threshold seems to depend significantly on the random seed and fluctuated in different simulations between $n=10$ and $n=50$.
\end{itemize}

In addition to discovering that the CLT becomes effective at relatively small $n$ for a variety of distributions, this exercise demonstrates the impact of the variance on the centering process. The $\chi^{2}(m)$ distribution approaches the N($m,2m$) as $m$ becomes large and its variance is $2m$.

\begin{figure}[htp]
\centering
\includegraphics[width=0.75\textwidth]{Figures/chi1_PDF.pdf}
\caption{Normal approximation for $\chi^{2}$-distribution}\label{fig.CLT5}
\end{figure}

The different panels of figure~\ref{fig.CLT6} demonstrate that the \emph{variance} and \emph{sample size} work as \emph{substitutes} in the centering process; a higher variance requires a smaller sample size in order for the $\chi^{2}(m)$-distribution to approximate the standard normal distribution. This largely due to the fact that centering is achieved by subtracting the mean and dividing by the standard deviation. A large variance implies a large denominator which in turn implies that a smaller number of observations are required for the distribution to be asymptotically distributed as N(0,1).

\begin{figure}[htp]
\centering
\includegraphics[width=0.49\textwidth]{Figures/chi2_CDF.pdf} \hfill
\includegraphics[width=0.49\textwidth]{Figures/chi3_CDF.pdf}\\
\includegraphics[width=0.49\textwidth]{Figures/chi4_CDF.pdf} \hfill
\includegraphics[width=0.49\textwidth]{Figures/chi5_CDF.pdf}\\
\caption{Variance impact and CLT for $\chi^{2}$-distribution}\label{fig.CLT6}
\end{figure}


\section{Matrix algebra}\label{chap_Review:Matrix}

\renewcommand\bibname{Further reading}
\begin{subappendices}
\bibliographystyle{econometrica}
\bibliography{Literature/OverLord}
\IfFileExists{Literature/Review.bbl}{\chapter{Review}\label{chap_Review}
\section{Preliminaries}\label{chap_Review:Prelim}
``Probability approach in econometrics'' by Trygve \citet{Haavelmo1944} stands at the beginning of a paradigm shift in statistical thinking; the \emph{classical} interpretation of statistics is giving way to the \emph{frequentist} interpretation. Under the classical interpretation, probabilities are linked to the principle of indifference and statistical events are governed by natural symmetry (e.g. coin tosses, die etc.) This constraining set-up of what gives rise to probabilities is one of the key problems of the classical view. 

Under the frequentist interpretation on the other hand, probabilities of events are linked to the frequency with which a given event is observed. This implies that

$$P(x)= \lim\limits_{n_{t} \rightarrow \infty} \frac{n_{x}}{n_{t}},$$ 

where $n_{x}$ is the number of observed trials and $n_{t}$ is the total number of trials. This probabilistic approach to statistics, which includes Bayesian statistics\index{Bayesian statistics}, represents the world in terms of ``degrees of belief'' (Bayesian priors).

More formally, Bayes' Theorem is expressed as the conditional probability of occurence A given B such that $$P(A|B) = \frac{P(B|A)P(A)}{P(B)}.$$\index{Bayes' Theorem|see{Bayesian statistics}}


\section{Statistics and probability theory}\label{chap_Review:Stats}
\subsection{Types of random variables}
A random variable, usually written as $X$, is a variable whose possible values are numerical outcomes of a random phenomenon. A random variable $X$  is a (real) number with the following characteristics:

\begin{inparaenum}[\itshape a\upshape)]
\item the value of $X$ is unknown;
\item the possible values $X$ can take are known; the set of such values is called support and denoted by $\bm{\mathcal{S}}_{X}$;
\item if $\bm{\mathcal{A}}$  is a subset of the set of real numbers $\mathbb{R}$  (i.e. $\bm{\mathcal{A}}\subseteq \mathbb{R}$), we know the probability that $X$ belongs to $\bm{\mathcal{A}}$ which is denoted by $P(X\in\bm{\mathcal{A}})$.
\item If (and when) the value $x$ of $X$  becomes known, such that $x=X$, $x$ is then referred to as the realization of $X$. Hence, the support $\bm{\mathcal{S}}_{X}$  is the set of all possible realizations of $X$.
\end{inparaenum}

Random variables generally fall into two categories, discrete and continuous:

\begin{itemize}
\item Discrete random variables\index{random variable!discrete}: Takes on only a countable number of distinct values such as $0,1,2,3,4 \ldots$ Discrete random variables are usually (but not necessarily) counts. If a random variable can take only a finite number of distinct values, then it must be discrete.
\item Continuous random variables\index{random variable!continuous}: Takes an infinite number of possible values. Continuous random variables are usually measurements. Examples include height, weight, the amount of sugar in an orange, the time required to run a mile.
\end{itemize}

In addition to distinguishing between discrete and continuous random variables, we can also categorise random variables into \emph{qualitative} variables or \emph{quantitative} variables. Qualitative\index{qualitative variable|see{random variable}} or categorical variables are either nominal, ordinal or dichotomous. 

\begin{itemize}
\item Nominal variables: \index{random variable!nominal}\index{categorical variable|see{random variable}} A variable for mutually exclusive, but not ordered categories such as, for example, your study might compare five different types of suburbs (inner-ring suburbs, commuter suburb, sitcom suburb, high-income suburb, minority suburb). 
\item An ordinal variable\index{random variable!ordinal} is a random variable where the order matters but not the difference between values. For example, quality of life rankings might express how nice a place is on a scale of 1 to 10. A score of 7 means a nicer place that a score of 5, and that is more than a score of 3. But the difference between the 7 and the 5 may not be the same as that between 5 and 3. The values simply express an order. 
\end{itemize}

Quantitative variables\index{quantitative variable|see{random variable}} are either interval variables where differentiation or transformation preserves the units (such as in the case of temperature measurements in C, F) or ratio variables where a value of 0 means that nothing of that variable is present. Thus temperature measured in Kelvin or height, for example, are ratio variables.

\begin{itemize}
\item An interval variable\index{random variable!interval} is a measurement where the difference between two values is meaningful. The difference between a temperature of 100 degrees and 90 degrees is the same difference as between 90 degrees and 80 degrees. 
\item When working with ratio variables\index{random variable!ratio}, you can look at the ratio of two measurements. A weight of 4 grams is twice a weight of 2 grams, because weight is a ratio variable. A temperature of 100 degrees Celsius is not twice as hot as 50 degrees C, because temperature C is not a ratio variable. Similarly, a pH of 3 is not twice as acidic as a pH of 6, because pH is an interval variable and not a ratio variable. 
\end{itemize}

\begin{table}[htp] 
\centering
\caption{Basic statistical operations by variable type}
\footnotesize
\begin{tabular}{lcccc}\label{tbl_Vars}
\\ \hline
\hline
&\multicolumn{4}{c}{\emph{Random variable type}}\\
\emph{Operation}&Nominal&Ordinal&Ratio&Interval\\
\cline{2-5}
Frequency distribution&Yes&Yes&Yes&Yes\\
Median and percentiles&No&Yes&Yes&Yes\\
Addition/subtraction&No&No&Yes&Yes\\
Mean, standard deviation, standard error of the mean&No&No&Yes&Yes\\
Ratio/coefficient of variation&No&No&No&Yes\\
\hline
\hline \\
\normalsize
\end{tabular}
\end{table}


\subsection{Expectations}
\subsubsection{Expectation of a discrete random variable}
\begin{eqnarray}
\mathbb{E}(X) &=& x_{1}p_{1} + x_{2}p_{2} + x_{3}p_{3} +\ldots + x_{n}p_{n}\nonumber \\
&=& \sum^{n}_{i=1}x_{i}p_{i}
\end{eqnarray}

Recall that in a finite sample, $\mathbb{E}(X)$ denotes the sample mean and the sum of all probabilities must be equal to one, i.e. $\sum_{i}p_{i}=1$.

\paragraph{Example} Consider the following gamble where an unbiased coin with heads ($H$) and tails ($T$) is tossed $n$ times until tails ($T$) appears. If $T$ shows up at the $n$-th toss, the payoff is $\$2^{n-1}$. How much would you be willing to put up as a wager to participate in this game?

\begin{eqnarray*}
\mathbb{E}(X) &=& \$2^{0}\cdot\frac{1}{2} + \$2^{1}\cdot\frac{1}{2}\cdot\frac{1}{2} + \$2^{2}\cdot\frac{1}{2}\cdot\frac{1}{2}\cdot\frac{1}{2} + \ldots \\
&=& \$1\cdot\frac{1}{2} + \$2\cdot\frac{1}{4} + \$4\cdot\frac{1}{8} + \ldots\\
&=& \sum^{\infty}_{i=1}\$ 0.5 = \$\infty
\end{eqnarray*}

This is the \emph{St. Petersburg paradox}\index{St. Petersburg paradox} famously defined and described by Daniel Bernoulli\index{Bernoulli, Daniel} (Swiss mathematician (1700--1782), father of modern fluid dynamics) to illustrate the failure of expected value principle (as opposed to expected utility) in decision theory under uncertainty. Thus the expected payout value of a St. Petersburg game is an infinite amount, yet clearly the willingness to pay (WTP)\index{WTP|see{willingness to pay}}\index{willingness to pay} to participate in such a game must be (much) less than its expected value. The classical resolution of the paradox thus involves the explicit introduction of a utility function\index{utility}, an expected utility hypothesis, and the presumption of diminishing marginal utility of money. This implies $\mathbb{E}(U) < \mathbb{E}(X)$.

\paragraph{Example} What is the expected value of throwing a single die (or a pair of dice)?

\begin{eqnarray*}
\mathbb{E}(X) = \frac{1}{6}\cdot 1 + \frac{1}{6}\cdot 2 + \frac{1}{6}\cdot 3 + \ldots + \frac{1}{6}\cdot 6 = 3.5
\end{eqnarray*}

Thus the sample mean $\mu_{X} = 3.5$ and analogously, for a pair of dice, $\mu_{X} = 7$. A table with the bivariate probabilities is shown in table~\ref{tbl_Dice} below:


\begin{table}[htp] 
\centering
\caption{Bivariate probabilities for a pair of dice and expected value of $X$}
\footnotesize
\bigskip
\begin{tabular}{lll|lll}\label{tbl_Dice}
\\ \hline
\hline
$X$&$p$&$Xp$&$X$&$p$&$Xp$\\
\hline
$x_{1}$ & $p_{1}$ & $x_{1}p_{1}$ & 2 & 1/36 & 2/36\\
$x_{2}$ & $p_{2}$ & $x_{2}p_{2}$ & 3 & 2/36 & 6/36\\
$x_{3}$ & $p_{3}$ & $x_{3}p_{3}$ & 4 & 3/36 & 12/36\\
$\ldots$ & $\ldots$ & $\ldots$ & 5 & 4/36 & 20/36\\
$\ldots$ & $\ldots$ & $\ldots$ & 6 & 5/36 & 30/36\\
$\ldots$ & $\ldots$ & $\ldots$ & 7 & 6/36 & 42/36\\
$\ldots$ & $\ldots$ & $\ldots$ & 8 & 5/36 & 40/36\\
$\ldots$ & $\ldots$ & $\ldots$ & 9 & 4/36 & 36/36\\
$\ldots$ & $\ldots$ & $\ldots$ & 10 & 3/36 & 30/36\\
$\ldots$ & $\ldots$ & $\ldots$ & 11 & 2/36 & 22/36\\
$x_{n}$ & $p_{n}$ & $x_{n}p_{n}$ & 12 & 1/36 & 6/36\\

\hline
\multicolumn{3}{l}{Total: $\mathbb{E}(X)=\sum^{n}_{i=1}x_{i}p_{i}$}&&\multicolumn{2}{r}{252/36 = 7}\\
\hline \\
\normalsize
\end{tabular}
\end{table}


\subsection{Expectations rules}
\begin{enumerate}
\item $\mathbb{E}[X + Y + Z] = \mathbb{E}(X) + \mathbb{E}(Y) + \mathbb{E}(Z);$
\item $\mathbb{E}[bX] = b\mathbb{E}(X);$
\item $\mathbb{E}[b] = b$
\end{enumerate}

\paragraph{Example} Different rules can be combined, e.g. $\mathbb{E}[a + bX] = a + b\mathbb{E}(X)$. Let a normally distributed random variable $V$ can have any mean $\mu$ and any positive variance $\sigma^{2}$. Such a random variable is said to follow the $\mathcal{N}(\mu, \sigma^{2})$ distribution. A standard normal variable therefore has the $\mathcal{N}(0, 1)$ distribution.

Suppose that $V$ has the standard normal distribution. What are the corresponding mean $\mu$ and variance $\sigma^{2}$ of a new random variable $Z \equiv \mu + \sigma V$? For the mean, we use the fact that $\mathbb{E}\left[V\right]=0$ which yields

\begin{eqnarray*}
\mathbb{E}\left[Z\right] &=& \mathbb{E}\left[\mu + \sigma V\right] = \mathbb{E}\left[\mu\right] + \mathbb{E}\left[\sigma V\right] = \mu + \sigma\mathbb{E}\left[V\right]\\
&=& \mu.
\end{eqnarray*}

Similarly, when deriving the variance, recall that $\text{var}(V)=\mathbb{E}\left[V^{2}\right]=1$ such that

\begin{eqnarray*}
    \mathbb{E}\left[(Z-E\left[Z\right])^{2}\right] &=& \mathbb{E}\left[Z^{2}-2Z\mathbb{E}\left[Z\right]+ (\mathbb{E}\left[Z\right])^{2}\right]= \mathbb{E}\left[Z^{2} -2Z\mu + \mu^{2}\right]\\
    &=& \mathbb{E}\left[Z^{2}\right] -2\mu\mathbb{E}\left[Z\right] + \mathbb{E}\left[\mu^{2}\right]= \mathbb{E}\left[Z^{2}\right] -2\mu^{2} + \mu^{2}\\
    &=& \mathbb{E}\left[(\mu^{2}+ 2\mu\sigma V + \sigma^{2}V^{2})\right] -\mu^{2} \\
    &=& \mu^{2} + 2\mu\sigma \mathbb{E}\left[V\right] +\sigma^{2}\mathbb{E}\left[V^{2}\right] -\mu^{2} =  \\
    &=& \sigma^{2}.
    \end{eqnarray*}


\subsubsection{Variance of a discrete random variable}
\begin{eqnarray}
\textrm{Var}(X) = \sigma^{2}_{X} &=& \mathbb{E}[(X-\mu_{X})^{2}]\nonumber \\
&=& (x_{1} - \mu_{X})^{2}p_{1} + (x_{2} - \mu_{X})^{2}p_{2} + \ldots + (x_{n} - \mu_{X})^{2}p_{n} \nonumber\\
&=& \sum^{n}_{i=1}(x_{i} - \mu_{X})^{2}p_{i}
\end{eqnarray}

Derivation of alternative expression for the variance:

\begin{eqnarray}
\textrm{Var}(X) = \sigma^{2}_{X} &=& \mathbb{E}[(X-\mu_{X})^{2}]\nonumber \\
&=& \mathbb{E}[(X-\mu_{X})((X-\mu_{X})] \nonumber\\
&=& \mathbb{E}[X^{2}-2\mu_{X}X + \mu_{X}^{2}] \nonumber\\
&=& \mathbb{E}[X^{2}] -\mathbb{E}[2\mu_{X}X] + \mathbb{E}[\mu_{X}^{2}] \nonumber\\
&=& \mathbb{E}[X^{2}] -2\mu_{X}\mathbb{E}[X] + \mu_{X}^{2} \nonumber\\
&=& \mathbb{E}[X^{2}] -2\mu_{X}^{2} + \mu_{X}^{2} = \mathbb{E}[X^{2}] - \mu_{X}^{2}
\end{eqnarray}

\subsubsection{Fixed and random components of variables}
Let the random variable $X$ be decomposed into two parts, namely a constant mean $\mu_{X}$ and a variable component $u$. Think of $u$ as measurement error or a stochastic disturbance term with $\mathbb{E}(u)= \mathbb{E}(X-\mu_{X}) = \mathbb{E}(X) -\mu_{X}= 0$.

Because all variation in $X$ is due to $u$, the variance of $X$, $\sigma^{2}_{X}$, is identical to the variance of the disturbance term, $\sigma_{u}$. More formally, we can write

\begin{eqnarray*}
\sigma^{2}_{X}&=& \mathbb{E}[(X-\mu_{X})^{2}] = \mathbb{E}(u^{2})\\
\sigma^{2}_{u}&=& \mathbb{E}[(u-\mu_{u})^{2}] = \mathbb{E}(u^{2}).
\end{eqnarray*}

\section{Probability distributions}
All random variables (discrete and continuous) have a probability density function (PDF)\index{PDF|see{density function}}\index{density function!probability} and a cumulative density function (CDF)\index{CDF|see{density function}}\index{density function!cumulative}. These are functions that give the probability that the random variable $X$ is less than or equal to $x$, for every value $x$. For a discrete random variable, the cumulative distribution function is found by summing up the probabilities.


\begin{figure}[htp]
\caption{Probability density function with $f(x)\equiv F'(x)$ }\label{fig_PDF}
\begin{center}
\includegraphics[width=0.7\textwidth]{"Figures/PDFjc"}
\end{center}
\footnotesize
\end{figure}


\begin{figure}[htp]
\caption{Cumulative density function with $\int f(x) \equiv F(x)$}\label{fig_CDF}
\begin{center}
\includegraphics[width=0.7\textwidth]{"Figures/CDFjc"}
\end{center}
\footnotesize
\end{figure}

Using R, it is fairly straightforward to produce the PDF and CDF of a standard normal distribution:

\begin{lstlisting}
#Probability density function with shaded area
cord.x <- c(-2,seq(-2,-1,0.01),-1) 
cord.y <- c(0,dnorm(seq(-2,-1,0.01)),0) 
curve(dnorm(x,0,1),xlim=c(-3,3),main='PDF', xlab='Marginal Density') 
polygon(cord.x,cord.y,col='blue')


#Cumulative density function
x <- rnorm(10000)
plot(ecdf(x), ylab="Probability")
\end{lstlisting}


\begin{eqnarray*}
\text{Pr.}[A\leq x\leq B]=\int_{A}^{B} f(x)dx
\end{eqnarray*}

$\therefore$ we can write the expectation of a continuous random variable as:

\begin{eqnarray*}
\text{CDF}=\int_{-\infty}^{x} f(x)dx = F(B)-F(A)=\int_{-\infty}^{x}F'(x)dx
\end{eqnarray*}

Equivalently, in the case of a continuous random variable the variance can be expressed as:

\begin{eqnarray}
\mathbb{E}\left[x\right]=\int x f(x)dx = \mu_{x}
\end{eqnarray}
\text{similarly}
\begin{eqnarray}
\mathbb{E}\left[(x-\mu_{x})^{2}\right]=\sigma_{x}^{2}=\int(x-\mu_{x})^{2}f(x)dx
\end{eqnarray}



\subsection{Expectation of a continuous random variable}
In the case of a continuous random variable, we can express the expectation in terms of a function of the random variable such that

\begin{eqnarray}
\mathbb{E}(X) &=& \mathbb{E}[g(X)] = g(x_{1})p_{1} + g(x_{2})p_{2} + \ldots + g(x_{n})p_{n} \nonumber\\
&=& \sum^{n}_{i=1}g(x_{i})p_{i}
\end{eqnarray}


\subsection{Lindberg-L\'{e}vy Central Limit Theorem}
For the random variable $x_{t}$ which is \emph{i.i.d.} with mean $\mu$ and variance $\sigma^{2}$, the Lindberg-L\'{e}vy central limit theorem (CLT) states that the quantity

\begin{eqnarray}
z_{t} \equiv \frac{1}{\sqrt{n}}\sum^{n}_{t=1}\frac{x_{t}-\mu}{\sigma}\label{eqn.CLT}
\end{eqnarray}

is asymptotically distributed as N(0,1). The impact of the central limit theorem on different distributions across three sample sizes, $n=10, 50, 100$, is shown below.


\begin{figure}[htp]
\centering
\includegraphics[width=0.49\textwidth]{Figures/binomial_PDF.pdf} \hfill
\includegraphics[width=0.49\textwidth]{Figures/binomial_CDF.pdf}\\
\includegraphics[width=0.49\textwidth]{Figures/expon_PDF.pdf} \hfill
\includegraphics[width=0.49\textwidth]{Figures/expon_CDF.pdf}\\
\caption{Normal approximation for binomial and exponential distribution}\label{fig.CLT1}
\end{figure}

\begin{figure}[htp]
\centering
\includegraphics[width=0.49\textwidth]{Figures/t_PDF.pdf} \hfill
\includegraphics[width=0.49\textwidth]{Figures/t_CDF.pdf}\\
\includegraphics[width=0.49\textwidth]{Figures/f_PDF.pdf} \hfill
\includegraphics[width=0.49\textwidth]{Figures/f_CDF.pdf}\\
\caption{Normal approximation for $t$- and $F$-distribution}\label{fig.CLT2}
\end{figure}

\begin{figure}[htp]
\centering
\includegraphics[width=0.49\textwidth]{Figures/gamma_PDF.pdf} \hfill
\includegraphics[width=0.49\textwidth]{Figures/gamma_CDF.pdf}\\
\includegraphics[width=0.49\textwidth]{Figures/beta_PDF.pdf} \hfill
\includegraphics[width=0.49\textwidth]{Figures/beta_CDF.pdf}\\
\caption{Normal approximation for $\Gamma$- and Beta distribution}\label{fig.CLT3}
\end{figure}

Overall, this exercise demonstrates how powerful the CLT really is; for all six different distributions, $1/\sqrt(n)$ times the sum of $n$ centered random variables approximately follows a normal distribution even at $n$ as small as 10. Figures~\ref{fig.CLT1} to figure~\ref{fig.CLT3} show the empirical distribution functions for the six distributions.

\begin{itemize}
\item[\emph{Binomial}] The approximation to the standard normal distribution increases dramatically at $n=50$ as the step function is still visible at $n=10$.
\item[\emph{Exponential}] CLT takes effect visibly at $n=50$
\item[\emph{t}] Even at $n=10$, the centered random variables from the $t$-distribution approximate N(0,1) well. It is hardly surprising since the two distributions are very similar except that the t-distribution displays wider tails.
\item[\emph{F}] Centering for random variables from the $F$-distribution approximates the normal distribution only at comparatively large $n$, with the CLT fully taking effect at $n=100$
\item[\emph{Gamma}] Centered random variables from the $\Gamma$-distribution approximate N(0,1) reasonably well at the smallest $n$ with the CLT kicking more reliably at $n=50$.
\item[\emph{Beta}] Even at $n=10$, the centered random variables from the Beta-distribution approximate N(0,1) well. However for this distribution, the CLT-threshold seems to depend significantly on the random seed and fluctuated in different simulations between $n=10$ and $n=50$.
\end{itemize}

In addition to discovering that the CLT becomes effective at relatively small $n$ for a variety of distributions, this exercise demonstrates the impact of the variance on the centering process. The $\chi^{2}(m)$ distribution approaches the N($m,2m$) as $m$ becomes large and its variance is $2m$.

\begin{figure}[htp]
\centering
\includegraphics[width=0.75\textwidth]{Figures/chi1_PDF.pdf}
\caption{Normal approximation for $\chi^{2}$-distribution}\label{fig.CLT5}
\end{figure}

The different panels of figure~\ref{fig.CLT6} demonstrate that the \emph{variance} and \emph{sample size} work as \emph{substitutes} in the centering process; a higher variance requires a smaller sample size in order for the $\chi^{2}(m)$-distribution to approximate the standard normal distribution. This largely due to the fact that centering is achieved by subtracting the mean and dividing by the standard deviation. A large variance implies a large denominator which in turn implies that a smaller number of observations are required for the distribution to be asymptotically distributed as N(0,1).

\begin{figure}[htp]
\centering
\includegraphics[width=0.49\textwidth]{Figures/chi2_CDF.pdf} \hfill
\includegraphics[width=0.49\textwidth]{Figures/chi3_CDF.pdf}\\
\includegraphics[width=0.49\textwidth]{Figures/chi4_CDF.pdf} \hfill
\includegraphics[width=0.49\textwidth]{Figures/chi5_CDF.pdf}\\
\caption{Variance impact and CLT for $\chi^{2}$-distribution}\label{fig.CLT6}
\end{figure}


\section{Matrix algebra}\label{chap_Review:Matrix}

\renewcommand\bibname{Further reading}
\begin{subappendices}
\bibliographystyle{econometrica}
\bibliography{Literature/OverLord}
\IfFileExists{Literature/Review.bbl}{\input{Literature/Review.bbl}}{\typeout{Need to run bibtex}}
\end{subappendices} }{\typeout{Need to run bibtex}}
\end{subappendices} }{\typeout{Need to run bibtex on Review}}
\end{subappendices}

% =====================================================================
%  Chapter 2: A Matrix Algebra Toolkit
%  Source: P3 transcription (matrix sections)
% =====================================================================

\chapter{A Matrix Algebra Toolkit}
\label{chap_Matrix}

\section{Vectors, Matrices, and Basic Operations}
\label{sec_VecMat}

A column vector $\mathbf{a}_i = (a_{1i},\, a_{2i},\, \ldots,\, a_{ni})'$
stacked horizontally gives the $n \times k$ matrix:
\begin{equation}
  A = \bigl[\mathbf{a}_1\;\mathbf{a}_2\;\cdots\;\mathbf{a}_k\bigr]
  = \begin{bmatrix}
      a_{11} & a_{12} & \cdots & a_{1k} \\
      a_{21} & a_{22} &        & \vdots \\
      \vdots &        & \ddots &        \\
      a_{n1} & \cdots &        & a_{nk}
    \end{bmatrix}_{n \times k}
\end{equation}

The corresponding $j$-th row vector is $\boldsymbol{\alpha}_j' =
(a_{j1},\, a_{j2},\, \ldots,\, a_{jk})$.

\noindent\textbf{Inner product:}
$\mathbf{a}'\mathbf{b} = \mathbf{b}'\mathbf{a} =
\sum_{i=1}^{n} a_i b_i$

\noindent\textbf{The special vector $\boldsymbol{\iota}$:}
$\boldsymbol{\iota}'\mathbf{a} = \sum a_i = n\bar{a}$ (sum of elements);
$\boldsymbol{\iota}'\boldsymbol{\iota} = n$.

\subsection{Key Identities for the Linear Model}

With $X_{n \times k} = [\mathbf{x}_1';\mathbf{x}_2';\cdots;\mathbf{x}_n']$:
\begin{equation}
  X'X = \sum_{i=1}^{n}\mathbf{x}_i\mathbf{x}_i'
  = \begin{bmatrix}
      \sum x_1^2   & \sum x_1 x_2 & \cdots & \sum x_1 x_k \\
      \sum x_2 x_1 & \sum x_2^2   &        & \vdots        \\
      \vdots       &              & \ddots &               \\
      \sum x_k x_1 & \cdots       &        & \sum x_k^2
    \end{bmatrix}_{k \times k}
\end{equation}
\begin{equation}
  X'\mathbf{y} = \sum_{i=1}^{n}\mathbf{x}_i y_i \qquad (k \times 1)
\end{equation}


\section{The Determinant and Invertibility}
\label{sec_Determinant}

For a $k \times k$ matrix $A$:\quad $AA^{-1} = A^{-1}A = I_k$.

Matrix $A$ is \textbf{non-singular}\index{matrix!non-singular} if it is of
full rank, i.e.\ $\operatorname{rank}(A) = k$ (no linearly dependent columns).
Non-singularity is equivalent to invertibility and links directly to the
problem of \textbf{multicollinearity}\index{multicollinearity}.

The \textbf{determinant}\index{matrix!determinant} measures the ``volume'' of
a matrix---geometrically, the signed area of the parallelogram spanned by its
column vectors:
\begin{equation}
  \det(A) \neq 0 \iff A \text{ is non-singular}
\end{equation}

For $A = \begin{bmatrix} a_{11} & a_{12} \\ a_{21} & a_{22} \end{bmatrix}$:
\begin{equation}
  \det(A) = a_{11}a_{22} - a_{12}a_{21}
\end{equation}

The inverse:
\begin{equation}
  A^{-1} = \frac{1}{\det(A)}\operatorname{adj}(A)
  = \frac{1}{\det A}
    \begin{pmatrix} a_{22} & -a_{12} \\ -a_{21} & a_{11} \end{pmatrix},
  \qquad \det(A^{-1}) = \frac{1}{\det(A)}
\end{equation}

\begin{figure}[htp]
\centering
\begin{tikzpicture}[>=Stealth, scale=1.2]
  \coordinate (O)   at (0,   0);
  \coordinate (A1)  at (2.8, 0.4);
  \coordinate (A2)  at (0.6, 2.2);
  \coordinate (TOP) at ($(A1)+(A2)$);
  \fill[gray!20] (O) -- (A1) -- (TOP) -- (A2) -- cycle;
  \draw[dashed, gray!60] (A1) -- (TOP) -- (A2);
  \draw[->, very thick] (O) -- (A1) node[right] {$\mathbf{a}_1$};
  \draw[->, very thick] (O) -- (A2) node[above left] {$\mathbf{a}_2$};
  \node[below left] at (O) {$O$};
  \node at (1.6, 1.1) {\small $\det(A) = a_{11}a_{22} - a_{12}a_{21}$};
\end{tikzpicture}
\caption{The determinant as the signed area of the parallelogram spanned
by the column vectors $\mathbf{a}_1$ and $\mathbf{a}_2$. Linear dependence
collapses the parallelogram to zero area, making $\det(A) = 0$.}
\label{fig_Determinant}
\end{figure}


\section{Orthogonal Projections}
\label{sec_Projections}

Two \textbf{projection matrices}\index{projection matrix} are central to the
geometry of the linear model (developed fully in Chapter~\ref{chap_LinearModel}):

\begin{tcolorbox}[keyresult]
\begin{alignat*}{2}
  P_X &= X(X'X)^{-1}X'
    \qquad &&\text{(hat matrix; projects onto column space of }X\text{)} \\
  M_X &= I_n - P_X
    \qquad &&\text{(annihilator; orthogonal complement)}
\end{alignat*}
$P_X + M_X = I_n$ \quad (complementary projection matrices)
\end{tcolorbox}

Key properties (proofs in Chapter~\ref{chap_LinearModel}):
\begin{itemize}
  \item Idempotency: $P_X P_X = P_X$;\quad $M_X M_X = M_X$
  \item Orthogonality: $P_X M_X = 0$
  \item Annihilation: $M_X X = 0$
\end{itemize}

Both matrices are \textbf{symmetric} ($P_X' = P_X$, $M_X' = M_X$) and
\textbf{positive semi-definite}.

\paragraph{Why $(X'X)^{-1}$ acts as a corrective lens.} Suppose
hypothetically that $(X'X)^{-1} = I$. This would require: (i) off-diagonal
elements equal zero (variables are mutually uncorrelated); (ii) diagonal
elements equal one ($\mu_x = 0$, $\sigma_x = 1/\sqrt{N}$). Then $\hat\beta =
IX'\mathbf{y} = X'\mathbf{y}$ (the naïve regression coefficient). The matrix
$(X'X)^{-1}$ therefore compresses, expands, and rotates the naïve coefficients
to account for the scale and mutual correlations of the regressors.
By extension, $s^2(X'X)^{-1}$ provides the full variance--covariance matrix
of $\hat\beta$.


\section{Centering: The $M_\iota$ Matrix}
\label{sec_Centering}

Let $\bar{x} = \tfrac{1}{n}\sum x_i$. The centred variable
$\mathbf{z} \equiv \mathbf{x} - \bar{x}\boldsymbol{\iota}$ is orthogonal
to the constant vector $\boldsymbol{\iota}$:
\begin{equation}
  \boldsymbol{\iota}'\mathbf{z}
  = \boldsymbol{\iota}'(\mathbf{x} - \bar{x}\boldsymbol{\iota})
  = n\bar{x} - \bar{x}n = 0
\end{equation}

Centering is performed by the orthogonal projection matrix:

\begin{tcolorbox}[keyresult]
\begin{equation}
  M_\iota = I_n - P_\iota
          = I - \boldsymbol{\iota}(\boldsymbol{\iota}'\boldsymbol{\iota})^{-1}
            \boldsymbol{\iota}'
          = I - \frac{\boldsymbol{\iota}\boldsymbol{\iota}'}{n}
\end{equation}
\end{tcolorbox}

\noindent\textit{Proof:} $M_\iota\mathbf{x} = (I-P_\iota)\mathbf{x} =
\mathbf{x} - \boldsymbol{\iota}(\boldsymbol{\iota}'\boldsymbol{\iota})^{-1}
\boldsymbol{\iota}'\mathbf{x} = \mathbf{x} - \bar{x}\boldsymbol{\iota} =
\mathbf{z}$.\hfill$\square$

The TSS decomposition uses $M_\iota$ directly:
\begin{equation}
  \text{TSS} = \mathbf{y}'M_\iota\mathbf{y}
  = \sum_{i=1}^{n}(y_i - \bar{y})^2
\end{equation}


\section{Eigenvalues and the Spectral Radius}
\label{sec_Eigenvalues}

For a square matrix $A$, a scalar $\lambda$ is an \textbf{eigenvalue}
\index{eigenvalue} if there exists a non-zero vector $\mathbf{v}$ such that
$A\mathbf{v} = \lambda\mathbf{v}$. The set of eigenvalues is the
\textbf{spectrum} of $A$; the largest in absolute value is the
\textbf{spectral radius}, $\rho(A) = \max|\lambda_i|$.

\noindent\textbf{Why eigenvalues matter for spatial analysis.}
The SAR model $\mathbf{y} = \rho W\mathbf{y} + \mathbf{e}$ requires
$(I - \rho W)$ to be invertible. This holds if and only if $\rho \neq
1/\lambda_i(W)$ for all eigenvalues $\lambda_i$ of $W$. For a
row-standardised weight matrix, $\lambda_{\max}(W) = 1$, so the
stationarity condition reduces to $|\rho| < 1$ (Chapter~\ref{chap_Spatial}).

\renewcommand\bibname{Further reading}
\begin{subappendices}
\bibliographystyle{econometrica}
\bibliography{Literature/OverLord}
\IfFileExists{Literature/Matrix.bbl}{\input{Literature/Matrix.bbl}}{\typeout{Need bibtex on Matrix}}
\end{subappendices}

% =====================================================================
%  Chapter 3: Data, Measurement, and Open Infrastructure
%  Source: P3 transcription (ACS, GDP/GDI, price indices) + new content
% =====================================================================

\chapter{Data, Measurement, and Open Infrastructure}
\label{chap_Data}

\section{The Data Landscape for Urban and Regional Research}
\label{sec_DataLandscape}

Modern urban and regional research rests on an unusually rich base of public
data infrastructure. The US federal statistical system produces consistent,
publicly accessible data on population, housing, labour markets, income,
prices, and the built environment at fine geographic resolution---often down
to the census block or census tract. This infrastructure is not merely a
convenience; it is the empirical foundation that makes regional policy
informatics possible without proprietary datasets or expensive licensing fees.

This chapter provides a practical guide to the key data sources, their
structure and limitations, and the open-source R and Python tools for
accessing them. It also covers measurement concepts---price indices, national
accounts, margins of error---that recur throughout the book.

\noindent\textbf{A note on geographic aggregation.} Before working with any
spatial data, it is essential to understand the areal unit at which
observations are recorded and the implications of aggregation choices. The
Modifiable Areal Unit Problem (MAUP)\index{MAUP|see{Modifiable Areal Unit
Problem}}\index{Modifiable Areal Unit Problem} arises because statistical
results can change substantially when the same data are aggregated to
different zone systems---from census tracts to ZIP codes to counties to
metropolitan areas. The MAUP is treated at length in
Chapter~\ref{chap_Spatial}; here we note only that the choice of geographic
unit is never purely technical, and results should be routinely checked for
sensitivity to that choice.


\section{Survey Data and Statistical Precision}
\label{sec_SurveyData}

\subsection{American Community Survey (ACS) and Margins of Error}

ACS tables report \textbf{margins of error} (MOE) that permit construction
of 90\% confidence bands around estimates. The relationship between MOE
and standard error is:

\begin{tcolorbox}[keyresult]
\begin{equation}
  \text{s.e.} = \frac{\text{MOE}_\text{ACS}}{1.645}
  \qquad\Longleftrightarrow\qquad
  \text{MOE} = \text{s.e.} \times z\text{-score}
\end{equation}
\end{tcolorbox}

\noindent\textbf{Testing differences in population means.} When comparing
an estimate $M_1$ from geography 1 with estimate $M_2$ from geography 2,
using standard errors $s_1$ and $s_2$ derived from the ACS MOEs:
\begin{equation}
  \left|\frac{M_1 - M_2}{\sqrt{s_1^2 + s_2^2}}\right| > z_\text{confidence}
\end{equation}

If this test statistic exceeds the critical value (e.g.\ 1.645 for 90\%
confidence), the difference is statistically significant at that level.

The ACS is available at multiple geographic levels and in 1-year and 5-year
estimates. The 5-year estimates are more precise (lower MOE) but cover a
longer time window; the 1-year estimates are more timely but available only
for geographies with population $\geq 65{,}000$.


\section{National and Regional Economic Accounts}
\label{sec_NationalAccounts}

\subsection{GDP and Gross Domestic Income (GDI)}

Gross Domestic Income (GDI) measures output from the income side
(NIPA)\index{NIPA|see{National Income and Product Accounts}}. Its
components sum to GDP:

\begin{table}[htp]
\centering
\caption{GDI components as shares of GDP (approximate, 2003)}
\label{tbl_GDI}
\renewcommand{\arraystretch}{1.3}
\begin{tabular}{@{}lc@{}}
\toprule
Component & Approximate share \\
\midrule
Compensation of employees & 56\% \\
Taxes on production and imports (TOPI), less subsidies & 7\% \\
Net operating surplus (NOS) & 25\% \\
Consumption of fixed capital (CFC) & 12\% \\
Statistical discrepancy & $\approx 0.2\%$ \\
\midrule
\textbf{Total $\equiv$ GDP} & \textbf{100\%} \\
\bottomrule
\end{tabular}
\renewcommand{\arraystretch}{1}
\end{table}

\noindent where TOPI $\equiv$ sales taxes, property taxes, custom duties;
NOS $\equiv$ corporate profits, net earnings of unincorporated businesses,
government surplus; CFC $\equiv$ depreciation, capital expenditures, interest.

\subsection{Measuring Price Changes}

\begin{itemize}
  \item \textbf{CPI:}\index{CPI|see{Consumer Price Index}}
    Fixed basket of goods tracked over time; historically a
    Laspeyres index.
  \item \textbf{GDP deflator:}
    $\displaystyle\frac{\text{Nominal GDP}}{\text{Real GDP}}\times100$.
    Differs from CPI in that basket weights are \emph{not} fixed.
\end{itemize}

Two classical fixed-weight price indices:

\begin{table}[htp]
\centering
\caption{Fixed-weight price indices and their biases}
\label{tbl_PriceIndices}
\renewcommand{\arraystretch}{1.6}
\begin{tabular}{@{}llc@{}}
\toprule
Index & Formula & Bias \\
\midrule
\textbf{Paasche}\index{Paasche index}
  & $\dfrac{\displaystyle\sum p_{it}q_{it}}{\displaystyle\sum p_{i0}q_{it}}$
  & Understates inflation \\[4pt]
\textbf{Laspeyres}\index{Laspeyres index}
  & $\dfrac{\displaystyle\sum p_{it}q_{i0}}{\displaystyle\sum p_{i0}q_{i0}}$
  & Overstates inflation \\
\bottomrule
\end{tabular}
\renewcommand{\arraystretch}{1}
\end{table}

The BLS now uses the \textbf{Fisher ideal index}\index{Fisher ideal index},
the geometric mean of the Paasche ($P$) and Laspeyres ($L$) indices:
\begin{equation}
  \text{Fisher} = \sqrt{P \times L}
  \qquad\text{cf.}\quad
  \underbrace{\Bigl(\prod_{i=1}^{n}a_i\Bigr)^{1/n}}_{\text{geometric mean}}
  \;\text{vs.}\;
  \underbrace{\frac{1}{n}\sum_{i=1}^{n}a_i}_{\text{arithmetic mean}}
\end{equation}


\section{The Open Data Stack}
\label{sec_OpenData}

Table~\ref{tbl_OpenData} summarises the key public data sources for regional
analytics and the primary R packages for accessing each. Together these
sources enable a complete regional analytics workflow without any proprietary
data.

\begin{table}[htp]
\centering
\caption{Public data infrastructure for regional analytics}
\label{tbl_OpenData}
\renewcommand{\arraystretch}{1.4}
\begin{tabular}{@{}p{3.8cm}p{5.5cm}p{3cm}@{}}
\toprule
\textbf{Source} & \textbf{Content} & \textbf{R package} \\
\midrule
Census / ACS & Demographics, housing, income by geography & \texttt{tidycensus} \\
TIGER/Line & All US administrative geographies & \texttt{tigris} \\
FRED & Macro and regional time series & \texttt{fredr} \\
LEHD / LODES & Workplace--residence flows at block level & \texttt{lehdr} \\
BLS QCEW & Employment by county $\times$ industry $\times$ quarter & Direct API \\
BEA Regional Accounts & GDP by metro; personal income by county & \texttt{bea.R} \\
HUD portal & CHAS, HMDA, LIHTC, Fair Market Rents & \texttt{hudapi} \\
IPUMS & Harmonised microdata 1850--present & \texttt{ipumsr} \\
OpenStreetMap & Networks, amenities, land use & \texttt{osmdata} \\
Zillow Research & ZHVI, ZORI at ZIP code level & Direct download \\
EPA EJSCREEN & Environmental justice by census geography & Direct API \\
USDA ERS & Rural--urban classification; agricultural land & Direct download \\
DOT HPMS & Highway infrastructure & Direct download \\
FTA NTD & Transit system data & Direct download \\
\bottomrule
\end{tabular}
\renewcommand{\arraystretch}{1}
\end{table}

\noindent\textbf{The R stack that replaces ArcGIS.} The combination
\texttt{sf} + \texttt{terra} + \texttt{tidycensus} + \texttt{osmdata} +
\texttt{spdep} + \texttt{spatialreg} covers essentially all ArcGIS
social-science applications, with the crucial difference that every operation
is scriptable, reproducible, and version-controlled. This is a scientific
claim about the validity of methods---not merely a software preference. A
point-and-click GIS workflow produces irreproducible analysis; an R workflow
produces reproducible analysis.


\section{Reproducible Workflows: The Technical Infrastructure}
\label{sec_Reproducible}

The open-source commitment described in the Preface requires a minimum
technical infrastructure. Three tools form the core:

\begin{enumerate}
  \item \textbf{Quarto} (\texttt{quarto.org}): A literate programming
    platform that combines prose, code (R or Python), and mathematical
    notation in a single source file, rendering to PDF, HTML, or
    presentations. Every example in this book was written and rendered
    using Quarto.

  \item \textbf{GitHub}: Version control and collaboration. Maintaining a
    public repository of code and data supporting published analyses is
    increasingly treated as a condition of publication in economics journals
    and should be the norm in planning and regional science.

  \item \textbf{Package management}: \texttt{renv} (R) and
    \texttt{requirements.txt} / \texttt{pyproject.toml} (Python) pin
    package versions to ensure that an analysis runs identically on every
    machine, now and in future. Run \texttt{renv::snapshot()} after
    installing any new R package.
\end{enumerate}

\noindent\textbf{Agentic coding assistance.} Large language model assistants
(GitHub Copilot, Claude, ChatGPT) can generate competent spatial regression
code from a plain-English description. This capability makes the theoretical
content of this book \emph{more} important, not less. When the code is
generated automatically, the limiting factor on analysis quality becomes
whether the analyst can evaluate the output---does the weight matrix choice
make sense? Is the identification strategy credible? Does the coefficient
have a causal interpretation? These are questions that require the material
in Parts~II and~III; they cannot be answered by a coding assistant.

% TODO: Add worked example of tidycensus data pull and ACS MOE computation

\renewcommand\bibname{Further reading}
\begin{subappendices}
\bibliographystyle{econometrica}
\bibliography{Literature/OverLord}
\IfFileExists{Literature/Data.bbl}{\chapter{Data sources}\label{chap_Data}}{\typeout{Need bibtex on Data}}
\end{subappendices}


% ── Part II: The Linear Model and Causal Inference ────────────────────
\part{The Linear Model and Causal Inference}\label{part2}
% =====================================================================
%  Chapter 4: The Linear Model
%  Source: ClassNormReg.tex (Overleaf) + P3 transcription
% =====================================================================

\chapter[The Linear Model]{The Classical Linear Model}
\label{chap_LinearModel}

\section{The Regression Function}
\label{sec_RegFunction}

The general regression model is:
\begin{equation}
  y = m(\mathbf{x}) + e, \qquad
  m(\mathbf{x}) \equiv \E[y \mid \mathbf{x}]
  = \int_{-\infty}^{\infty} y\, f(y \mid \mathbf{x})\, dy
\end{equation}

By construction, the regression error $e = y - m(\mathbf{x})$ satisfies:
\begin{enumerate}[label=(\roman*)]
  \item $\E[e \mid \mathbf{x}] = 0$
  \item $\E[e] = 0$
  \item $\E[e\mathbf{x}] = 0$
\end{enumerate}

\subsection{Homoskedasticity vs.\ Heteroskedasticity}

The term \emph{skedasis} means dispersion. Two cases arise for the conditional
error variance $\E[e^2 \mid \mathbf{x}]$:

\begin{enumerate}[label=\alph*)]
  \item \textbf{Homoskedasticity:}\; $\E[e^2\mid\mathbf{x}] = \sigma^2$
    (constant conditional variance)
  \item \textbf{Heteroskedasticity:}\;
    $\E[e^2\mid\mathbf{x}] = \sigma^2(\mathbf{x})$
    (skedastic function; variance depends on $\mathbf{x}$)
\end{enumerate}

The unconditional variance via the law of iterated expectations:
\begin{equation}
  \sigma^2 = \E[e^2] = \E\!\bigl[\E[e^2 \mid \mathbf{x}]\bigr]
  = \E[\sigma^2(\mathbf{x})]
\end{equation}

The \textbf{linear model}\index{linear model} specifies:
\begin{alignat*}{2}
  y &= \mathbf{x}'\beta + e \\
  \E[e \mid \mathbf{x}] &= 0 \\
  \E[e^2 \mid \mathbf{x}] &= \sigma^2 \quad\text{or}\quad \sigma^2(\mathbf{x})
\end{alignat*}

\noindent\textbf{Modelling conditional variance.}\index{ARCH}\index{GARCH}
Robert Engle's ARCH model (1982), for which he received the 2003 Nobel Prize,
specifies $\E[e^2 \mid \mathbf{x}] = \sigma^2(\mathbf{x})$ as an explicit
function of past squared errors. The generalised GARCH model extends this to
allow the conditional variance to depend on past conditional variances as well.
These models are central to financial time series but also arise in spatial
panels where variance may cluster geographically.


\section{The Best Linear Predictor}
\label{sec_BLP}

Minimise the expected squared prediction error over $\beta$:
\begin{equation}
  F(\beta) = \E\!\bigl[(y - \mathbf{x}'\beta)^2\bigr]
  = \E[y^2] - 2\beta'\E[\mathbf{x}y] + \beta'\E[\mathbf{x}\mathbf{x}']\beta
\end{equation}

\noindent\textbf{First order condition:}
$\partial F / \partial\beta = -2\E[\mathbf{x}y] + 2\E[\mathbf{x}\mathbf{x}']\beta = 0$.

If $\E[\mathbf{x}\mathbf{x}']$ is invertible:
\begin{tcolorbox}[keyresult]
\begin{equation}
  \beta = \bigl(\E[\mathbf{x}\mathbf{x}']\bigr)^{-1}\E[\mathbf{x}y]
\end{equation}
\end{tcolorbox}

\subsection{The Empirical Analogue: OLS}

Replace population expectations with sample averages:
\begin{equation}
  F_n(\beta) = \frac{1}{n}\sum(y - \mathbf{x}'\beta)^2
  = \frac{1}{n}\,\text{SSE}_n(\beta)
\end{equation}

Setting $\partial\,\text{SSE}/\partial\beta = 0$:
\begin{equation}
  \hat\beta = \Bigl(\sum_i\mathbf{x}_i\mathbf{x}_i'\Bigr)^{-1}
              \Bigl(\sum_i\mathbf{x}_i y_i\Bigr)
\end{equation}


\section{The Model in Matrix Notation}
\label{sec_MatrixNotation}

The system of $n$ equations $y_i = \mathbf{x}_i'\beta + e_i$ is written
compactly as:
\begin{equation}
  \underset{n\times1}{\mathbf{y}}
  = \underset{n\times k}{X}\;\underset{k\times1}{\beta}
  + \underset{n\times1}{\mathbf{e}}
\end{equation}

\begin{tcolorbox}[keyresult]
\textbf{OLS estimator (matrix form):}
\begin{equation}
  \hat\beta = (X'X)^{-1}X'\mathbf{y}
  \label{eq_OLS}
\end{equation}
Residuals: $\hat{\mathbf{e}} = \mathbf{y} - X\hat\beta$.\quad
Variance estimate: $\hat\sigma^2 = n^{-1}\hat{\mathbf{e}}'\hat{\mathbf{e}}$.
\end{tcolorbox}

\noindent\textbf{Three ways to write the linear model:}
\begin{enumerate}
  \item Equation by equation:\; $y_i = \mathbf{x}_i'\beta + e_i$
  \item Matrix form:\; $\mathbf{y} = X\beta + \mathbf{e}$
  \item Scalar form:\; $y_i = \beta_0 + \beta_1 x_{1i} + \cdots
    + \beta_{k-1}x_{k-1,i} + e_i$
\end{enumerate}


\section{Sampling Distribution of $\hat\beta$}
\label{sec_SamplingDist}

$\hat\beta$ is a random variable with its own $\mu_{\hat\beta}$ and
$\sigma^2_{\hat\beta}$, and therefore has a \textbf{sampling
distribution}\index{sampling distribution}.

\subsection{Unbiasedness of OLS}

\begin{equation}
  \E(\hat\beta \mid X) = \beta \qquad\text{and}\qquad \E(\hat\beta) = \beta
\end{equation}

\noindent\textit{Proof:}
\begin{align*}
  \E(\hat\beta \mid X)
  &= \E\!\bigl[(X'X)^{-1}X'\mathbf{y} \mid X\bigr] \\
  &= (X'X)^{-1}X'\,\E[\mathbf{y} \mid X]
   = (X'X)^{-1}X'X\beta = \beta \qquad\square
\end{align*}

\subsection{Variance of $\hat\beta$}

Under homoskedasticity, $\E[\mathbf{e}\mathbf{e}'] = \sigma^2 I_n$:

\begin{tcolorbox}[keyresult]
\begin{equation}
  V(\hat\beta) = \sigma^2(X'X)^{-1}
  \label{eq_VarBeta}
\end{equation}
\end{tcolorbox}

\noindent\textit{Proof:} $\hat\beta - \beta = (X'X)^{-1}X'\mathbf{e}$, so:
\begin{align*}
  V(\hat\beta)
  &= \E\!\bigl[(\hat\beta-\beta)(\hat\beta-\beta)'\bigr]
   = (X'X)^{-1}X'\,\E[\mathbf{e}\mathbf{e}']\,X(X'X)^{-1} \\
  &= (X'X)^{-1}X'\,\sigma^2 I\,X(X'X)^{-1}
   = \sigma^2(X'X)^{-1} \qquad\square
\end{align*}

\subsection{Estimating $\sigma^2$}

OLS does not directly provide an estimate of $\sigma^2$. Using residuals
$\hat{e}_i$ in place of errors $e_i$ yields a downward-biased estimate:
\begin{equation}
  \E[\hat\sigma^2] = \frac{n-k}{n}\,\sigma^2 < \sigma^2
\end{equation}

\begin{tcolorbox}[keyresult]
\begin{equation}
  s^2 \equiv \frac{1}{n-k}\sum_{i=1}^{n}\hat{e}_i^2
  = \frac{n}{n-k}\,\hat\sigma^2
  \qquad \text{(unbiased estimator of }\sigma^2\text{)}
\end{equation}
\end{tcolorbox}

Note: while $s^2$ is unbiased for $\sigma^2$, $s$ is \emph{not} unbiased
for $\sigma$.


\section{Geometry of OLS}
\label{sec_GeometryOLS}

Expanding $\hat\beta = (X'X)^{-1}X'\mathbf{y}$:
\begin{align*}
  \hat\beta
  &= (X'X)^{-1}X'(X\beta + \mathbf{e}) \\
  &= \beta + (X'X)^{-1}X'\mathbf{e}
\end{align*}
Fitted values: $\hat{\mathbf{y}} = X\hat\beta = X(X'X)^{-1}X'\mathbf{y} =
P_X\mathbf{y}$.

\subsection{Projection Matrices and Their Relationships}
\label{subsec_Projections}

The \textbf{hat matrix}\index{hat matrix} and its complement
(Section~\ref{sec_Projections}) satisfy additional structural relationships
when $X = [X_1\;X_2]$. Defining $P_1 = X_1(X_1'X_1)^{-1}X_1'$:
\begin{align}
  P_X P_1 &= P_1 \qquad\text{since}\quad P_1 < P_X \label{eq_PP1} \\
  M_X M_1 &= M_X \qquad\text{since}\quad M_1 > M_X \label{eq_MM1}
\end{align}

The orthogonal decomposition $\mathbf{y} = P_X\mathbf{y} + M_X\mathbf{y}
= \hat{\mathbf{y}} + \hat{\mathbf{e}}$ is illustrated in
Figure~\ref{fig_GeomOLS}.

\begin{figure}[htp]
\centering
\begin{tikzpicture}[>=Stealth, scale=1.3]
  \fill[gray!12]
    (-0.5,-0.35) -- (3.7,-0.35) -- (4.2,0.95) -- (0.0,0.95) -- cycle;
  \node[gray!70, font=\small] at (1.8, -0.05) {column space of $X$};
  \coordinate (O)   at (1.4, 0.15);
  \fill (O) circle (1.5pt) node[below left] {$O$};
  \coordinate (Y)   at (3.4, 3.1);
  \coordinate (PXY) at (3.4, 0.60);
  \draw[->, very thick]     (O) -- (Y)   node[right] {$\mathbf{y}$};
  \draw[->, thick]          (O) -- (PXY) node[below right]
    {$\hat{\mathbf{y}} = P_X\mathbf{y}$};
  \draw[->, thick, dashed] (PXY) -- (Y)  node[midway, right=2pt]
    {$\hat{\mathbf{e}} = M_X\mathbf{y}$};
  \draw ($(PXY)+(-0.17,0)$) -- ++(0,0.17) -- ++(0.17,0);
  \coordinate (X1) at (2.8, 0.42);
  \coordinate (X2) at (1.9, 0.80);
  \draw[->, thick] (O) -- (X1) node[below] {$X_1$};
  \draw[->, thick] (O) -- (X2) node[above left] {$X_2$};
\end{tikzpicture}
\caption{Geometry of OLS. $\hat{\mathbf{y}} = P_X\mathbf{y}$ is the
orthogonal projection of $\mathbf{y}$ onto the column space of $X$; the
residual vector $\hat{\mathbf{e}} = M_X\mathbf{y}$ is perpendicular to
that space. Note $\hat{\mathbf{y}}'\hat{\mathbf{e}} = 0$.}
\label{fig_GeomOLS}
\end{figure}


\section{Goodness of Fit}
\label{sec_GoodnessOfFit}

\subsection{The TSS Decomposition}

\begin{tcolorbox}[keyresult]
\begin{alignat*}{2}
  \text{TSS} &= \mathbf{y}'M_\iota\mathbf{y}
    = \textstyle\sum(y_i-\bar{y})^2 \\
  \text{ESS} &= \mathbf{y}'M_\iota P_X\mathbf{y}
    = \textstyle\sum(\hat{y}_i-\bar{y})^2 \\
  \text{RSS} &= \mathbf{y}'M_X\mathbf{y}
    = \textstyle\sum(y_i-\hat{y}_i)^2
\end{alignat*}
\vspace{-4pt}
\begin{equation}
  R^2 = \frac{\text{ESS}}{\text{TSS}} = 1 - \frac{\text{RSS}}{\text{TSS}}
\end{equation}
\end{tcolorbox}

\subsection{Adjusted $R^2$}

\begin{equation}
  \bar{R}^2 = 1 - (1-R^2)\frac{n-1}{n-k}
\end{equation}

$\bar{R}^2$ increases only if the $t$-statistic of the newly added variable
exceeds $|1|$ in absolute value. An increase in $R^2$ is \emph{not}
indicative of improved specification---even poorly specified equations can
achieve high $R^2$.


\section{The Frisch--Waugh--Lovell Theorem}
\label{sec_FWL}

\begin{tcolorbox}[definition]
\textbf{Theorem (FWL; see \citealt{Lovell2008,Davidson2004}).} If
$Y = X_1\beta_1 + X_2\beta_2 + \mathbf{e}$ ($X_1$ is $n\times k_1$;
$X_2$ is $n\times k_2$), then the OLS estimate of $\beta_2$ from the full
regression equals the OLS estimate from the auxiliary regression:
\begin{equation}
  M_{X_1}Y = M_{X_1}X_2\beta_2 + \mathbf{v}, \quad \mathbf{v} = M_{X_1}\mathbf{e}
\end{equation}
where $M_{X_1} = I - X_1(X_1'X_1)^{-1}X_1'$.
\end{tcolorbox}

\noindent\textbf{Proof.} Write $M_{X_1} \equiv M_1$.
Premultiply $y = X_1\hat\beta_1 + X_2\hat\beta_2 + M_Xy$ by $X_2'M_1$:
\begin{equation}
  X_2'M_1y
  = \underbrace{X_2'M_1X_1\hat\beta_1}_{(1)}
  + \underbrace{X_2'M_1X_2\hat\beta_2}_{(2)}
  + \underbrace{X_2'M_1M_Xy}_{(3)}
\end{equation}
Two properties eliminate terms (1) and (3):
\begin{enumerate}
  \item $M_1X_1 = 0$ $\Rightarrow$ term (1) $= 0$.
  \item $M_1M_X = M_X$ (since $M_X < M_1$) $\Rightarrow$
    $X_2'M_1M_Xy = X_2'M_Xy = X_2'\hat{\mathbf{e}} = 0$.
\end{enumerate}

\begin{tcolorbox}[keyresult]
\begin{equation}
  \hat\beta_2 = (X_2'M_1X_2)^{-1}X_2'M_1y,
  \qquad
  \Var(\hat\beta_2) = \sigma^2(X_2'M_1X_2)^{-1}
  \label{eq_FWL}
\end{equation}
\end{tcolorbox}

\noindent\textbf{Interpretation.} The FWL theorem says that the coefficient
on $X_2$ in the full regression is numerically identical to the coefficient
from regressing the residuals of $Y$ on $X_1$ onto the residuals of $X_2$
on $X_1$. In other words, OLS \emph{partials out} all other regressors before
estimating any particular coefficient. A regression coefficient does not
measure the raw correlation between $Y$ and $X_2$; it measures the
correlation after removing everything both $Y$ and $X_2$ share with $X_1$.


\section{Specification Issues}
\label{sec_Specification}

\subsection{Overspecification}

Variables not belonging to the true DGP are mistakenly included.
Overspecification is \textbf{not} misspecification.

Suppose we estimate $y = X\beta + Z\gamma + \mathbf{e}$,
$\mathbf{e}\sim\text{iid}(0,\sigma^2I)$, but the true DGP is
$y = X\beta_0 + \mathbf{e}$. The overspecified model nests the true model
(setting $\gamma=0$), so $\hat\beta$ is unbiased.

\begin{tcolorbox}[keyresult]
\textbf{Overspecification:} unbiased, (weakly) consistent, nominally
inefficient estimates.
\end{tcolorbox}

\subsection{Underspecification and Omitted Variable Bias}
\label{subsec_OVB}

We estimate $y = X\beta + \mathbf{e}$ but the true DGP is:
\begin{equation}
  y = X\beta_0 + Z\gamma_0 + \mathbf{e},
  \quad \mathbf{e}\sim\text{iid}(0,\sigma_0^2I)
\end{equation}

\noindent\textbf{Omitted variable bias:}\index{omitted variable bias}
\begin{align}
  \E(\hat\beta)
  &= \beta_0 + \underbrace{(X'X)^{-1}X'Z\gamma_0}_{\text{OVB}}
\end{align}

\begin{tcolorbox}[keyresult]
The OVB is zero if and only if:
\begin{enumerate}
  \item $X'Z = 0$ (omitted variable is orthogonal to included regressors), or
  \item $\gamma_0 = 0$ (the omitted variable has no effect on $y$).
\end{enumerate}
\end{tcolorbox}

\begin{tcolorbox}[keyresult]
\textbf{Underspecification:} biased, consistent, linear estimator.
\end{tcolorbox}


\section{Functional Form and Nonlinearity}
\label{sec_FunctionalForm}

The linear model requires linearity in \emph{parameters}, not in variables.
This distinction permits substantial flexibility.

\subsection{Polynomial Terms}

\begin{equation}
  y = \beta_0 + \beta_1 x + \beta_2 x^2 + e
\end{equation}

Substituting $z = x^2$ preserves linearity in the parameters while capturing
a nonlinear relationship in $x$.

\subsection{Log-linearisation: The Cobb-Douglas Production Function}

The \textbf{Cobb-Douglas}\index{Cobb-Douglas production function} function:
\begin{equation}
  y = AK^\alpha L^\beta e, \qquad\text{with CES: }\alpha = 1-\beta
\end{equation}

Log-linearisation yields a standard linear regression model:
\begin{equation}
  \ln y = \ln A + (1-\beta)\ln K + \beta\ln L + \ln e
\end{equation}

\noindent\textbf{Constant returns to scale} ($F(aK,aL) = aF(K,L)$):
\begin{equation}
  F(aK,aL) = A(aK)^\alpha(aL)^{1-\alpha} = a\,AK^\alpha L^{1-\alpha}
  \;\checkmark
\end{equation}
If $\alpha + \beta > 1$: increasing returns to scale.

\noindent\textbf{Elasticity:}
\begin{equation}
  \frac{dy/dL}{y/L}
  = \frac{\text{marginal product}}{\text{average product}}
  = \beta \quad\text{(constant factor proportion)}
\end{equation}

\noindent\textbf{Derivation from CES.} The Cobb-Douglas arises as a limiting
case of the CES function $Y = A[\alpha K^\rho + (1-\alpha)L^\rho]^{\gamma/\rho}$
as $\rho \to 0$. Applying L'H\^opital's rule to $\ln Y$:
\begin{equation}
  \lim_{\rho\to0}\ln Y = \ln A + \alpha\ln K + (1-\alpha)\ln L
  \;\Longrightarrow\; Y = AK^\alpha L^{1-\alpha}
\end{equation}


\section{Maximum Likelihood Estimation and OLS}
\label{sec_MLE}

The \textbf{classical normal regression} (CNR) model adds the assumption that
$e_i \mid \mathbf{x}_i \sim \mathcal{N}(0, \sigma^2)$, implying:
\begin{equation}
  y_i \mid \mathbf{x}_i \sim \mathcal{N}(\mathbf{x}_i'\beta, \sigma^2)
\end{equation}

The likelihood function for observation $i$:
\begin{equation}
  \mathcal{L}_i = f(y_i) = \frac{1}{\sqrt{2\pi\sigma^2}}
  \exp\!\left(-\frac{(y_i - \mathbf{x}_i'\beta)^2}{2\sigma^2}\right)
\end{equation}

The \textbf{maximum likelihood estimator}\index{maximum likelihood
estimator (MLE)} $(\hat\beta_{mle}, \hat\sigma^2_{mle})$ maximises
$\log\mathcal{L}(\beta,\sigma^2)$. Since this is a function of $\beta$ only
through $\text{SSE}_n(\beta)$, maximising the likelihood is equivalent to
minimising the sum of squared errors:

\begin{tcolorbox}[keyresult]
\begin{equation}
  \hat\beta_{mle} = \hat\beta_{ols}
\end{equation}
\end{tcolorbox}

Under the normality assumption, OLS is not merely the best linear unbiased
estimator (BLUE by the Gauss-Markov theorem); it is the minimum variance
unbiased estimator among \emph{all} unbiased estimators.


\section{General Method of Moments}
\label{sec_GMM}

OLS exploits the moment condition $\E[\mathbf{x}_i e_i] = 0$.
The \textbf{General Method of Moments}\index{GMM|see{General Method of
Moments}}\index{General Method of Moments} (GMM) generalises this by
allowing for a vector of moment conditions $\E[\mathbf{z}_i g(\mathbf{y}_i,
\boldsymbol{\theta})] = \mathbf{0}$, where $\mathbf{z}_i$ may be instruments
distinct from regressors $\mathbf{x}_i$. OLS is a special case of GMM with
$\mathbf{z}_i = \mathbf{x}_i$; IV/2SLS is another special case with
$\mathbf{z}_i$ containing external instruments (Chapter~\ref{chap_BeyondOLS}).

% TODO: Expand GMM with J-test for over-identification

\renewcommand\bibname{Further reading}
\begin{subappendices}
\bibliographystyle{econometrica}
\bibliography{Literature/OverLord}
\IfFileExists{Literature/ClassNormReg.bbl}{\chapter[The CNR model]{The classical normal regression (CNR) model}\label{chap_ClassNormReg}

\section{The CNR model}
\subsection{Basic set-up}
\subsection{The model in matrix notation}

\section{Different approaches to OLS}
\subsection{Least squares estimator}

\subsection{Geometric approach to OLS}
\subsubsection{Projection matrices and their relationships}
Let the matrix $\boldsymbol{P_{X}}$ be defined as

\begin{eqnarray}
\boldsymbol{P_{X}} \equiv \boldsymbol{X}(\boldsymbol{X}'\boldsymbol{X})^{-1}\boldsymbol{X}'.\label{eqn_CNR_Px}
\end{eqnarray}

$\boldsymbol{P_{X}}$ belongs to the general class of so-called projection matrices which have a number of interesting properties. Most importantly, the projection matrix $\boldsymbol{P_{X}}$ is a square matrix of dimension $n \times n$ that gives a vector space projection from $\mathbb{R}^{n}$ to a given subspace. Projection matrices are said to be \emph{idempotent} which implies that $\boldsymbol{P_{X}}\boldsymbol{P_{X}}=\boldsymbol{P_{X}}$.

This important property is straightforward to demonstrate in the case of equation \eqref{eqn_CNR_Px}, where we can use the fact that $$(\boldsymbol{X}'\boldsymbol{X})^{-1}\boldsymbol{X}'\boldsymbol{X}=\boldsymbol{I}$$ in order to show that $$\boldsymbol{P_{X}}\boldsymbol{P_{X}}=\boldsymbol{X}(\boldsymbol{X}'\boldsymbol{X})^{-1}\boldsymbol{X}'\boldsymbol{X}(\boldsymbol{X}'
\boldsymbol{X})^{-1}\boldsymbol{X}'=\boldsymbol{X}(\boldsymbol{X}'\boldsymbol{X})^{-1}\boldsymbol{X}'=\boldsymbol{P_{X}}$$. Furthermore, projection matrices can generally be classified either as \emph{orthogonal} projection matrices or \emph{oblique} (non-orthogonal) projection matrices.

In addition to being square, a projection matrix is also a symmetric matrix if and only if the vector space projection is orthogonal. $\boldsymbol{P_{X}}$ is an orthogonal projection matrix that spans the space $\mathcal{S}(\boldsymbol{X})$; consequently, any point in $\boldsymbol{X}$ that is projected by $\boldsymbol{P_{X}}$ will be projected onto itself. This implies that $\boldsymbol{P_{X}}\boldsymbol{X} = \boldsymbol{X}$. Again, this is easily verified by using equation \eqref{eqn_CNR_Px} which yields $\boldsymbol{P_{X}}\boldsymbol{X}=\boldsymbol{X}(\boldsymbol{X}'\boldsymbol{X})^{-1}\boldsymbol{X}'\boldsymbol{X} = \boldsymbol{X}$.

$\boldsymbol{M_{X}}$ is the orthogonal complement to $\boldsymbol{P_{X}}$, such that $\boldsymbol{M_{X}} = (\boldsymbol{I} - \boldsymbol{P_{X}})$. Consequently,

\begin{eqnarray}
\boldsymbol{P_{X}}\boldsymbol{M_{X}} &=& \boldsymbol{P_{X}}(\boldsymbol{I} - \boldsymbol{P_{X}}) \nonumber \\
&=& \boldsymbol{P_{X}}-\boldsymbol{P_{X}}\boldsymbol{P_{X}} = \boldsymbol{P_{X}}-\boldsymbol{P_{X}}=\boldsymbol{0}.
\end{eqnarray}

Since $\boldsymbol{P_{X}}$ and $\boldsymbol{M_{X}}$ are orthogonal to each other, any vector in $\mathcal{S}(\boldsymbol{X})$ will either be projected onto itself (i.e. $\boldsymbol{P_{X}}\boldsymbol{X} = \boldsymbol{X}$) or it will be annihilated ($\boldsymbol{M_{X}}\boldsymbol{X} = \boldsymbol{0}$).

Now, suppose that the matrix $\boldsymbol{X}$ can be decomposed into the submatrices $\boldsymbol{X}_{1}$ and $\boldsymbol{X}_{2}$, such that $\boldsymbol{X} = [\boldsymbol{X}_{1}\hspace{1ex}\boldsymbol{X}_{2}]$. Analogous to $\boldsymbol{P_{X}}$, we can now define the projection matrix $\boldsymbol{P}_{1} = \boldsymbol{X}_{1}(\boldsymbol{X}_{1}'\boldsymbol{X}_{1})^{-1}\boldsymbol{X}_{1}'$. This yields the following important relationships:

\begin{eqnarray}
\boldsymbol{P_{X}}\boldsymbol{P}_{1}=\boldsymbol{P}_{1} \qquad &\text{since}& \qquad \boldsymbol{P}_{1}<\boldsymbol{P_{X}}\\
\boldsymbol{M_{X}}\boldsymbol{M}_{1}=\boldsymbol{M_{X}} \qquad &\text{since}& \qquad \boldsymbol{M}_{1}>\boldsymbol{M_{X}}
\end{eqnarray}


\subsubsection{Geometry of the Frisch-Waugh-Lovell theorem}
See \citet{Lovell2008} for a simple proof that does not rely on matrix algebra. More comprehensive treatment of the Frisch-Waugh-Lovell theorem (FWL) can be found in \citet{Davidson2004}.

\subsection{General Method of Moments (GMM)}

\subsection{Maximum Likelihood Estimation and OLS}
The maximum likelihood estimate (MLE) is that assumed value of an unknown population parameter which maximises the probability of obtaining the sample data that is actually observed.

The normal regression model is the linear regression model under the restriction that the error term $e_{i}$ is independent of $\boldsymbol{x}_{i}$ and has the distribution $(0, \sigma^{2})$ Thus the standard set-up of OLS that we have just seen can easily be re-written as

\begin{eqnarray}
e_{i} | \boldsymbol{x}_{i} \sim (0, \sigma^{2})
\end{eqnarray}

Recall from section~\ref{chap_Review:Stats} that the mean and variance of a variable $V = \mu_{0} + e$ where $e \sim (0,\sigma^{2})$, is $\mu_{V} = \mu$ and $\sigma^{2}_{V}=\sigma^{2}_{e}$. This implies that

\begin{eqnarray}
y_{i} | \boldsymbol{x}_{i} \sim (\boldsymbol{x}_{i}'\boldsymbol{\beta}, \sigma^{2})
\end{eqnarray}

Equivalently, we can rewrite the linear model as $\boldsymbol{y} \sim N{\boldsymbol{\Omega}, \boldsymbol{\Omega}}$ where $\boldsymbol{\Omega}$ is assumed to meet the Gauss-Markov assumptions.

Given the parametric nature of the CNR, we can use likelihood methods\index{likelihood methods} for estimating the parameters. The likelihood function for each observation is

\begin{eqnarray}
\mathcal{L}_{i}  = f(y_{i}) = \frac{1}{\sqrt{2 \pi \sigma^{2}}}e^{-\frac{(y_{i} - \boldsymbol{x}'_{i}\boldsymbol{\beta})^{2}}{2\sigma^{2}}}
\end{eqnarray}

The maximum likelihood estimator\index{maximum likelihood estimator (MLE)|see{likelihood methods}} (MLE; $\hat{\boldsymbol{\beta}}_{mle}, \sigma^{2}_{mle}$) maximises $\log \mathcal{L}(\boldsymbol{\beta}, \sigma^{2})$. Since this is a function of $\boldsymbol{\beta}$ only through the sum of squared errors $SSE_{n}(\boldsymbol{\beta})$, maximising the likelihood is equivalent to minimising the sum of squared errors. Therefore

\begin{eqnarray}
\hat{\boldsymbol{\beta}}_{mle} = \boldsymbol{\hat{\beta}}_{ols}.
\end{eqnarray}


\begin{subappendices}
\section{Code for Spatial Stats in R}
\lstinputlisting{"Programs/UP504_OLS.R"}

\renewcommand\bibname{Further reading}

\bibliographystyle{econometrica}
\bibliography{Literature/OverLord}
\IfFileExists{Literature/ClassNormReg.bbl}{\chapter[The CNR model]{The classical normal regression (CNR) model}\label{chap_ClassNormReg}

\section{The CNR model}
\subsection{Basic set-up}
\subsection{The model in matrix notation}

\section{Different approaches to OLS}
\subsection{Least squares estimator}

\subsection{Geometric approach to OLS}
\subsubsection{Projection matrices and their relationships}
Let the matrix $\boldsymbol{P_{X}}$ be defined as

\begin{eqnarray}
\boldsymbol{P_{X}} \equiv \boldsymbol{X}(\boldsymbol{X}'\boldsymbol{X})^{-1}\boldsymbol{X}'.\label{eqn_CNR_Px}
\end{eqnarray}

$\boldsymbol{P_{X}}$ belongs to the general class of so-called projection matrices which have a number of interesting properties. Most importantly, the projection matrix $\boldsymbol{P_{X}}$ is a square matrix of dimension $n \times n$ that gives a vector space projection from $\mathbb{R}^{n}$ to a given subspace. Projection matrices are said to be \emph{idempotent} which implies that $\boldsymbol{P_{X}}\boldsymbol{P_{X}}=\boldsymbol{P_{X}}$.

This important property is straightforward to demonstrate in the case of equation \eqref{eqn_CNR_Px}, where we can use the fact that $$(\boldsymbol{X}'\boldsymbol{X})^{-1}\boldsymbol{X}'\boldsymbol{X}=\boldsymbol{I}$$ in order to show that $$\boldsymbol{P_{X}}\boldsymbol{P_{X}}=\boldsymbol{X}(\boldsymbol{X}'\boldsymbol{X})^{-1}\boldsymbol{X}'\boldsymbol{X}(\boldsymbol{X}'
\boldsymbol{X})^{-1}\boldsymbol{X}'=\boldsymbol{X}(\boldsymbol{X}'\boldsymbol{X})^{-1}\boldsymbol{X}'=\boldsymbol{P_{X}}$$. Furthermore, projection matrices can generally be classified either as \emph{orthogonal} projection matrices or \emph{oblique} (non-orthogonal) projection matrices.

In addition to being square, a projection matrix is also a symmetric matrix if and only if the vector space projection is orthogonal. $\boldsymbol{P_{X}}$ is an orthogonal projection matrix that spans the space $\mathcal{S}(\boldsymbol{X})$; consequently, any point in $\boldsymbol{X}$ that is projected by $\boldsymbol{P_{X}}$ will be projected onto itself. This implies that $\boldsymbol{P_{X}}\boldsymbol{X} = \boldsymbol{X}$. Again, this is easily verified by using equation \eqref{eqn_CNR_Px} which yields $\boldsymbol{P_{X}}\boldsymbol{X}=\boldsymbol{X}(\boldsymbol{X}'\boldsymbol{X})^{-1}\boldsymbol{X}'\boldsymbol{X} = \boldsymbol{X}$.

$\boldsymbol{M_{X}}$ is the orthogonal complement to $\boldsymbol{P_{X}}$, such that $\boldsymbol{M_{X}} = (\boldsymbol{I} - \boldsymbol{P_{X}})$. Consequently,

\begin{eqnarray}
\boldsymbol{P_{X}}\boldsymbol{M_{X}} &=& \boldsymbol{P_{X}}(\boldsymbol{I} - \boldsymbol{P_{X}}) \nonumber \\
&=& \boldsymbol{P_{X}}-\boldsymbol{P_{X}}\boldsymbol{P_{X}} = \boldsymbol{P_{X}}-\boldsymbol{P_{X}}=\boldsymbol{0}.
\end{eqnarray}

Since $\boldsymbol{P_{X}}$ and $\boldsymbol{M_{X}}$ are orthogonal to each other, any vector in $\mathcal{S}(\boldsymbol{X})$ will either be projected onto itself (i.e. $\boldsymbol{P_{X}}\boldsymbol{X} = \boldsymbol{X}$) or it will be annihilated ($\boldsymbol{M_{X}}\boldsymbol{X} = \boldsymbol{0}$).

Now, suppose that the matrix $\boldsymbol{X}$ can be decomposed into the submatrices $\boldsymbol{X}_{1}$ and $\boldsymbol{X}_{2}$, such that $\boldsymbol{X} = [\boldsymbol{X}_{1}\hspace{1ex}\boldsymbol{X}_{2}]$. Analogous to $\boldsymbol{P_{X}}$, we can now define the projection matrix $\boldsymbol{P}_{1} = \boldsymbol{X}_{1}(\boldsymbol{X}_{1}'\boldsymbol{X}_{1})^{-1}\boldsymbol{X}_{1}'$. This yields the following important relationships:

\begin{eqnarray}
\boldsymbol{P_{X}}\boldsymbol{P}_{1}=\boldsymbol{P}_{1} \qquad &\text{since}& \qquad \boldsymbol{P}_{1}<\boldsymbol{P_{X}}\\
\boldsymbol{M_{X}}\boldsymbol{M}_{1}=\boldsymbol{M_{X}} \qquad &\text{since}& \qquad \boldsymbol{M}_{1}>\boldsymbol{M_{X}}
\end{eqnarray}


\subsubsection{Geometry of the Frisch-Waugh-Lovell theorem}
See \citet{Lovell2008} for a simple proof that does not rely on matrix algebra. More comprehensive treatment of the Frisch-Waugh-Lovell theorem (FWL) can be found in \citet{Davidson2004}.

\subsection{General Method of Moments (GMM)}

\subsection{Maximum Likelihood Estimation and OLS}
The maximum likelihood estimate (MLE) is that assumed value of an unknown population parameter which maximises the probability of obtaining the sample data that is actually observed.

The normal regression model is the linear regression model under the restriction that the error term $e_{i}$ is independent of $\boldsymbol{x}_{i}$ and has the distribution $(0, \sigma^{2})$ Thus the standard set-up of OLS that we have just seen can easily be re-written as

\begin{eqnarray}
e_{i} | \boldsymbol{x}_{i} \sim (0, \sigma^{2})
\end{eqnarray}

Recall from section~\ref{chap_Review:Stats} that the mean and variance of a variable $V = \mu_{0} + e$ where $e \sim (0,\sigma^{2})$, is $\mu_{V} = \mu$ and $\sigma^{2}_{V}=\sigma^{2}_{e}$. This implies that

\begin{eqnarray}
y_{i} | \boldsymbol{x}_{i} \sim (\boldsymbol{x}_{i}'\boldsymbol{\beta}, \sigma^{2})
\end{eqnarray}

Equivalently, we can rewrite the linear model as $\boldsymbol{y} \sim N{\boldsymbol{\Omega}, \boldsymbol{\Omega}}$ where $\boldsymbol{\Omega}$ is assumed to meet the Gauss-Markov assumptions.

Given the parametric nature of the CNR, we can use likelihood methods\index{likelihood methods} for estimating the parameters. The likelihood function for each observation is

\begin{eqnarray}
\mathcal{L}_{i}  = f(y_{i}) = \frac{1}{\sqrt{2 \pi \sigma^{2}}}e^{-\frac{(y_{i} - \boldsymbol{x}'_{i}\boldsymbol{\beta})^{2}}{2\sigma^{2}}}
\end{eqnarray}

The maximum likelihood estimator\index{maximum likelihood estimator (MLE)|see{likelihood methods}} (MLE; $\hat{\boldsymbol{\beta}}_{mle}, \sigma^{2}_{mle}$) maximises $\log \mathcal{L}(\boldsymbol{\beta}, \sigma^{2})$. Since this is a function of $\boldsymbol{\beta}$ only through the sum of squared errors $SSE_{n}(\boldsymbol{\beta})$, maximising the likelihood is equivalent to minimising the sum of squared errors. Therefore

\begin{eqnarray}
\hat{\boldsymbol{\beta}}_{mle} = \boldsymbol{\hat{\beta}}_{ols}.
\end{eqnarray}


\begin{subappendices}
\section{Code for Spatial Stats in R}
\lstinputlisting{"Programs/UP504_OLS.R"}

\renewcommand\bibname{Further reading}

\bibliographystyle{econometrica}
\bibliography{Literature/OverLord}
\IfFileExists{Literature/ClassNormReg.bbl}{\input{Literature/ClassNormReg.bbl}}{\typeout{Need to run bibtex}}
\end{subappendices}
}{\typeout{Need to run bibtex}}
\end{subappendices}
}{\typeout{Need bibtex on ClassNormReg}}
\end{subappendices}

% =====================================================================
%  Chapter 5: Causal Inference and Identification
%  Source: P2 transcription (complete)
% =====================================================================

\chapter{Causal Inference and Identification}
\label{chap_Causal}

\section{The Goals of Empirical Research}
\label{sec_Goals}

Not all empirical questions are causal. \citet{Angrist2009a} organise the
goals of applied econometric work into a hierarchy:

\begin{figure}[htp]
\centering
\begin{tikzpicture}[>=Stealth,
    node distance=1.2cm and 2.2cm,
    every node/.style={draw, rounded corners, align=center, font=\small,
                       minimum width=2.8cm, minimum height=0.65cm}]
  % Root
  \node (root) {Empirical research};
  % Level 1
  \node (sum)  [below left=of root,  xshift=0.5cm]  {Summarising data};
  \node (inf)  [below right=of root, xshift=-0.5cm] {Inference};
  \draw[->] (root) -- (sum);
  \draw[->] (root) -- (inf);
  % Level 2 under Inference
  \node (desc) [below left=of inf,  xshift=0.3cm]  {Descriptive};
  \node (cf)   [below right=of inf, xshift=-0.3cm] {Counterfactual};
  \draw[->] (inf) -- (desc);
  \draw[->] (inf) -- (cf);
  % Level 3 under Counterfactual
  \node (pred) [below left=of cf,   xshift=0.5cm]  {Prediction\\(what-if)};
  \node (caus) [below right=of cf,  xshift=-0.5cm] {Causal\\identification};
  \draw[->] (cf) -- (pred);
  \draw[->] (cf) -- (caus);
\end{tikzpicture}
\caption{Taxonomy of the goals of empirical research.}
\label{fig_EmpiricalGoals}
\end{figure}

The distinction between descriptive inference and causal identification is
central. A descriptive regression summarises conditional means:
$\E[Y\mid X=x]$ tells us what values of $Y$ tend to accompany given values
of $X$, but nothing about what would happen to $Y$ if we \emph{changed} $X$.
Causal identification requires an additional assumption---a source of
exogenous variation that allows us to move from correlation to causation.

\subsection{Standard Regression Notation}

We can write the regression model in several equivalent ways:
\begin{alignat*}{2}
  Y_i &= m(\mathbf{x}) + \varepsilon && \\
      &= X_i\beta + \varepsilon_i && \quad\text{(standard form)} \\
  Y_i &\sim f_N(y_i \mid \mu_i, \sigma^2), \quad \mu_i = X_i\beta
    && \quad\text{(distributional form)}
\end{alignat*}

More generally, for a response variable $y_i$ following distribution
$f(y_i \mid \theta_i, \alpha)$ with systematic component $\theta_i =
g(X_i, \beta)$:

\begin{table}[htp]
\centering
\caption{Generalised regression notation across model families}
\label{tbl_RegNotation}
\renewcommand{\arraystretch}{1.4}
\begin{tabular}{@{}lllll@{}}
\toprule
\textbf{Model} & \textbf{Distribution} $f$ & \textbf{Link} $g$ &
\textbf{Range of $y$} & \textbf{Estimator} \\
\midrule
OLS          & Normal    & Identity  & $(-\infty, +\infty)$ & OLS / MLE \\
Logit        & Bernoulli & Logistic  & $\{0,1\}$             & MLE \\
Probit       & Bernoulli & Normal CDF & $\{0,1\}$            & MLE \\
Poisson      & Poisson   & Log       & $\{0,1,2,\ldots\}$   & MLE \\
Neg.\ Binomial & NB      & Log       & $\{0,1,2,\ldots\}$   & MLE \\
Gamma        & Gamma     & Log / Inv & $(0,+\infty)$         & MLE \\
Tobit        & Normal    & Identity  & $[0,+\infty)$         & MLE \\
Spatial Lag  & Normal    & Identity  & $(-\infty,+\infty)$   & MLE / IV \\
\bottomrule
\end{tabular}
\renewcommand{\arraystretch}{1}
\end{table}

\subsection{Two Types of Uncertainty}

\begin{itemize}
  \item \textbf{Estimation uncertainty:} arises because we observe a finite
    sample; standard errors shrink as $n \to \infty$.
  \item \textbf{Fundamental uncertainty:} arises because $Y$ is inherently
    stochastic; this uncertainty does not shrink with $n$. It is reflected in
    the residual variance $\sigma^2$.
\end{itemize}

\begin{figure}[htp]
\centering
\begin{tikzpicture}[>=Stealth, scale=1.0]
  % Three panels: linear, logistic, exponential systematic components
  % Panel 1: Linear
  \begin{scope}[xshift=0cm]
    \draw[->] (-0.2,0) -- (2.8,0) node[right] {\small$x$};
    \draw[->] (0,-0.2) -- (0,2.4) node[above] {\small$\mu$};
    \draw[very thick] (0,0.3) -- (2.5,2.2);
    \node[below] at (1.25,-0.3) {\small Linear};
  \end{scope}
  % Panel 2: Logistic
  \begin{scope}[xshift=4.0cm]
    \draw[->] (-0.2,0) -- (2.8,0) node[right] {\small$x$};
    \draw[->] (0,-0.2) -- (0,2.4) node[above] {\small$\mu$};
    \draw[very thick, domain=0:2.5, samples=60]
      plot (\x, {2.1/(1+exp(-4*(\x-1.25)))});
    \draw[dashed, gray] (-0.1, 2.1) -- (2.7, 2.1);
    \node[below] at (1.25,-0.3) {\small Logistic};
  \end{scope}
  % Panel 3: Exponential
  \begin{scope}[xshift=8.0cm]
    \draw[->] (-0.2,0) -- (2.8,0) node[right] {\small$x$};
    \draw[->] (0,-0.2) -- (0,2.4) node[above] {\small$\mu$};
    \draw[very thick, domain=0:2.3, samples=60]
      plot (\x, {0.3*exp(1.0*\x)});
    \node[below] at (1.25,-0.3) {\small Exponential};
  \end{scope}
\end{tikzpicture}
\caption{Three systematic component forms: linear (identity link), logistic
(logit link), and exponential (log link).}
\label{fig_SystematicComp}
\end{figure}


\section{The Fundamental Problem of Causal Inference}
\label{sec_FPCI}

\subsection{Potential Outcomes and the Rubin Causal Model}

Let $D_i \in \{0,1\}$ be a binary treatment indicator. Each unit $i$ has
two \textbf{potential outcomes}\index{potential outcomes}:
\begin{itemize}
  \item $Y_{0i}$: outcome if untreated ($D_i = 0$)
  \item $Y_{1i}$: outcome if treated ($D_i = 1$)
\end{itemize}

The observed outcome is (\citealt{Holland1986}; Rubin Causal Model):

\begin{tcolorbox}[keyresult]
\begin{equation}
  Y_i = Y_{0i} + (Y_{1i} - Y_{0i})D_i
  \label{eq_PotentialOutcomes}
\end{equation}
\end{tcolorbox}

The \textbf{individual treatment effect} is $Y_{1i} - Y_{0i}$. The
\textbf{fundamental problem of causal inference} is that we observe only one
potential outcome for each unit---either $Y_{0i}$ or $Y_{1i}$, never both.

\subsection{The Average Treatment Effect on the Treated (ATT)}

\begin{tcolorbox}[keyresult]
\begin{align}
  \underbrace{\E[Y_i \mid D_i=1] - \E[Y_i \mid D_i=0]}_{\text{observed difference}}
  &= \underbrace{\E[Y_{1i} - Y_{0i} \mid D_i=1]}_{\text{ATT}}
   + \underbrace{\E[Y_{0i} \mid D_i=1] - \E[Y_{0i} \mid D_i=0]}_{\text{selection bias}}
  \label{eq_ATT}
\end{align}
\end{tcolorbox}

The observed difference between treated and untreated units confounds the
causal effect (ATT) with selection bias---the difference in potential
untreated outcomes between those who were and were not treated.


\section{The Selection Problem in Urban Contexts}
\label{sec_Selection}

\subsection{The Hospital Paradox}

A striking illustration: patients in hospitals have worse health outcomes than
patients not in hospitals. Does this mean hospitals are harmful? No---sicker
people are more likely to \emph{select into} hospital treatment. The
selection mechanism determines who receives treatment, creating the
confounding. The same logic applies throughout urban and regional analysis:
neighbourhoods, cities, and regions are not random assignments.

\begin{table}[htp]
\centering
\caption{The selection problem: conventional estimates vs.\ corrected estimates}
\label{tbl_Selection1}
\renewcommand{\arraystretch}{1.4}
\begin{tabular}{@{}lcc@{}}
\toprule
& \textbf{Treated (hospital)} & \textbf{Untreated (no hospital)} \\
\midrule
Average health outcome       & Low    & High \\
Na\"{i}ve treatment effect   & \multicolumn{2}{c}{Negative (spurious)} \\
True treatment effect (ATT)  & \multicolumn{2}{c}{Positive (causal)} \\
Selection bias               & \multicolumn{2}{c}{Negative (sicker select in)} \\
\bottomrule
\end{tabular}
\renewcommand{\arraystretch}{1}
\end{table}

\begin{table}[htp]
\centering
\caption{REO (real estate owned) vs.\ conventional property sales:
a selection problem in housing markets}
\label{tbl_Selection2}
\renewcommand{\arraystretch}{1.4}
\begin{tabular}{@{}lcc@{}}
\toprule
Property characteristic & REO sale & Conventional sale \\
\midrule
Seller motivation       & Distressed (bank) & Voluntary \\
Condition at sale       & Often degraded    & Market-maintained \\
Observed sale price     & Discounted        & Full market value \\
Na\"{i}ve ``REO discount'' & Negative      & Baseline \\
True REO effect         & \multicolumn{2}{c}{Requires identification strategy} \\
\bottomrule
\end{tabular}
\renewcommand{\arraystretch}{1}
\end{table}

\subsection{Geographic Sorting and the Endogeneity of Place}

Cities and neighbourhoods are not random assignments---they are the outcomes
of sorting by households, firms, and governments over decades and centuries.
When we observe that neighbourhood $A$ has better outcomes than neighbourhood
$B$, we cannot distinguish the causal effect of neighbourhood characteristics
from the selection process that determined who lives where. The standard
spatial regression coefficient conflates both. This is the core identification
challenge of regional policy informatics.


\section{The Experimental Ideal}
\label{sec_Experimental}

\subsection{Random Assignment}

Under random assignment, $D_i$ is independent of potential outcomes:
\begin{equation}
  (Y_{0i}, Y_{1i}) \ind D_i
\end{equation}
This implies $\E[Y_{0i}\mid D_i=1] = \E[Y_{0i}\mid D_i=0]$, so the
selection bias term in equation~\eqref{eq_ATT} vanishes:
\begin{equation}
  \E[Y_i \mid D_i=1] - \E[Y_i \mid D_i=0] = \text{ATT}
  = \E[Y_{1i} - Y_{0i} \mid D_i=1]
\end{equation}

Note: randomisation \emph{spreads} selection bias symmetrically across
treatment and control groups---it does not remove selection bias from
individual units but makes it zero \emph{in expectation}.

\paragraph{Example: The Perry Pre-School Programme (Heckman et al., 2010)}
\index{Perry Pre-School}

The Perry Pre-School Programme (1962) was a randomised controlled trial
(RCT) that randomly assigned disadvantaged African-American children in
Ypsilanti, Michigan to either a high-quality pre-school programme (treatment)
or a control group. Follow-up studies by \citet{Heckman2010} traced
participants into adulthood, finding substantial long-run effects on education,
employment, criminal behaviour, and earnings---effects that far exceeded the
short-run IQ gains that were the original target of the programme. The Perry
RCT remains one of the most influential examples of a field experiment in
social science, demonstrating both the power of random assignment and the
importance of measuring long-run outcomes.

\subsection{Regression with Constant Treatment Effect}

Under random assignment and assuming constant treatment effects:

\begin{tcolorbox}[keyresult]
\begin{equation}
  Y_i = \alpha + \beta D_i + \eta_i
  \label{eq_TreatmentReg}
\end{equation}
where $\beta = \E[Y_{1i} - Y_{0i}]$ is the Average Treatment Effect (ATE)
and $\alpha = \E[Y_{0i}]$ is the mean counterfactual outcome.
\end{tcolorbox}


\section{The Conditional Independence Assumption}
\label{sec_CIA}

When randomisation is not available, causal identification typically requires
one of three strategies: CIA (this section), instrumental variables
(Chapter~\ref{chap_BeyondOLS}), or quasi-experimental designs
(Chapter~\ref{chap_BeyondOLS}).

\begin{tcolorbox}[keyresult]
\textbf{Conditional Independence Assumption (CIA):}
\begin{equation}
  \{Y_{0i}, Y_{1i}\} \ind D_i \mid X_i
  \label{eq_CIA}
\end{equation}
Selection into treatment is independent of potential outcomes \emph{after
conditioning on observable covariates} $X_i$.
\end{tcolorbox}

Under the CIA, the regression with control variables $X_i$:
\begin{equation}
  Y_i = \alpha + \beta D_i + X_i\gamma' + \eta_i
\end{equation}
\emph{has a causal interpretation}---$\beta$ estimates the average treatment
effect conditional on $X_i$.

\noindent\textbf{Limitations of the CIA.} The CIA requires that we have
observed and included all confounders---variables that affect both treatment
assignment and outcomes. In the presence of \emph{unobserved} confounders,
the CIA fails and regression estimates are biased. This is the case in most
spatial settings, where geographic sorting creates selection on unobservables
(amenity preferences, ability, social capital) that are not captured in
administrative data.


\section{The TSS Decomposition and $R^2$}
\label{sec_TSS}

The decomposition of total variation into explained and residual components
follows directly from the orthogonality of $\hat{\mathbf{y}}$ and
$\hat{\mathbf{e}}$:

\begin{align*}
  \sum(y_i - \bar{y})^2
  &= \sum(\hat{y}_i - \bar{y})^2 + \sum(y_i - \hat{y}_i)^2 \\
  \text{TSS} &= \text{ESS} + \text{RSS}
\end{align*}

\noindent\textit{Proof:}
\begin{align*}
  \sum(y_i - \bar{y})^2
  &= \sum\bigl[(\hat{y}_i-\bar{y}) + \hat{e}_i\bigr]^2 \\
  &= \sum(\hat{y}_i-\bar{y})^2 + 2\sum(\hat{y}_i-\bar{y})\hat{e}_i
     + \sum\hat{e}_i^2 \\
  &= \sum(\hat{y}_i-\bar{y})^2 + \sum\hat{e}_i^2
\end{align*}
where the cross term vanishes because $\hat{\mathbf{y}}'\hat{\mathbf{e}} = 0$
and $\boldsymbol{\iota}'\hat{\mathbf{e}} = 0$ (the constant is always
included, so the residuals sum to zero). $\square$

\renewcommand\bibname{Further reading}
\begin{subappendices}
\bibliographystyle{econometrica}
\bibliography{Literature/OverLord}
\IfFileExists{Literature/Causal.bbl}{\input{Literature/Causal.bbl}}{\typeout{Need bibtex on Causal}}
\end{subappendices}

% =====================================================================
%  Chapter 6: Beyond OLS — The Identification Toolkit
%  Source: BeyondOLS.tex (stub) + P3 transcription (IV, heterosk., GLS)
% =====================================================================

\chapter[Beyond OLS]{Beyond OLS: The Identification Toolkit}
\label{chap_BeyondOLS}

\section{Endogeneity: Sources and Consequences}
\label{sec_Endogeneity}

OLS is inconsistent when the key assumption $\E[X'\mathbf{e}] = 0$ fails.
This \textbf{endogeneity}\index{endogeneity} can arise from three distinct
sources.

\subsection{Omitted Variable Bias Revisited}

As shown in Section~\ref{subsec_OVB}, failing to include a variable $Z$ that
is correlated with both $X$ and $Y$ biases $\hat\beta$. The direction of bias
is determined by the sign of $(X'X)^{-1}X'Z\gamma_0$---the product of the
regression of the omitted variable on the included regressors and the effect
of the omitted variable on $Y$.

\subsection{Measurement Error}

True regressors $X^*$ are not observed; only $X = X^* + u$ is available,
where $u$ is measurement error independent of $X^*$ and $\mathbf{y}$.
Substituting:
\begin{equation}
  \mathbf{y} = X^*\beta + \mathbf{e} = (X-u)\beta + \mathbf{e}
             = X\beta + (\mathbf{e} - u\beta) = X\beta + \mathbf{v}
\end{equation}

The bias in OLS arises from $\E[X'\mathbf{v}]$:
\begin{align}
  \E[X'\mathbf{v}]
  &= \E\!\bigl[(X^*+u)'(\mathbf{e}-u\beta)\bigr] \nonumber \\
  &= \E[X^{*\prime}\mathbf{e}] + \E[u'\mathbf{e}]
     - \E[X^{*\prime}u]\beta - \E[u'u]\beta \nonumber \\
  &= -\E[u'u]\,\beta
\end{align}

This is non-zero whenever $\beta \neq 0$ and $\Var(u) > 0$, so OLS is
inconsistent. The \textbf{attenuation bias}\index{attenuation bias} pulls
$\hat\beta$ toward zero:
\begin{equation}
  \hat\beta \xrightarrow{p} \beta^* = \beta - \E\!\bigl[(X'X)^{-1}u'u\bigr]\beta
\end{equation}

\subsection{Simultaneous Equations}

When $X$ and $Y$ are jointly determined---as in a supply-and-demand system
or in a spatial model where $y_i$ and $y_j$ are mutually dependent---OLS
conflates the structural parameters of multiple equations. For example:
\begin{align*}
  q_i &= -\beta_1 p_i + e_{1i} \quad\text{(demand)} \\
  q_i &= \beta_2 p_i + e_{2i} \quad\text{(supply)}
\end{align*}
OLS on either equation alone recovers a mixture of both structural
coefficients. In each case $\hat\beta$ is \textbf{biased and
inconsistent}:\; $\partial y/\partial x = \beta + \partial e/\partial x
\neq \beta$.


\section{Instrumental Variables}
\label{sec_IV}

\subsection{The IV Estimator}

An instrument $Z$ is valid if: (a) \textbf{relevance}---$Z$ is correlated
with $X$; (b) \textbf{exclusion restriction}---$Z$ affects $Y$ only through
$X$ (not directly).

\begin{tcolorbox}[keyresult]
\begin{equation}
  \hat\beta_{IV} = (Z'X)^{-1}Z'\mathbf{y}
  \label{eq_IV}
\end{equation}
\end{tcolorbox}

Classic example \citep{Angrist1991}: month of birth and distance from
college instrument for years of schooling in a wage equation, with unobserved
ability creating endogeneity of the schooling variable.

\subsection{Two-Stage Least Squares (2SLS)}

With endogenous regressors $X_2$ and instruments $Z_2$, where $X = [X_1\;X_2]$
and $Z = [X_1\;Z_2]$ ($X_1$ instruments for itself):

\begin{enumerate}
  \item \textbf{First stage:} Regress $X$ on $Z$:\quad
    $\hat\Gamma = (Z'Z)^{-1}Z'X$,\quad $\hat X = Z\hat\Gamma$
  \item \textbf{Second stage:} Regress $\mathbf{y}$ on $\hat X$:\quad
    $\hat\beta = (\hat X'\hat X)^{-1}\hat X'\mathbf{y}$
\end{enumerate}

Scrutinising $\hat X = [\hat X_1,\hat X_2] = [P_ZX_1, P_ZX_2]
= [X_1, P_ZX_2]$: \textbf{only endogenous variables are replaced with their
fitted values}, not exogenous ones. The 2SLS estimator is:
\begin{equation}
  \hat\beta_{2SLS} = (X'P_ZX)^{-1}X'P_Z\mathbf{y}
  = \bigl[X'Z(Z'Z)^{-1}Z'X\bigr]^{-1}X'Z(Z'Z)^{-1}Z'\mathbf{y}
\end{equation}

In the just-identified case (one instrument per endogenous variable),
$\hat\beta_{2SLS} = \hat\beta_{IV}$.

\begin{tcolorbox}[keyresult]
The IV/2SLS estimator is biased in finite samples but \textbf{consistent}
under instrument validity---in contrast to OLS which is inconsistent under
endogeneity.
\end{tcolorbox}

\subsection{Spatial Lags as Instruments}

In spatial models, the endogenous variable is the spatial lag $Wy$.
A common IV strategy uses higher-order spatial lags $W^2y$, $WX$, or $W^2X$
as instruments---variables that are correlated with $Wy$ through the spatial
propagation mechanism but excluded from the structural equation. This strategy
is formalised in the SAR-IV estimator (Chapter~\ref{chap_Spatial}).


\section{Differences-in-Differences}
\label{sec_DiD}

When panel data are available and a treatment is applied at a specific point
in time, \textbf{differences-in-differences}\index{differences-in-differences}
(DiD) removes time-invariant confounders:
\begin{equation}
  Y_{it} = \alpha_i + \lambda_t + \beta(D_i \times \text{Post}_t) + \varepsilon_{it}
\end{equation}
where $\alpha_i$ are unit fixed effects, $\lambda_t$ are time fixed effects,
and $D_i \times \text{Post}_t$ is the interaction of treatment group and
post-treatment period. $\beta$ is identified under the
\textbf{parallel trends assumption}: absent treatment, treated and control
units would have followed the same time trend.

\noindent\textbf{Geographic DiD.} Administrative boundaries (census tract
borders, school district lines, city limits) often create sharp discontinuities
in treatment assignment---properties just inside versus just outside a
rezoning boundary, for example. These boundaries can serve as regression
discontinuity (RD) designs that exploit spatial proximity to identify causal
effects. See \citet{Black1999} for a classic application to school quality
capitalisation.

% TODO: Add staggered DiD and Callaway-Sant'Anna correction


\section{Heteroskedasticity and Robust Inference}
\label{sec_Heterosk}

Under heteroskedasticity, $\E[e_i^2\mid\mathbf{x}_i] = \sigma_i^2$:
\begin{equation}
  y_i = X\beta + e_i, \quad
  e_i \sim \text{id}(0, \Omega), \quad
  \Omega \equiv \operatorname{diag}(\sigma_1^2,\ldots,\sigma_n^2)
\end{equation}

OLS is still \textbf{unbiased} because unbiasedness depends only on first
moments ($\E[\mathbf{e}]=0$). But the variance is no longer $\sigma^2(X'X)^{-1}$:
\begin{align}
  V(\hat\beta)
  &= (X'X)^{-1}X'\,\E[\mathbf{e}\mathbf{e}']\,X(X'X)^{-1} \nonumber \\
  &= (X'X)^{-1}X'\Omega X(X'X)^{-1}
  \label{eq_SandwichVar}
\end{align}

This is the \textbf{sandwich (HC) estimator}\index{sandwich estimator}. OLS
standard errors computed under the homoskedastic formula $s^2(X'X)^{-1}$ are
incorrect under heteroskedasticity---typically underestimating the true
standard errors.

\begin{tcolorbox}[keyresult]
Heteroskedasticity-consistent (HC) standard errors from equation
\eqref{eq_SandwichVar} are valid under heteroskedasticity of unknown form.
Use \texttt{vcovHC()} in R or \texttt{HC3} as default.
\end{tcolorbox}


\section{Generalised Least Squares}
\label{sec_GLS}

If $\Omega$ is known (or consistently estimated), we can eliminate
heteroskedasticity by pre-multiplying the model by $\Omega^{-1/2}$:

\begin{tcolorbox}[keyresult]
\begin{equation}
  \hat\beta_{GLS} = (X'\Omega^{-1}X)^{-1}X'\Omega^{-1}\mathbf{y}
\end{equation}
\end{tcolorbox}

GLS is the Best Linear Unbiased Estimator (BLUE) when $\Omega$ is known.
In practice $\Omega$ must be estimated (Feasible GLS), introducing estimation
error but retaining consistency.

\noindent\textbf{Spatial GLS.} In spatial models the error covariance
$\Omega$ has a specific parametric structure determined by the spatial weight
matrix $W$ and the error dependence parameter $\lambda$. The Spatial Error
Model (SEM, Chapter~\ref{chap_Spatial}) is a GLS model with
$\Omega = [(I-\lambda W)'(I-\lambda W)]^{-1}\sigma^2$.


\section{Panel Data Methods}
\label{sec_Panel}

Panel data follow the same units (individuals, firms, cities) over multiple
time periods. The key advantage is the ability to control for time-invariant
unobserved heterogeneity.

\subsection{Fixed Effects}

The \textbf{within estimator}\index{fixed effects estimator} removes
unit-specific means:
\begin{equation}
  (y_{it} - \bar{y}_i) = (x_{it} - \bar{x}_i)'\beta + (e_{it} - \bar{e}_i)
\end{equation}

This eliminates any unit-level time-invariant confounder
$\alpha_i$---whether or not it is observed. The fixed effects estimator is
consistent under the weaker assumption that $\E[x_{it}e_{it}] = 0$ (strict
exogeneity within unit), not $\E[x_{it}\alpha_i] = 0$.

\subsection{Random Effects}

The random effects estimator treats unit effects $\alpha_i$ as random
variables drawn from a distribution, rather than fixed parameters to be
estimated. It is more efficient than fixed effects but requires the stronger
assumption that $\alpha_i \ind X_{it}$---unit effects are uncorrelated with
regressors. The \textbf{Hausman test} discriminates between the two:
if fixed effects and random effects estimates differ significantly, the random
effects assumption is rejected.

\subsection{Two-Way Fixed Effects and DiD}

The two-way FE model adds time fixed effects $\lambda_t$:
\begin{equation}
  y_{it} = \alpha_i + \lambda_t + X_{it}'\beta + e_{it}
\end{equation}
Under parallel trends, this is numerically equivalent to the
differences-in-differences estimator described above. Recent work by
\citet{Callaway2021} shows that the standard two-way FE estimator can produce
misleading estimates when treatment adoption is staggered across units; they
propose a cohort-based correction.

% TODO: Add Callaway-Sant'Anna worked example with urban policy data

\renewcommand\bibname{Further reading}
\begin{subappendices}
\bibliographystyle{econometrica}
\bibliography{Literature/OverLord}
\IfFileExists{Literature/BeyondOLS.bbl}{\chapter{Beyond OLS}\label{chap_BeyondOLS}
\section{Measurement error and identification}
\section{Instrumental variables (IV), Two-Stage Least Squares (2SLS)}




\renewcommand\bibname{Further reading}
\begin{subappendices}
\bibliographystyle{econometrica}
\bibliography{Literature/OverLord}
\IfFileExists{Literature/BeyondOLS.bbl}{\chapter{Beyond OLS}\label{chap_BeyondOLS}
\section{Measurement error and identification}
\section{Instrumental variables (IV), Two-Stage Least Squares (2SLS)}




\renewcommand\bibname{Further reading}
\begin{subappendices}
\bibliographystyle{econometrica}
\bibliography{Literature/OverLord}
\IfFileExists{Literature/BeyondOLS.bbl}{\input{Literature/BeyondOLS.bbl}}{\typeout{Need to run bibtex}}
\end{subappendices}
}{\typeout{Need to run bibtex}}
\end{subappendices}
}{\typeout{Need bibtex on BeyondOLS}}
\end{subappendices}


% ── Part III: The Spatial Economy ─────────────────────────────────────
\part{The Spatial Economy}\label{part3}
% =====================================================================
%  Chapter 7: Spatial Equilibrium and Urban Structure
%  Source: SpatialStats.tex (intro prose) + new framework
%  NOTE: This chapter requires the most new writing. Sections marked
%        % TODO are high-priority development targets.
% =====================================================================

\chapter{Spatial Equilibrium and Urban Structure}
\label{chap_SpatialEquil}

Quantitative analysis of spatial data has advanced greatly in recent years,
largely as the result of the rapidly decreasing cost of computational power,
the increased availability of high-resolution data, and methodological
innovation in spatial econometrics. Indeed, recent analytical and
technological advances are increasingly blurring the traditional boundaries
in spatial data analysis where GIS\index{geographic information systems},
remote sensing\index{remote sensing} and spatial statistics\index{spatial
statistics} were essentially treated as separate fields
\citep[e.g.][]{Davis2002a,O'Sullivan2003,Campbell2006}.

Nowhere in the social sciences, perhaps, is this silent revolution felt more
than in urban and regional planning. While planners have usually emphasised
GIS applications for land-use and environmental planning, the discipline has
never embraced all three sub-fields of quantitative spatial data analysis in
a systematic way. With the rapid deployment of ``Big Data''\index{Big Data}
and the emergence of ``policy informatics''\index{policy informatics} as a new
paradigm among computational social scientists, planners should increasingly
see themselves at the forefront of integrated spatial data methods
\citep{Anselin2012,Barrett2011}.

\section{Why Spatial Equilibrium?}
\label{sec_WhySpatEq}

The \textbf{spatial equilibrium}\index{spatial equilibrium} framework provides
the theoretical spine that connects urban economics, regional science,
economic geography, and real estate into a single coherent argument. The
central insight, due to \citet{Rosen1979} and \citet{Roback1982}, is that
workers and firms sort across locations until utility and profits are equalised
at the margin. In equilibrium, observed geographic patterns of wages, rents,
and employment reflect underlying fundamentals---productivity differences,
amenity values, and the costs of transportation and housing.

Three implications follow:

\begin{enumerate}
  \item \textbf{Wages, rents, and quality of life are jointly determined.}
    High-amenity cities will have lower wages and higher rents, all else
    equal, because workers are willing to accept lower pay to live there.
    Observed wage or rent levels alone are not informative about the
    underlying distribution of productivity and amenity.

  \item \textbf{Observed patterns are equilibrium outcomes.} If we observe
    a spatial concentration of activity, it is because the benefits of
    concentration (agglomeration economies) outweigh the costs (congestion,
    land prices). Spatial econometrics tests whether and how adjustment
    toward equilibrium propagates across space.

  \item \textbf{Policy interventions change the equilibrium.} A place-based
    subsidy shifts the equilibrium allocation of workers and firms. The
    incidence of the subsidy---who captures the benefit---depends on the
    elasticities of housing supply and labour demand. Identifying these
    elasticities is a central task of regional policy informatics.
\end{enumerate}

% TODO: Add Rosen-Roback model derivation (wage-rent gradient)


\section{The Urban Land Market}
\label{sec_LandMarket}

\subsection{The von Thünen Model and Bid-Rent Gradients}

The earliest formal spatial equilibrium model is due to \citet{VonThunen1826}.
In von Thünen's model, a city (central market) is surrounded by an
agricultural hinterland. Farmers bid for land based on the profit from
agricultural production after transport costs. Land closer to the city
commands higher rents because transport costs are lower; the equilibrium
rent gradient is downward-sloping from the centre.

The generalisation to urban land is immediate: replace agricultural produce
with commuting costs. Workers bid for residential land based on their surplus
income after paying rent and commuting costs. The equilibrium bid-rent
function for land at distance $d$ from the city centre is:
\begin{equation}
  r(d) = r_0 - t \cdot d
\end{equation}
where $r_0$ is the rent at the city centre and $t$ is the commuting cost per
unit distance.

\subsection{The Monocentric City: Alonso-Muth-Mills}

The \textbf{Alonso-Muth-Mills model}\index{monocentric city model} formalises
the urban land market equilibrium with utility-maximising households,
profit-maximising developers, and a competitive land market
\citep{Alonso1964, Muth1969, Mills1967}. In equilibrium, the spatial
gradient of house prices compensates for the gradient in commuting costs;
household utility is equalised across all locations.

% TODO: Full AMM derivation with utility function and equilibrium conditions


\section{Agglomeration and the Concentration of Economic Activity}
\label{sec_Agglomeration}

\subsection{Why Do Cities Exist?}

The existence of cities requires an explanation: why do workers and firms
pay the premium of dense urban land when they could spread out? Three
categories of agglomeration economies\index{agglomeration economies} explain
the productivity premium of density \citep{Duranton2004}:

\begin{enumerate}
  \item \textbf{Sharing:} Dense markets allow firms to share indivisible
    inputs (specialised suppliers, infrastructure) and workers to share a
    thick labour market with lower search costs.
  \item \textbf{Matching:} Thick markets improve the quality of
    worker-firm matches, increasing productivity and wages. Cities are
    matching engines \citep{Helsley1990}.
  \item \textbf{Learning:} Knowledge spillovers diffuse more readily over
    short distances. Workers in dense environments learn from each other
    through informal contact \citep{Glaeser1992}.
\end{enumerate}

\subsection{The New Economic Geography}

\citet{Krugman1991} formalised the conditions under which economic activity
concentrates in a core region rather than distributing uniformly across space.
In Krugman's model, increasing returns to scale in manufacturing create
incentives to concentrate production; forward and backward linkages between
firms and workers amplify the initial concentration. The outcome depends on
the balance between \emph{centripetal forces} (scale economies, linkages)
and \emph{centrifugal forces} (high land prices, congestion, iceberg transport
costs).

% TODO: Add NEG bifurcation diagram and formal model sketch


\section{The Housing Market as a Spatial Equilibrium Price Surface}
\label{sec_HousingMarket}

\subsection{Hedonic Pricing: The Lancaster-Rosen Model}

A property is a bundle of characteristics---location, size, quality,
neighbourhood attributes, school quality. The \textbf{hedonic regression}
\index{hedonic regression} estimates the implicit price of each characteristic:
\begin{equation}
  \ln P_i = X_i'\beta + \varepsilon_i
\end{equation}
where $P_i$ is the sale price of property $i$ and $X_i$ is a vector of
property and neighbourhood characteristics \citep{Lancaster1966, Rosen1974}.
The hedonic coefficient $\beta_k$ measures the percentage change in price
associated with a unit change in characteristic $k$---the \emph{implicit
price} of that characteristic.

\noindent\textbf{Identification challenge.} The implicit price of an amenity
(e.g.\ proximity to a park, school quality) may be confounded by unobserved
neighbourhood quality if high-amenity neighbourhoods systematically attract
higher-income residents. IV strategies---using boundary discontinuities
or policy changes as instruments---are required for causal identification
\citep{Black1999}.

\subsection{Repeat-Sales Indices}

The \textbf{repeat-sales method}\index{repeat-sales index} \citep{Bailey1963,
Case1989} estimates price change from properties that sold more than once,
controlling for unobserved time-invariant quality by taking first differences:
\begin{equation}
  \ln P_{it} - \ln P_{is} = \sum_{\tau=s+1}^{t} \gamma_\tau D_\tau + \varepsilon_{it}
\end{equation}
where $D_\tau$ are time dummies and $\gamma_\tau$ traces the aggregate price
index. This is the methodology underlying the Case-Shiller index and the
FHFA House Price Index.

% TODO: Add worked example with Zillow ZHVI data in R


\section{Labour Market Geography}
\label{sec_LabourGeo}

\subsection{Commuting Zones and Local Labour Markets}

\citet{Autor2013} delineate \textbf{commuting zones}\index{commuting zones}
as the natural unit of local labour market analysis---functional economic
areas that contain the bulk of within-labour-market commuting flows.
The LEHD/LODES data (Chapter~\ref{chap_Data}) provide the empirical basis
for constructing commuting zones at fine geographic resolution using
workplace-to-residence flow data.

\subsection{The Spatial Mismatch Hypothesis}

\citet{Kain1968} argued that residential segregation combined with suburban
employment decentralisation created a geographic mismatch between where
low-income minority workers lived and where jobs were located.
Testing this hypothesis requires both a measure of job accessibility and
a credible instrument for residential location---since households sort into
locations partly on the basis of employment access. The LEHD/LODES data,
combined with the routing tools in the \texttt{osmdata} and \texttt{dodgr}
R packages, provide the inputs for constructing contemporary spatial
mismatch measures.

% TODO: Add commuting zone delineation example with LODES data

\renewcommand\bibname{Further reading}
\begin{subappendices}
\bibliographystyle{econometrica}
\bibliography{Literature/OverLord}
\IfFileExists{Literature/SpatialEquil.bbl}{\input{Literature/SpatialEquil.bbl}}{\typeout{Need bibtex on SpatialEquil}}
\end{subappendices}

% =====================================================================
%  Chapter 8: Measuring the Spatial Economy
%  Source: EconAnalysis.tex (Overleaf) — complete content
% =====================================================================

\chapter{Measuring the Spatial Economy}
\label{chap_EconAnalysis}

\section{Modelling and Measuring Economic Activity}
\label{sec_EconActivity}

Quantitative measures of economic structure serve two purposes in regional
analysis. First, they describe the industrial composition of a local economy
relative to a reference geography (usually the national average), revealing
its specialisations and dependencies. Second, they form the empirical basis
for testing theories of agglomeration, cluster formation, and regional growth.

\subsection{Measures of Spatial Concentration and Specialisation}

\subsubsection{Locational Measures}

\paragraph{Location Quotient}\index{location quotient}

The \textbf{location quotient} (LQ) is the standard measure of employment
distribution that controls for the size of the region. The relative
concentration of industry $i$ in region $j$ is:

\begin{tcolorbox}[keyresult]
\begin{equation}
  LQ_i = \frac{E_{ij}/E_{in}}{E_j/E_n}
  \label{eq_LQ}
\end{equation}
\end{tcolorbox}

\noindent where $E_{ij}$ is employment in industry $i$ in region $j$,
$E_j$ is total employment in region $j$, $E_{in}$ is national employment
in industry $i$, and $E_n$ is total national employment.

An $LQ > 1$ indicates above-average concentration of industry $i$ in region
$j$; industries with $LQ > 1.25$ are generally viewed as constituting the
region's export-oriented \textbf{economic base}\index{economic base}.
An $LQ < 1$ indicates a sector that serves primarily local demand.

\paragraph{Horizontal Clustering}\index{clustering!horizontal}

\textbf{Horizontal clustering} (HC) is an alternative measure that accounts
for possible agglomeration effects in terms of the absolute number of jobs
in a specific industry. It is the arithmetic equivalent of a bubble chart
that accounts for an industry's gravitational pull over and above what is
expected at the national level:
\begin{equation}
  HC_i = E_{ij} - \hat{E}_{ij}
\end{equation}
where $\hat{E}_{ij} = E_j \times (E_{in}/E_n)$ is the expected employment
level obtained by setting $\widehat{LQ}_i = 1$.

\subsubsection{Shift-Share Analysis}\index{shift-share analysis}

Shift-share analysis decomposes regional employment growth into three
components:
\begin{equation}
  \Delta E_{ij} = \underbrace{E_{ij}\,\frac{\Delta E_n}{E_n}}_{\text{National growth}}
  + \underbrace{E_{ij}\!\left(\frac{\Delta E_{in}}{E_{in}} - \frac{\Delta E_n}{E_n}\right)}_{\text{Industry mix}}
  + \underbrace{E_{ij}\!\left(\frac{\Delta E_{ij}}{E_{ij}} - \frac{\Delta E_{in}}{E_{in}}\right)}_{\text{Competitive effect}}
\end{equation}

The \textbf{national growth} component attributes growth to the overall
national trend. The \textbf{industry mix} (or proportionality) component
asks whether the region specialises in fast- or slow-growing industries
nationally. The \textbf{competitive} (or differential) component captures
whether local industries outperform their national counterparts---the
residual attributable to location-specific factors.

% TODO: Add worked example with BLS QCEW data

\subsubsection{Agglomeration Measures}

Location measures generally ignore geographic dispersion and the uneven
distribution of employment within the subregions of an area under study.
\textbf{Agglomeration measures}\index{agglomeration} are designed to overcome
this shortcoming.

\paragraph{Herfindahl-Hirschman Index}\index{Herfindahl-Hirschman index}

The \textbf{HHI} measures employment concentration across subregions of an
MSA:
\begin{equation}
  HHI_i = \sum_{j=1}^{n} s_j^2
\end{equation}
where $s_j$ is industry $i$'s LQ at the level of subregion $j$ relative to
the MSA, and $n$ is the number of counties within the MSA.\footnote{For the
Blacksburg MSA, $n=4$ (Giles, Montgomery, Pulaski counties and the city of
Radford).} The index equals 1 under absolute concentration and $1/n$ under
uniform dispersion.

\paragraph{Locational Gini Coefficient}\index{Gini coefficient!locational}

The \textbf{locational Gini coefficient} (LGC) measures industrial
concentration within a region more rigorously than the HHI:
\begin{equation}
  LGC_i = \frac{\sum_{i=1}^{n}\sum_{j=1}^{n}|x_i - x_j|}{2n(n-1)\mu}
\end{equation}
where $x_i$ and $x_j$ are the LQs of industries $i$ and $j$ in each of the
subregions, $\mu$ is the mean LQ in the MSA, and $n$ is the number of
subregions.

\paragraph{Ellison-Glaeser Index}\index{Ellison-Glaeser index}

The \textbf{Ellison-Glaeser} (EG) index \citep{Ellison1997} measures how
industrial concentration patterns differ from a situation where firms locate
randomly:
\begin{equation}
  EG_i = \frac{\displaystyle\sum_{j=1}^{n}(s_j - x_j)^2
               - (1-\sum_{j=1}^{n}x_j^2)\sum_{k=1}^{m}z_k^2}
              {\displaystyle(1-\sum_{j=1}^{n}x_j^2)(1-\sum_{k=1}^{m}z_k^2)}
\end{equation}
where $s_j$ is the employment share of subregion $j$ in industry $i$;
$x_j$ is the overall employment share of subregion $j$; $z_k$ is the share
of 6-digit NAICS subsector establishments in sector $k$; and $m$ is the
number of 6-digit subsectors.

The EG index accounts for both geographic concentration and industrial
concentration (plant-level returns to scale), yielding a ``dartboard''
measure: if the observed concentration exceeds what would arise if firms
chose locations randomly, the excess is attributed to agglomeration
economies.


\section{Measuring Inequality}
\label{sec_Inequality}

% TODO: Lorenz curves, Gini coefficient, Theil index
% TODO: Spatial decomposition: between-region vs. within-region inequality

Inequality measurement in a spatial context requires decomposing overall
inequality into components attributable to differences \emph{across} regions
and differences \emph{within} regions. The Theil index is particularly useful
because it admits an exact additive decomposition:
\begin{equation}
  T = T_\text{between} + T_\text{within}
\end{equation}

% TODO: Worked example with ACS income data by census tract


\section{Spatial Concentration and Agglomeration}
\label{sec_SpatialConc}

% TODO: Connect EG index to NEG theory; add empirical evidence on
%       US industrial clusters (Porter 2003, Glaeser et al. 1992)


\section{Economic Impact Analysis}
\label{sec_EconImpact}

\subsection{Input-Output Modelling}
\label{subsec_IO}

The \textbf{Leontief input-output model}\index{input-output model} describes
the interdependencies among sectors of an economy. Let $A$ be the
\textbf{technical coefficients matrix}\index{technical coefficients matrix},
where $a_{ij}$ is the amount of good $i$ required to produce one unit of
good $j$. The model is:
\begin{equation}
  \mathbf{x} = A\mathbf{x} + \mathbf{f}
  \qquad\Longrightarrow\qquad
  (I - A)\mathbf{x} = \mathbf{f}
\end{equation}
where $\mathbf{x}$ is the vector of total output by sector and $\mathbf{f}$
is the vector of final demand. The solution:

\begin{tcolorbox}[keyresult]
\begin{equation}
  \mathbf{x} = (I - A)^{-1}\mathbf{f} = L\mathbf{f}
\end{equation}
where $L = (I-A)^{-1}$ is the \textbf{Leontief inverse}\index{Leontief
inverse} (total requirements matrix).
\end{tcolorbox}

Each element $l_{ij}$ of $L$ gives the total output of sector $i$ required,
directly and indirectly, to deliver one unit of sector $j$'s output to
final demand.

\paragraph{Regional multipliers.}

The \textbf{output multiplier} for sector $j$ is $\sum_i l_{ij}$ (column
sum of $L$): the total output generated across all sectors per unit of
final demand for sector $j$. BEA's RIMS II
multipliers\index{RIMS II multipliers} are derived from the national I-O
accounts supplemented with regional employment data.

\noindent\textbf{Open-source I-O analysis.} BEA's publicly available
detailed use tables, combined with the \texttt{leontief} package in R,
permit full regional multiplier estimation without proprietary software:

\begin{lstlisting}[language=R, caption={Open I-O analysis with BEA data}]
library(leontief)
# Read BEA Use Table (downloaded from bea.gov)
use_table <- read.csv("bea_use_table.csv", row.names = 1)
A <- technical_coefficients(use_table)  # compute A matrix
L <- leontief_inverse(A)                # compute L = (I-A)^{-1}
multipliers <- colSums(L)               # output multipliers
\end{lstlisting}

\subsection{Applied Tools}
\label{subsec_AppliedIO}

\subsubsection{RIMS II}\index{RIMS II}

The Bureau of Economic Analysis (BEA) Regional Input-Output Modelling System
(RIMS II) provides multipliers for any area (state, MSA, county) for any
industry defined in the BEA national I-O accounts. Multipliers are available
for output, earnings, and employment. They are derived from two primary data
sources: the national I-O accounts and the BEA Regional Economic Accounts.

\subsubsection{IMPLAN}\index{IMPLAN}

IMPLAN is a proprietary social accounting matrix (SAM) model that extends
the standard I-O framework by adding household income and government accounts.
It is widely used in planning practice for economic impact assessments.
The key advantage over RIMS II is its spatial flexibility (any custom
geographic boundary) and its ability to model household spending patterns.
The key disadvantage is cost and the ``black box'' nature of proprietary
multipliers.

\subsubsection{FedFIT}\index{FedFIT}

FedFIT (Federal Footprint in the Territories) is a BEA tool for estimating
the economic contribution of federal government spending in US territories
and states.

\paragraph{A note on interpreting multipliers.}

Common mistakes in economic impact analysis include: (i)
\textbf{double-counting}---each sector's contribution should be computed once;
(ii) \textbf{ignoring displacement}---new activity may displace existing
activity rather than generating net new output; (iii) \textbf{treating I-O
estimates as causal}---I-O multipliers describe input requirements under a
fixed technology, not causal effects of policy interventions. For causal
analysis of economic development interventions, quasi-experimental
identification strategies (Chapter~\ref{chap_BeyondOLS}) are required.

\renewcommand\bibname{Further reading}
\begin{subappendices}
\bibliographystyle{econometrica}
\bibliography{Literature/OverLord}
\IfFileExists{Literature/EconAnalysis.bbl}{\chapter{Economic Analysis}\label{chap_EconAnalysis}

\section{Modelling and measuring economic activity}
\subsection{Measures of spatial concentration and specialisation}
The various quantitative measures of economic concentration\index{measures of economic concentration} are defined in detail in this section.

\subsubsection{Locational measures}

\paragraph{Location quotient}
The location quotient\index{location quotient|see{measures of economic concentration}}\index{measures of economic concentration!location quotient} (LQ) is the standard measure of employment distribution that controls for the size of the region. The relative
concentration of industry $i$ in region $j$ is defined as

\begin{eqnarray}
LQ_{i} = \frac{(E_{ij}/E_{in})}{(E_{j}/E_{n})}
\end{eqnarray}

\noindent where, $E_{ij}$ is employment in industry $i$ in region $j$, $E_{j}$ is total employment in region $j$, $E_{in}$ is national employment in industry $i$, and $E_{n}$ is total national employment. Thus, a LQ of greater than one indicates that there is an above average proportion of employment in a given industry in a given region. Industries with an LQ above 1.25 are generally viewed
as constituting the export-oriented economic base.

\paragraph{Horizontal clustering}
Horizontal clustering\index{clustering!horizontal|see{measures of economic concentration}}\index{measures of economic concentration!horizontal clustering} (HC) is an alternative measure that accounts for possible agglomeration effects in terms of the numbers of jobs in a specific industry. As such, it is the arithmetic equivalent to a bubble chart\index{bubble chart} which accounts for an industry's gravitational pull over and above what is expected to occur at the national level. HC can then be written as

\begin{eqnarray}
HC_{i} = E_{ij} - \widehat{E}_{ij},
\end{eqnarray}

\noindent where $\widehat{E}_{ij} = E_{j}\times \frac{E_{in}}{E_{n}}$ by setting $\widehat{LQ}_{i}=1$.

\subsubsection{Shift-share analysis}


\subsubsection{Agglomeration measures}
Location measures generally ignore geographic dispersion and uneven distribution of employment within subregions of an area under study.
Agglomeration measures\index{agglomeration}\index{agglomeration|seealso{measures of economic concentration}} are designed to overcome this shortcoming.

\paragraph{Herfindahl-Hirschman Index}
In the context of industry employment concentration within an MSA the simplest such measure is the Herfindahl-Hirschman index\index{Herfindahl-Hirschman index|see{measures of economic concentration}}\index{measures of economic concentration!Herfindahl-Hirschman index} (HHI). It is defined as

\begin{eqnarray}
HHI_{i} = \sum^{n}_{i=1}s_{i}^{2},
\end{eqnarray}

\noindent where $s_{i}$ is simply industry $i$'s LQ at the level of the subregion compared to the MSA and $n$ is the number of counties
within the MSA.\footnote{In the context of MSA-level data, counties are commonly assumed to be the subregions. For the Blacksburg MSA,
we would set $n=4$ since it includes Giles, Montgomery, Pulaski counties and the city of Radford.} The index is equal to 1 if there is absolute concentration and it takes a value of $\frac{1}{n}$ if employment in the industry is equally dispersed across the MSA.

\paragraph{Locational Gini Coefficient}
The locational Gini coefficient\index{Gini coefficient!locational} (LGC) also accounts for agglomeration and concentration within a specific region, but in a most sophisticated way than the HHI. It is defined as

\begin{eqnarray}
LGC_{i} =
\frac{\sum_{i=1}^{n}\sum_{j=1}^{n}|x_{i}-x_{j}|}{2n(n-1)\mu},
\end{eqnarray}

\noindent where $x_{i}$ and $x_{j}$ are the LQs of industries $i$ and $j$ in each of the subregions, $\mu$ is the mean of the LQs in the MSA and $n$ is the number of counties.

\paragraph{Ellison-Glaeser Index}
The Ellison-Glaeser index\index{Ellison-Glaeser index|see{measures of economic concentration}}\index{measures of economic concentration!Ellison-Glaeser index} (EG) measures how industrial concentration patterns differ from a situation where firms in a purely random manner. The index is defined as follows

\begin{equation}
EG_{i} = \frac{\sum_{i=1}^{n}(s_{i}-x_{i})^2 - (1-
\sum^{n}_{i=1}x^{2}_{i})\sum^{m}_{j=1}z^{2}_{j}}{(1-
\sum^{n}_{i=1}x^{2}_{i})(1- \sum^{m}_{j=1}z^{2}_{j})},
\end{equation}

\noindent where $s_{i}$, $x_{i}$ and $n$ have the same definitions as above; $z_{i}$ is the share of the 6-digit NAICS subsector
establishments in sector $j$ and $m$ is the number of 6-digit subsectors.


\section{Measuring inequality}

\section{Spatial concentration and agglomeration}

\section{Economic impact analysis}
\subsection{Input-output modelling}
\subsection{Applied tools}
\subsubsection{RIMS II}
\subsubsection{IMPLAN}
\subsubsection{FedFIT}

\renewcommand\bibname{Further reading}
\begin{subappendices}
\bibliographystyle{econometrica}
\bibliography{Literature/OverLord}
\IfFileExists{Literature/EconAnalysis.bbl}{\chapter{Economic Analysis}\label{chap_EconAnalysis}

\section{Modelling and measuring economic activity}
\subsection{Measures of spatial concentration and specialisation}
The various quantitative measures of economic concentration\index{measures of economic concentration} are defined in detail in this section.

\subsubsection{Locational measures}

\paragraph{Location quotient}
The location quotient\index{location quotient|see{measures of economic concentration}}\index{measures of economic concentration!location quotient} (LQ) is the standard measure of employment distribution that controls for the size of the region. The relative
concentration of industry $i$ in region $j$ is defined as

\begin{eqnarray}
LQ_{i} = \frac{(E_{ij}/E_{in})}{(E_{j}/E_{n})}
\end{eqnarray}

\noindent where, $E_{ij}$ is employment in industry $i$ in region $j$, $E_{j}$ is total employment in region $j$, $E_{in}$ is national employment in industry $i$, and $E_{n}$ is total national employment. Thus, a LQ of greater than one indicates that there is an above average proportion of employment in a given industry in a given region. Industries with an LQ above 1.25 are generally viewed
as constituting the export-oriented economic base.

\paragraph{Horizontal clustering}
Horizontal clustering\index{clustering!horizontal|see{measures of economic concentration}}\index{measures of economic concentration!horizontal clustering} (HC) is an alternative measure that accounts for possible agglomeration effects in terms of the numbers of jobs in a specific industry. As such, it is the arithmetic equivalent to a bubble chart\index{bubble chart} which accounts for an industry's gravitational pull over and above what is expected to occur at the national level. HC can then be written as

\begin{eqnarray}
HC_{i} = E_{ij} - \widehat{E}_{ij},
\end{eqnarray}

\noindent where $\widehat{E}_{ij} = E_{j}\times \frac{E_{in}}{E_{n}}$ by setting $\widehat{LQ}_{i}=1$.

\subsubsection{Shift-share analysis}


\subsubsection{Agglomeration measures}
Location measures generally ignore geographic dispersion and uneven distribution of employment within subregions of an area under study.
Agglomeration measures\index{agglomeration}\index{agglomeration|seealso{measures of economic concentration}} are designed to overcome this shortcoming.

\paragraph{Herfindahl-Hirschman Index}
In the context of industry employment concentration within an MSA the simplest such measure is the Herfindahl-Hirschman index\index{Herfindahl-Hirschman index|see{measures of economic concentration}}\index{measures of economic concentration!Herfindahl-Hirschman index} (HHI). It is defined as

\begin{eqnarray}
HHI_{i} = \sum^{n}_{i=1}s_{i}^{2},
\end{eqnarray}

\noindent where $s_{i}$ is simply industry $i$'s LQ at the level of the subregion compared to the MSA and $n$ is the number of counties
within the MSA.\footnote{In the context of MSA-level data, counties are commonly assumed to be the subregions. For the Blacksburg MSA,
we would set $n=4$ since it includes Giles, Montgomery, Pulaski counties and the city of Radford.} The index is equal to 1 if there is absolute concentration and it takes a value of $\frac{1}{n}$ if employment in the industry is equally dispersed across the MSA.

\paragraph{Locational Gini Coefficient}
The locational Gini coefficient\index{Gini coefficient!locational} (LGC) also accounts for agglomeration and concentration within a specific region, but in a most sophisticated way than the HHI. It is defined as

\begin{eqnarray}
LGC_{i} =
\frac{\sum_{i=1}^{n}\sum_{j=1}^{n}|x_{i}-x_{j}|}{2n(n-1)\mu},
\end{eqnarray}

\noindent where $x_{i}$ and $x_{j}$ are the LQs of industries $i$ and $j$ in each of the subregions, $\mu$ is the mean of the LQs in the MSA and $n$ is the number of counties.

\paragraph{Ellison-Glaeser Index}
The Ellison-Glaeser index\index{Ellison-Glaeser index|see{measures of economic concentration}}\index{measures of economic concentration!Ellison-Glaeser index} (EG) measures how industrial concentration patterns differ from a situation where firms in a purely random manner. The index is defined as follows

\begin{equation}
EG_{i} = \frac{\sum_{i=1}^{n}(s_{i}-x_{i})^2 - (1-
\sum^{n}_{i=1}x^{2}_{i})\sum^{m}_{j=1}z^{2}_{j}}{(1-
\sum^{n}_{i=1}x^{2}_{i})(1- \sum^{m}_{j=1}z^{2}_{j})},
\end{equation}

\noindent where $s_{i}$, $x_{i}$ and $n$ have the same definitions as above; $z_{i}$ is the share of the 6-digit NAICS subsector
establishments in sector $j$ and $m$ is the number of 6-digit subsectors.


\section{Measuring inequality}

\section{Spatial concentration and agglomeration}

\section{Economic impact analysis}
\subsection{Input-output modelling}
\subsection{Applied tools}
\subsubsection{RIMS II}
\subsubsection{IMPLAN}
\subsubsection{FedFIT}

\renewcommand\bibname{Further reading}
\begin{subappendices}
\bibliographystyle{econometrica}
\bibliography{Literature/OverLord}
\IfFileExists{Literature/EconAnalysis.bbl}{\input{Literature/EconAnalysis.bbl}}{\typeout{Need to run bibtex}}
\end{subappendices} }{\typeout{Need to run bibtex}}
\end{subappendices} }{\typeout{Need bibtex on EconAnalysis}}
\end{subappendices}

% =====================================================================
%  Chapter 9: Spatial Dependence and Autocorrelation
%  Source: SpatialStats.tex (motivations prose) + P4 transcription (full)
% =====================================================================

\chapter{Spatial Dependence and Autocorrelation}
\label{chap_Spatial}

\section{Motivations for Going Spatial}
\label{sec_MotivationsSpatial}

Spatial autocorrelation\index{spatial autocorrelation} can be loosely defined
as ``the coincidence of value similarity with locational similarity''
\citep{Anselin2002}. The problem of spatial dependence is best viewed as a
normal extension of Tobler's First Law of Geography:

\begin{quote}
``everything is related to everything else, but near things are more related
than distant things'' \citep[p.236]{Tobler1970}.
\end{quote}

Three main reasons motivate the use of spatial methods:

\begin{enumerate}[label=(\roman*)]
  \item \textbf{The independence assumption is not valid:} Attributes of
    observation $i$ may influence attributes of observation $j$.
  \item \textbf{Spatial heterogeneity:} The magnitude and direction of the
    treatment effect may vary across space.
  \item \textbf{Omitted variable bias:} Some common unobserved or latent
    variables are shared among spatial neighbours.
\end{enumerate}

Recall the non-spatial DGP:
\begin{equation}
  y_i = X_i\beta + e_i, \qquad e_i \sim \text{iid}(0,\sigma^2)
  \label{eq_NonSpatDGP}
\end{equation}

With assumptions: (1) observed values at location $i$ are independent of
values at location $j$; (2) $\E[e_ie_j] = \E[e_i]\E[e_j] = 0$.

\begin{tcolorbox}[definition]
\textbf{MAUP Warning.} The choice of spatial weight matrix is not merely
technical---it is an inferential decision with the same logical status as the
choice of instruments in IV estimation. Results should always be checked for
sensitivity to weight matrix specification.
\end{tcolorbox}


\section{The Spatial DGP}
\label{sec_SpatialDGP}

With 2 neighbours $i$ and $j$, observations depend on each other:
\begin{alignat*}{2}
  y_i &= \alpha_j y_j + X_i\beta + e_i, \qquad &e_i &\sim \text{iid}(0,\sigma^2) \\
  y_j &= \alpha_i y_i + X_j\beta + e_j, \qquad &e_j &\sim \text{iid}(0,\sigma^2)
\end{alignat*}

\noindent\textbf{Problem.} With more spatial dependencies, there is little
practical usefulness since the system has \emph{more parameters than
observations}: potentially $n^2 - n$ relations to be estimated.

\noindent\textbf{Assumptions of the spatial DGP:}
\begin{itemize}
  \item Observations at location $i$ depend on those at location $j$ and
    vice versa.
  \item The DGP is \emph{simultaneous} across space.
\end{itemize}

The overparametrisation problem is solved by imposing structure in a
\textbf{spatial autoregressive process}\index{spatial autoregressive process}.
Introducing a scalar spatial correlation coefficient $\rho$:

\begin{tcolorbox}[keyresult]
\begin{equation}
  y_i = \rho\sum_{j=1}^{n} W_{ij}\,y_j + e_i
\end{equation}
where $\sum_j W_{ij}y_j$ is the \textbf{spatial lag}\index{spatial lag} of
$y_i$---the weighted average of neighbouring observations.
\end{tcolorbox}


\section{The Spatial Weight Matrix}
\label{sec_WeightMatrix}

In matrix notation ($W$ of size $n\times n$; $y$ of size $n\times 1$):
\begin{equation}
  y = \rho Wy + e, \qquad e \sim N(0, \sigma^2 I_n)
\end{equation}

To construct $W$, begin with a \textbf{contiguity matrix}
$C$\index{contiguity matrix}, where rows represent observation $i$ and
columns represent neighbour $j$ (1 = neighbour, 0 = not).

\subsection{Example: Four Contiguous Regions}

Consider four regions $R_1$--$R_4$ arranged in a chain. First-order
neighbours ($N_1$) share a boundary; $R_1$ is its own second-order
neighbour (the neighbour of its first-order neighbour $R_2$).

\begin{figure}[htp]
\centering
\begin{tikzpicture}[>=Stealth, node distance=0cm]
  \node[draw, minimum width=1.2cm, minimum height=0.8cm] (R1) {$R_1$};
  \node[draw, minimum width=1.2cm, minimum height=0.8cm, right=of R1] (R2) {$R_2$};
  \node[draw, minimum width=1.2cm, minimum height=0.8cm, right=of R2] (R3) {$R_3$};
  \node[draw, minimum width=1.2cm, minimum height=0.8cm, right=of R3] (R4) {$R_4$};
  \draw[<->, thick] ([yshift=6pt]R1.north) to[bend left=40]
    node[above] {$N_1$} ([yshift=6pt]R2.north);
  \draw[<->, thick] ([yshift=-6pt]R2.south) to[bend right=40]
    node[below] {$N_2$} ([yshift=-6pt]R3.south);
  \draw[<->, thick] ([yshift=-6pt]R3.south) to[bend right=40]
    node[below] {$N_2$} ([yshift=-6pt]R4.south);
\end{tikzpicture}
\caption{Four contiguous regions with first-order ($N_1$) and
higher-order ($N_2$) neighbours.}
\label{fig_ContiguousRegions}
\end{figure}

The contiguity matrix:
\begin{equation}
  C = \begin{bmatrix}
    0 & 1 & 0 & 0 \\
    1 & 0 & 1 & 0 \\
    0 & 1 & 0 & 1 \\
    0 & 0 & 1 & 0
  \end{bmatrix}
\end{equation}

Create $W$ by \textbf{row-standardisation} (row sums normalised to one):
\begin{equation}
  W = \begin{bmatrix}
    0   & 1   & 0   & 0   \\
    1/2 & 0   & 1/2 & 0   \\
    0   & 1/2 & 0   & 1/2 \\
    0   & 0   & 1   & 0
  \end{bmatrix}
\end{equation}

\noindent\textbf{Verifying the spatial lag $Wy$:}
\begin{equation}
  Wy = \begin{bmatrix}y_2\\(y_1+y_3)/2\\(y_2+y_4)/2\\y_3\end{bmatrix}
\end{equation}

Each element of $Wy$ is the (row-weighted) average of the neighbours of the
corresponding observation.

\subsection{Weight Matrix Types}

\begin{itemize}
  \item \textbf{Rook contiguity:} Shares a boundary (no diagonal touches)
  \item \textbf{Queen contiguity:} Shares a boundary or corner
  \item \textbf{$k$-nearest neighbours:} Fixed number of nearest units
  \item \textbf{Distance decay:} $w_{ij} = d_{ij}^{-\delta}$ for some
    $\delta > 0$, often with a distance cutoff
  \item \textbf{Row-standardised:} Rows sum to one; enables spatial lag
    interpretation as a local average
\end{itemize}


\section{The Spatial Autoregressive (SAR) DGP}
\label{sec_SARDGP}

$\rho$ is a scalar parameter of spatial dependence:
\begin{equation}
  \rho \;\begin{cases}
    = 0 & \text{no spatial autocorrelation} \\
    > 0 & \text{positive spatial autocorrelation (clustering)} \\
    < 0 & \text{negative spatial autocorrelation (segregation)}
  \end{cases}
\end{equation}

The full \textbf{SA DGP} with constant term and covariates:
\begin{align}
  y &= \rho Wy + \boldsymbol{\iota}\alpha + e \\
  (I_n - \rho W)y &= \boldsymbol{\iota}\alpha + e \\
  y &= (I_n - \rho W)^{-1}\boldsymbol{\iota}\alpha + (I_n - \rho W)^{-1}e
\end{align}

Assuming $|\rho| < 1$ (stationarity condition; see Section~\ref{sec_Eigenvalues}
for the general condition $|\rho| < 1/\lambda_{\max}(W)$), the inverse can be
expressed as a \textbf{geometric matrix series}:
\begin{equation}
  (I_n - \rho W)^{-1} = I_n + \rho W + \rho^2 W^2 + \rho^3 W^3 + \cdots
\end{equation}

Therefore:
\begin{equation}
  y = \boldsymbol{\iota}\alpha + \rho W\boldsymbol{\iota}\alpha
      + \rho^2 W^2\boldsymbol{\iota}\alpha + \cdots
      + e + \rho We + \rho^2 W^2 e + \cdots
\end{equation}

The first part converges to $(1-\rho)^{-1}\boldsymbol{\iota}\alpha$ (since
$\alpha$ is scalar and $W\boldsymbol{\iota} = \boldsymbol{\iota}$):

\begin{tcolorbox}[keyresult]
\begin{equation}
  y = \frac{1}{1-\rho}\,\boldsymbol{\iota}\alpha
      + \underbrace{e + \rho We + \rho^2 W^2 e + \cdots}_{\text{violation of Gauss-Markov}}
\end{equation}
\end{tcolorbox}

This is a \textbf{strong violation} of the Gauss-Markov assumption of
independent error terms. Therefore \textbf{OLS is biased and inconsistent}
under a SAR DGP.

\subsection{SAR Model with Covariates}

\begin{equation}
  y = \rho Wy + X\beta + e
  \;\Longrightarrow\;
  y = (I_n-\rho W)^{-1}X\beta + (I-\rho W)^{-1}e, \quad e\sim N(0,\sigma^2I)
\end{equation}

Estimate three parameters: $\beta$, $\sigma^2$, and $\rho$.
If $\rho=0$, the SAR model reduces to standard OLS.

The SAR estimator can be written as:

\begin{tcolorbox}[keyresult]
\begin{equation}
  \hat\beta_\text{SAR} = \hat\beta_\text{OLS} - \hat\rho(X'X)^{-1}X'Wy
  \label{eq_SARbeta}
\end{equation}
\end{tcolorbox}


\section{Global Spatial Autocorrelation: Moran's $I$}
\label{sec_MoranI}

Moran's $I$ coefficient calculates the ratio of the product of a variable
and its spatial lag, adjusted for the sum of spatial weights, relative to
the overall variance:

\begin{tcolorbox}[keyresult]
\begin{equation}
  I = \frac{\displaystyle
       \frac{\displaystyle\sum_i\sum_j w_{ij}(y_i-\bar{y})(y_j-\bar{y})}
            {\displaystyle\sum_i\sum_j w_{ij}}}
      {\displaystyle\frac{\sum_i(y_i-\bar{y})^2}{n}}
  \;\in\; [-1,\,+1]
  \label{eq_MoranI}
\end{equation}
\end{tcolorbox}

$I = +1$: perfect spatial clustering; $I = -1$: perfect dispersion;
$I \approx 0$: spatial randomness. Moran's $I$ can equivalently be
interpreted as the \textbf{regression coefficient of a variable on its
spatial lag}---the slope of the Moran scatterplot.

\subsection{Local Moran's $I$}

The global statistic can be decomposed into contributions from individual
observations:
\begin{equation}
  I_i = \frac{(y_i-\bar{y})\displaystyle\sum_j w_{ij}(y_j-\bar{y})}
             {\displaystyle\frac{\sum_i(y_i-\bar{y})^2}{n}}
  \label{eq_LocalMoran}
\end{equation}

Local Moran's $I$ \citep{Anselin1995} is the basis of LISA (Local Indicators
of Spatial Association) analysis. A map of local $I_i$ $z$-scores identifies:
\begin{itemize}
  \item \textbf{HH} (high-high): high-value location surrounded by
    high-value neighbours---core of a spatial cluster
  \item \textbf{LL} (low-low): low-value surrounded by low-value neighbours
  \item \textbf{HL/LH} (high-low or low-high): spatial outliers---locations
    significantly different from their neighbourhood
\end{itemize}


\section{Estimating Models with a SAR DGP}
\label{sec_EstimatingSAR}

Recall: OLS is inconsistent and biased under a SAR DGP. Use alternative
estimators: \textbf{MLE} or \textbf{IV/2SLS}. 2SLS has similar properties
to MLE (consistency and asymptotic normality) but is considerably easier
to implement numerically.

\subsection{SAR-IV Estimator}

Address endogeneity due to spatial variable linkage by using lagged variables
as instruments. Let $X = [Wy\;\; X_1]$ and $Z = [L(y)\;\; X_1]$, where
$L(y)$ denotes higher-order spatial lags. Then:
\begin{equation}
  \hat\beta_\text{SAR-IV} = [X'P_ZX]^{-1}X'P_Zy,
  \qquad P_Z \equiv Z[Z'Z]^{-1}Z'
\end{equation}


\section{The Spatial Model Taxonomy}
\label{sec_SpatialTaxonomy}

All models below are expressed relative to OLS ($y = X\beta + e$):

\renewcommand{\arraystretch}{1.5}
\begin{table}[htp]
\centering
\caption{Spatial regression model taxonomy}
\label{tbl_SpatialModels}
\begin{tabular}{@{}lp{5.5cm}p{4.5cm}@{}}
\toprule
\textbf{Model} & \textbf{Equation} & \textbf{Notes} \\
\midrule
OLS  & $y = X\beta + e$ & Baseline; assumes spatial independence \\
SAR  & $y = \rho Wy + X\beta + e$ & Spatial lag in $y$; OLS biased \\
SEM  & $y = X\beta + (I-\lambda W)^{-1}e$ & Spatial error model; OLS unbiased but inefficient \\
SDM  & $y = \rho Wy + X\beta + WX\theta + e$ & Spatial Durbin model \\
SARMA & $y = \rho Wy + X\beta + (I-\theta W)^{-1}e$ & SAR + MA error \\
SAC  & $y = \rho W_1y + X\beta + (I-\lambda W_2)^{-1}e$ & Spatial dependence in $y$ and $e$; separate weight matrices \\
SDEM & $y = X\beta + WX\theta + (I-\lambda W)^{-1}e$ & Spatial Durbin error model \\
\bottomrule
\end{tabular}
\end{table}
\renewcommand{\arraystretch}{1}


\section{Geographically Weighted Regression}
\label{sec_GWR}

Standard OLS assumes $\beta$ is globally constant. GWR relaxes this by
allowing locally-varying parameter estimates:
\begin{equation}
  y_i = X\beta_i + e \qquad\text{with locally-varying }\beta_i
\end{equation}

\begin{tcolorbox}[keyresult]
\begin{equation}
  \hat\beta_\text{GWR}(i) = (X'W_iX)^{-1}X'W_iy
  \label{eq_GWR}
\end{equation}
\end{tcolorbox}

Weights $w_{ij}$ are calculated using a \textbf{Gaussian decay
kernel}\index{Gaussian kernel!GWR}:
\begin{equation}
  w_{ij} = e^{(-d_{ij}/h)^2}
\end{equation}
where $d_{ij}$ is the Euclidean distance between observations $i$ and $j$,
and $h$ is the user-selected \textbf{bandwidth}\index{bandwidth!GWR}
(determined by cross-validation).

\noindent\textbf{Interpreting spatial variation in coefficients.}
GWR maps of $\hat\beta_i$ show where the relationship between $X$ and $y$
is stronger or weaker. However, spatial variation in coefficients can reflect
either genuine spatial heterogeneity in the relationship or noise arising from
sparse data. Significance tests for GWR coefficients should be interpreted
cautiously, as they rely on asymptotic approximations that may not be valid
for small local sample sizes.


\section{A Brief Excursion on Map Projections}
\label{sec_Projections_chap9}

\subsection{Coordinate Systems}

\noindent\textbf{Lat:} North--south position $\equiv$ $y$-axis\hfill
\noindent\textbf{Lon:} East--west position $\equiv$ $x$-axis

Projection is necessary to represent the 3D globe on a 2D map; projections
always distort some property of the original globe.

\begin{figure}[htp]
\centering
\begin{tikzpicture}[>=Stealth, font=\small]
  % Cylindrical
  \begin{scope}[xshift=0cm]
    \draw[thick] (-0.7,-1.1) -- (-0.7,1.1);
    \draw[thick] ( 0.7,-1.1) -- ( 0.7,1.1);
    \draw[thick] (-0.7,1.1) arc[start angle=180,end angle=0,x radius=0.7,y radius=0.18];
    \draw[thick,dashed] (-0.7,1.1) arc[start angle=180,end angle=360,x radius=0.7,y radius=0.18];
    \draw[thick] (-0.7,-1.1) arc[start angle=180,end angle=360,x radius=0.7,y radius=0.18];
    \draw[thick,dashed] (-0.7,-1.1) arc[start angle=180,end angle=0,x radius=0.7,y radius=0.18];
    \draw[thick] (0,0) circle (0.55);
    \node[below] at (0,-1.3) {(i) Cylindrical};
  \end{scope}
  % Conic
  \begin{scope}[xshift=3.5cm]
    \draw[thick] (0,1.5) -- (-0.85,-0.3);
    \draw[thick] (0,1.5) -- ( 0.85,-0.3);
    \draw[thick] (-0.85,-0.3) arc[start angle=180,end angle=360,x radius=0.85,y radius=0.2];
    \draw[thick,dashed] (-0.85,-0.3) arc[start angle=180,end angle=0,x radius=0.85,y radius=0.2];
    \draw[thick] (0,-0.1) circle (0.55);
    \node[below] at (0,-1.0) {(ii) Conic};
  \end{scope}
  % Planar
  \begin{scope}[xshift=7.0cm]
    \fill[gray!15] (-0.9,0.2) -- (0,1.2) -- (0.9,0.2) -- (0,-0.8) -- cycle;
    \draw[thick]   (-0.9,0.2) -- (0,1.2) -- (0.9,0.2) -- (0,-0.8) -- cycle;
    \draw[thick] (0,0) circle (0.5);
    \node[below] at (0,-1.0) {(iii) Planar};
  \end{scope}
\end{tikzpicture}
\caption{Three main types of map projection.}
\label{fig_Projections}
\end{figure}

\subsection{Projections and What They Preserve}

\renewcommand{\arraystretch}{1.4}
\begin{table}[htp]
\centering
\caption{Projection names by preserved property}
\label{tbl_Projections}
\begin{tabular}{@{}llll@{}}
\toprule
\textbf{Preserves} & \textbf{Cylindrical} & \textbf{Conic} & \textbf{Planar} \\
\midrule
Area     & \texttt{cylegnal\_area} & \texttt{albers}  & \texttt{azequal\_area} \\
Angle    & \texttt{mercator} (UTM, omerc) & \texttt{lambert} & \texttt{stereographic} \\
Distance & ---                    & ---              & \texttt{azequidistant}  \\
\bottomrule
\end{tabular}
\end{table}
\renewcommand{\arraystretch}{1}

\subsection{The Tissot Indicatrix}

The \textbf{Tissot indicatrix}\index{Tissot indicatrix} characterises
distortions due to map projections. Infinitesimally small circles of equal
radius placed at regular intervals on the sphere become ellipses after
projection; their size, shape, and orientation reveal how the projection
distorts area, angle, and distance.

\renewcommand\bibname{Further reading}
\begin{subappendices}
\bibliographystyle{econometrica}
\bibliography{Literature/OverLord}
\IfFileExists{Literature/SpatialStats.bbl}{\chapter{Applied Spatial Data Analysis}\label{chap_SpatialStats}
Quantitative analysis of spatial data has advanced greatly in recent years, largely as the result of the rapidly decreasing cost of computational power, the increased availability of high-resolution data and methodological innovation in spatial econometrics. Indeed, recent analytical and technological advances are increasingly blurring the traditional boundaries in spatial data analysis where various applied methods in Geographic Information Systems (GIS)\index{GIS|see{geographic information systems}}\index{geographic information systems}, geology-oriented remote sensing\index{remote sensing} and spatial statistics\index{spatial statistics} were essentially treated as separate fields. The legacy of this separate treatment of spatial data is reflected by the fact that even the leading textbooks in this area do not focus on more than two of these three sub-fields \citep[e.g.][]{Davis2002a,O'Sullivan2003,Campbell2006}.

Nowhere in the social sciences, perhaps, is this silent revolution felt more than in urban and regional planning. While planners usually have emphasised the application to GIS for land-use and environmental planning purposes, sometimes in addition with a little attention being paid to either remote sensing and spatial statistics, the discipline has never embraced all three sub-fields of quantitative data analysis in a systematic way. 

However, with the rapid deployment of ``Big Data''\index{Big Data} and the emergence of ``policy informatics''\index{policy informatics} as a new paradigm among computational social scientists, planners should increasingly see themselves at the forefront of the application of such integrated spatial data methods.\footnote{See a recent study by McKinsey \& Co. on this \href{http://www.mckinsey.com/Insights/MGI/Research/Technology\_and\_Innovation/Big\_data\_The\_next\_frontier\_for\_innovation}{topic}. See \citet{Barrett2011} for a definition of the term and on recent developments in the area of policy informatics.} \citet{Anselin2012} provides a good overview that assesses the evolution of the way in which spatial data analytical methods have been incorporated into software tools over the past two decades.


\section{Motivations for ``going spatial''}

Spatial autocorrelation\index{spatial autocorrelation|see{autocorrelation}}\index{autocorrelation!spatial} can be loosely defined as ``the coincidence of value similarity with locational similarity'' \citep{Anselin2002}. The problem of spatial dependence then is best viewed as a normal extension of the first law in geography, whereby

\begin{quote}
``everything is related to everything else, but near things are more related than distant things'' \citep[p.236]{Tobler1970}.
\end{quote}

Broadly speaking, there are three main reasons for deploying spatial methods:

\begin{enumerate}[(i)]
\item \emph{The independence assumption is not valid}: The attributes of observation $i$ may influence the attributes of observation
$j$.
\item \emph{Spatial heterogeneity}: The magnitude and direction the the treatment effect may vary across space.
\item \emph{Omitted variable bias}: Some common unobserved or latent variables are shared among spatial neighbours.
\end{enumerate}


Recall that the matrix notation of the non-spatial DGP was defined as

\begin{eqnarray}
\boldsymbol{y} = \boldsymbol{X\beta} + \boldsymbol{e} \qquad \text{with} \qquad \boldsymbol{e}\sim\mathcal{N}(0,\sigma^{2}\boldsymbol{I}).\label{eqn_SS_DGP}
\end{eqnarray}

Similarly, also recall that the classical normal regression model (CNR) of equation~\ref{eqn_SS_DGP} is based on the following assumptions:

\begin{itemize}
\item Observed values at location $i$ are independent of values at location $j$.
\item The regression residuals associated observation $i$ are independent of those associated with observation $j$, i.e. $\mathbb{E}[e_{i}e_{j}]=\mathbb{E}[e_{i}]\mathbb{E}[e_{j}]=0$.
\end{itemize}

\begin{itemize}
\item Include the following literature: \citet{Bivand2008}, \citet{Bivand2008a}, \citet{Corrado2012}, \citet{Kelejian1998}, \citet{McMillen2010a}, \citet{Pinske2010}.
\end{itemize}




\section{Spatial point pattern analysis}
\section{Interpolation and geostatistics}
\section{Areal data and spatial autocorrelation}

\begin{subappendices}
\section{Code for Spatial Stats in R}
\lstinputlisting{"Programs/UP504_SpatStat1.R"}

\renewcommand\bibname{Further reading}

\bibliographystyle{econometrica}
\bibliography{Literature/OverLord}
\IfFileExists{Literature/SpatialStats.bbl}{\chapter{Applied Spatial Data Analysis}\label{chap_SpatialStats}
Quantitative analysis of spatial data has advanced greatly in recent years, largely as the result of the rapidly decreasing cost of computational power, the increased availability of high-resolution data and methodological innovation in spatial econometrics. Indeed, recent analytical and technological advances are increasingly blurring the traditional boundaries in spatial data analysis where various applied methods in Geographic Information Systems (GIS)\index{GIS|see{geographic information systems}}\index{geographic information systems}, geology-oriented remote sensing\index{remote sensing} and spatial statistics\index{spatial statistics} were essentially treated as separate fields. The legacy of this separate treatment of spatial data is reflected by the fact that even the leading textbooks in this area do not focus on more than two of these three sub-fields \citep[e.g.][]{Davis2002a,O'Sullivan2003,Campbell2006}.

Nowhere in the social sciences, perhaps, is this silent revolution felt more than in urban and regional planning. While planners usually have emphasised the application to GIS for land-use and environmental planning purposes, sometimes in addition with a little attention being paid to either remote sensing and spatial statistics, the discipline has never embraced all three sub-fields of quantitative data analysis in a systematic way. 

However, with the rapid deployment of ``Big Data''\index{Big Data} and the emergence of ``policy informatics''\index{policy informatics} as a new paradigm among computational social scientists, planners should increasingly see themselves at the forefront of the application of such integrated spatial data methods.\footnote{See a recent study by McKinsey \& Co. on this \href{http://www.mckinsey.com/Insights/MGI/Research/Technology\_and\_Innovation/Big\_data\_The\_next\_frontier\_for\_innovation}{topic}. See \citet{Barrett2011} for a definition of the term and on recent developments in the area of policy informatics.} \citet{Anselin2012} provides a good overview that assesses the evolution of the way in which spatial data analytical methods have been incorporated into software tools over the past two decades.


\section{Motivations for ``going spatial''}

Spatial autocorrelation\index{spatial autocorrelation|see{autocorrelation}}\index{autocorrelation!spatial} can be loosely defined as ``the coincidence of value similarity with locational similarity'' \citep{Anselin2002}. The problem of spatial dependence then is best viewed as a normal extension of the first law in geography, whereby

\begin{quote}
``everything is related to everything else, but near things are more related than distant things'' \citep[p.236]{Tobler1970}.
\end{quote}

Broadly speaking, there are three main reasons for deploying spatial methods:

\begin{enumerate}[(i)]
\item \emph{The independence assumption is not valid}: The attributes of observation $i$ may influence the attributes of observation
$j$.
\item \emph{Spatial heterogeneity}: The magnitude and direction the the treatment effect may vary across space.
\item \emph{Omitted variable bias}: Some common unobserved or latent variables are shared among spatial neighbours.
\end{enumerate}


Recall that the matrix notation of the non-spatial DGP was defined as

\begin{eqnarray}
\boldsymbol{y} = \boldsymbol{X\beta} + \boldsymbol{e} \qquad \text{with} \qquad \boldsymbol{e}\sim\mathcal{N}(0,\sigma^{2}\boldsymbol{I}).\label{eqn_SS_DGP}
\end{eqnarray}

Similarly, also recall that the classical normal regression model (CNR) of equation~\ref{eqn_SS_DGP} is based on the following assumptions:

\begin{itemize}
\item Observed values at location $i$ are independent of values at location $j$.
\item The regression residuals associated observation $i$ are independent of those associated with observation $j$, i.e. $\mathbb{E}[e_{i}e_{j}]=\mathbb{E}[e_{i}]\mathbb{E}[e_{j}]=0$.
\end{itemize}

\begin{itemize}
\item Include the following literature: \citet{Bivand2008}, \citet{Bivand2008a}, \citet{Corrado2012}, \citet{Kelejian1998}, \citet{McMillen2010a}, \citet{Pinske2010}.
\end{itemize}




\section{Spatial point pattern analysis}
\section{Interpolation and geostatistics}
\section{Areal data and spatial autocorrelation}

\begin{subappendices}
\section{Code for Spatial Stats in R}
\lstinputlisting{"Programs/UP504_SpatStat1.R"}

\renewcommand\bibname{Further reading}

\bibliographystyle{econometrica}
\bibliography{Literature/OverLord}
\IfFileExists{Literature/SpatialStats.bbl}{\input{Literature/SpatialStats.bbl}}{\typeout{Need to run bibtex}}
\end{subappendices} }{\typeout{Need to run bibtex}}
\end{subappendices} }{\typeout{Need bibtex on SpatialStats}}
\end{subappendices}


% ── Part IV: Frontiers ────────────────────────────────────────────────
\part{Frontiers}\label{part4}
% =====================================================================
%  Chapter 10: Time Series and Panel Data
%  Source: TimeSeries.tex (Overleaf stub, expanded) + new scaffold
%  NOTE: Section 10.1-10.4 need full development
% =====================================================================

\chapter{Time Series and Panel Data}
\label{chap_TimeSeries}

\section{Introduction and Key Concepts}
\label{sec_TSIntro}

Time series analysis studies the properties of sequences of observations
indexed by time. The key distinction from cross-sectional analysis is that
observations are not independent---the value at time $t$ is typically
correlated with past values (autocorrelation) and the statistical properties
of the series may change over time (non-stationarity).

\noindent\textbf{Why time series matters for regional analysis.}
Economic time series at the metropolitan and county level---employment,
income, housing prices, unemployment rates---are the primary inputs for
regional monitoring, forecasting, and policy evaluation. Understanding their
time-series properties is a prerequisite for valid statistical inference.

\subsection{Stationarity}

A time series $\{y_t\}$ is \textbf{weakly stationary}\index{stationarity}
if its mean, variance, and autocovariance structure do not change over time:
\begin{align}
  \E[y_t] &= \mu \quad\text{(constant mean)} \\
  \Var[y_t] &= \sigma^2 \quad\text{(constant variance)} \\
  \Cov[y_t, y_{t-k}] &= \gamma_k \quad\text{(depends only on lag }k\text{, not }t\text{)}
\end{align}

\subsection{Autocorrelation}

The \textbf{autocorrelation function} (ACF)\index{autocorrelation function}
at lag $k$ is:
\begin{equation}
  \rho_k = \frac{\Cov[y_t, y_{t-k}]}{\Var[y_t]} = \frac{\gamma_k}{\gamma_0}
  \;\in\; [-1,+1]
\end{equation}

The \textbf{partial autocorrelation function} (PACF) measures the
correlation at lag $k$ after controlling for all shorter lags.


\section{AR and ARMA Processes}
\label{sec_ARMA}

\subsection{The AR(1) Model}

The simplest time series model:
\begin{equation}
  y_t = \phi y_{t-1} + e_t, \qquad e_t \sim \text{iid}(0,\sigma^2)
\end{equation}

If $|\phi| < 1$: \emph{stationary}---the series mean-reverts.
If $\phi = 1$: \emph{unit root}---the series is a random walk and is
non-stationary.

\subsection{AR($p$) and ARMA($p$,$q$) Models}

The AR($p$) generalisation:
\begin{equation}
  y_t = \phi_1 y_{t-1} + \phi_2 y_{t-2} + \cdots + \phi_p y_{t-p} + e_t
\end{equation}

Adding $q$ moving-average terms gives the ARMA($p$,$q$) model:
\begin{equation}
  y_t = \sum_{i=1}^{p}\phi_i y_{t-i} + e_t + \sum_{j=1}^{q}\theta_j e_{t-j}
\end{equation}

The Box-Jenkins methodology selects $p$ and $q$ using the ACF and PACF
plots:
\begin{itemize}
  \item ACF cuts off at lag $q$, PACF decays: suggests MA($q$)
  \item PACF cuts off at lag $p$, ACF decays: suggests AR($p$)
  \item Both decay gradually: suggests ARMA($p$,$q$)
\end{itemize}

% TODO: Add R code for ARIMA modelling of regional employment series


\section{Non-Stationarity and Unit Roots}
\label{sec_UnitRoot}

\subsection{Spurious Regression}

If two time series are independently generated random walks (unit root
processes), regressing one on the other will typically produce a
statistically significant coefficient with high $R^2$---even though there
is no causal relationship. This is \textbf{spurious
regression}\index{spurious regression} \citep{Granger1974}.

\subsection{Augmented Dickey-Fuller Test}

The \textbf{Augmented Dickey-Fuller (ADF) test}\index{Dickey-Fuller test}
tests the null hypothesis of a unit root:
\begin{equation}
  \Delta y_t = \alpha + \beta t + \delta y_{t-1}
               + \sum_{j=1}^{k}\gamma_j \Delta y_{t-j} + e_t
\end{equation}
$H_0$: $\delta = 0$ (unit root); $H_1$: $\delta < 0$ (stationary).

\subsection{Cointegration and Error Correction}

Two non-stationary series are \textbf{cointegrated}\index{cointegration} if
a linear combination of them is stationary. Cointegration implies a
long-run equilibrium relationship; deviations from it are corrected over time
via an error correction mechanism:
\begin{equation}
  \Delta y_t = \alpha(y_{t-1} - \beta x_{t-1}) + \Gamma\Delta x_t + e_t
\end{equation}

% TODO: Add cointegration example with regional house prices and income


\section{Forecasting}
\label{sec_Forecasting}

\subsection{In-Sample Fit vs.\ Out-of-Sample Accuracy}

\begin{itemize}
  \item \textbf{In-sample fit} ($R^2$, AIC, BIC): measures how well the
    model fits the estimation sample. Subject to overfitting.
  \item \textbf{Out-of-sample accuracy} (RMSE, MAE on holdout data):
    measures genuine predictive performance.
\end{itemize}

\subsection{Forecast Evaluation}

The \textbf{Diebold-Mariano test} \citep{Diebold1995} compares the
predictive accuracy of two forecasting models:
\begin{equation}
  DM = \frac{\bar{d}}{\hat\sigma_{\bar{d}}/\sqrt{T}} \sim \mathcal{N}(0,1)
\end{equation}
where $\bar{d}$ is the mean loss differential between models.

% TODO: Expand principles of forecasting, including FRED regional data


\section{Panel Data Methods}
\label{sec_PanelData_TS}

Panel data methods were introduced in Section~\ref{sec_Panel} in the context
of causal identification. This section extends the treatment to dynamic
panels and the connection to time-series properties.

\subsection{Dynamic Panels}

When the lagged dependent variable $y_{i,t-1}$ is included as a regressor,
standard fixed effects estimation is inconsistent because $y_{i,t-1}$ is
correlated with $\alpha_i$ in the demeaned equation. The
\textbf{Arellano-Bond estimator}\index{Arellano-Bond estimator}
\citep{Arellano1991} uses lagged levels of $y$ as instruments for
$\Delta y_{i,t-1}$ in a first-differenced GMM framework.

\subsection{Spatial Panels}

Combining cross-sectional spatial dependence with time series produces
\textbf{spatial panel models}\index{spatial panel models}:
\begin{equation}
  y_{it} = \rho \sum_j w_{ij}y_{jt} + X_{it}'\beta + \alpha_i + \lambda_t + e_{it}
\end{equation}

This model allows simultaneous control for: (i) spatial spillovers via the
SAR term; (ii) unit-level unobserved heterogeneity via $\alpha_i$; and (iii)
common time shocks via $\lambda_t$.

% TODO: Add worked example with FRED metropolitan employment data

\renewcommand\bibname{Further reading}
\begin{subappendices}
\bibliographystyle{econometrica}
\bibliography{Literature/OverLord}
\IfFileExists{Literature/TimeSeries.bbl}{\chapter{Time Series Analysis}\label{chap_TimeSeries}
\section{Introduction and key concepts}
\section{Autocorrelation and non-stationarity}
\section{Principles of forecasting}



\renewcommand\bibname{Further reading}
\begin{subappendices}
\bibliographystyle{econometrica}
\bibliography{Literature/OverLord}
\IfFileExists{Literature/TimeSeries.bbl}{\chapter{Time Series Analysis}\label{chap_TimeSeries}
\section{Introduction and key concepts}
\section{Autocorrelation and non-stationarity}
\section{Principles of forecasting}



\renewcommand\bibname{Further reading}
\begin{subappendices}
\bibliographystyle{econometrica}
\bibliography{Literature/OverLord}
\IfFileExists{Literature/TimeSeries.bbl}{\input{Literature/TimeSeries.bbl}}{\typeout{Need to run bibtex}}
\end{subappendices} }{\typeout{Need to run bibtex}}
\end{subappendices} }{\typeout{Need bibtex on TimeSeries}}
\end{subappendices}

% =====================================================================
%  Chapter 11: Computational Thinking and Reproducible Research
%  Source: CompMethods.tex (Overleaf) + P4 transcription (KDE) + new
% =====================================================================

\chapter{Computational Thinking and Reproducible Research}
\label{chap_Computational}

\section{The \textit{Caveat Computor} Imperative Revisited}
\label{sec_CavetComputorRevisit}

The Preface introduced the \emph{caveat computor} principle: computational
abundance makes conceptual rigour more rather than less important. This
chapter develops that principle into a practical framework for computational
regional analysis.

\noindent\textbf{What agentic coding assistants can and cannot do.}
AI coding assistants (GitHub Copilot, Claude, ChatGPT) can generate correct,
runnable code for standard analyses from plain-English descriptions. They can
suggest data sources, identify relevant packages, and debug syntax errors.
What they cannot do is determine whether the weight matrix choice is
defensible, whether the identification strategy is credible, or whether
the coefficient has a valid causal interpretation. These judgements require
the theoretical knowledge developed in Parts~I and~II.

The most dangerous analyst is not the one who cannot code---it is the one
who produces polished output without understanding its assumptions. The
obligation of this book is to ensure that the ability to generate spatial
regression code is matched by the ability to evaluate it.


\section{Reproducible Research in Practice}
\label{sec_Reproducible_ch11}

\subsection{Minimal Working Examples}

A minimal working example\index{minimal working example} (MWE) or
short, self-contained, correct example (SSCCE) is the gold standard for
communicating about code. When debugging or seeking help, an MWE:

\begin{itemize}
  \item Is \textbf{minimal}: the problem is isolated; nothing distracts from
    the error. Remove everything that is not essential to reproducing the issue.
  \item Is \textbf{working}: the code can be copied, pasted directly into a
    file, and executed without modification.
  \item Is \textbf{correct} (as far as possible): it correctly represents
    the intent of the analysis.
\end{itemize}

\noindent\textbf{Resources:}
\begin{itemize}
  \item \href{http://www.sscce.org}{sscce.org}: the authoritative reference
  \item \href{https://stackoverflow.com/help/minimal-reproducible-example}{%
    Stack Overflow guidelines}
\end{itemize}

\subsection{Literate Programming: Quarto}

Quarto (\texttt{quarto.org}) is a literate programming platform that
combines prose, code (R, Python, Julia), and mathematical notation in a
single source file, rendering to PDF, HTML, Word, or presentations. The
key workflow:

\begin{enumerate}
  \item Write code and analysis in a \texttt{.qmd} file
  \item \texttt{quarto render report.qmd} produces the output
  \item Commit \texttt{.qmd} + \texttt{renv.lock} + \texttt{requirements.txt}
    to GitHub
  \item Anyone can reproduce the analysis: \texttt{git clone} then
    \texttt{quarto render}
\end{enumerate}

\subsection{Version Control and Open Publication}

The \emph{American Economic Review} requires submission of all data and code
sufficient for replication before publication \citep{AER2024}. Urban
planning and geography journals have been slower to adopt such standards.
This book takes the position that reproducibility is a scientific norm, not
a courtesy: an analysis that cannot be replicated independently is not
science.


\section{Nonparametric Density Estimation}
\label{sec_KDE}

This section provides a technical introduction to kernel density estimation
(KDE), which arises naturally in spatial analysis (e.g.\ estimating the
density surface of events, such as crime or business locations) and in
distributional analysis (e.g.\ visualising the distribution of house prices
across neighbourhoods).

\subsection{The Histogram as a Rectangular KDE}

The histogram is the simplest nonparametric approach to density estimation.
It partitions the support of $x$ into non-overlapping bins of width $2h$
and counts observations in each bin. This is equivalent to a KDE with a
rectangular (uniform) kernel:
\begin{equation}
  K(x) = \begin{cases}
    \tfrac{1}{2} & \text{if } |x| < 1 \\
    0            & \text{otherwise}
  \end{cases}
\end{equation}
and bandwidth $h = \text{bin width}/2$.

\subsection{The General Kernel Density Estimator}

The kernel density estimator has two key ingredients:
\begin{enumerate}[label=(\roman*)]
  \item \textbf{Bandwidth} $h$: width of bin divided by 2; determines the
    size of the neighbourhood around evaluation point $x_0$.
  \item \textbf{Kernel} $K$: a weight-assigning or smoothing function
    that satisfies $\int K(x)\,dx = 1$.
\end{enumerate}

\begin{tcolorbox}[keyresult]
\begin{equation}
  \widehat{\text{pdf}}(x_0)
  = \frac{1}{Nh}\sum_{i=1}^{N} K\!\left(\frac{x_i - x_0}{h}\right)
  \label{eq_KDE}
\end{equation}
\end{tcolorbox}

\noindent Unlike the histogram, bins can \textbf{overlap}: most kernels
assign higher weight to observations near the center $x_0$ and let weights
decline with distance.

\subsection{Alternative Kernels}

\begin{figure}[htp]
\centering
\begin{tikzpicture}[>=Stealth, font=\small]
  \def\W{1.8}
  \def\H{1.0}
  \def\sep{3.8}

  \begin{scope}[xshift=0cm]
    \draw[->] (-\W-0.1,0) -- (\W+0.2,0);
    \draw[->] (0,-0.1) -- (0,\H+0.25);
    \draw[very thick] (-\W,0)--(-1,0)--(-1,\H)--(1,\H)--(1,0)--(\W,0);
    \node[below] at (0,-0.2) {Rectangular};
  \end{scope}

  \begin{scope}[xshift=\sep cm]
    \draw[->] (-\W-0.1,0) -- (\W+0.2,0);
    \draw[->] (0,-0.1) -- (0,\H+0.25);
    \draw[very thick, domain=-1:1, samples=40]
      plot (\x, {\H*(1-\x*\x)});
    \draw[very thick] (-\W,0)--(-1,0) (1,0)--(\W,0);
    \node[below] at (0,-0.2) {Epanechnikov};
  \end{scope}

  \begin{scope}[xshift=2*\sep cm]
    \draw[->] (-\W-0.1,0) -- (\W+0.2,0);
    \draw[->] (0,-0.1) -- (0,\H+0.25);
    \draw[very thick, domain=-1:1, samples=50]
      plot (\x, {\H*(1-\x*\x)^2});
    \draw[very thick] (-\W,0)--(-1,0) (1,0)--(\W,0);
    \node[below] at (0,-0.2) {Biweight};
  \end{scope}

  \begin{scope}[xshift=3*\sep cm]
    \draw[->] (-\W-0.1,0) -- (\W+0.2,0);
    \draw[->] (0,-0.1) -- (0,\H+0.25);
    \draw[very thick] (-1,0)--(0,\H)--(1,0);
    \draw[very thick] (-\W,0)--(-1,0) (1,0)--(\W,0);
    \node[below] at (0,-0.2) {Triangular};
  \end{scope}

  \begin{scope}[xshift=4*\sep cm]
    \draw[->] (-\W-0.1,0) -- (\W+0.2,0);
    \draw[->] (0,-0.1) -- (0,\H+0.25);
    \draw[very thick, domain=-1.8:1.8, samples=60]
      plot (\x, {\H*exp(-2.5*\x*\x)});
    \node[below] at (0,-0.2) {Gaussian};
  \end{scope}
\end{tikzpicture}
\caption{Five common kernel functions. All are centred at $x_0 = 0$ and
assign higher weight to nearby observations.}
\label{fig_Kernels}
\end{figure}

\subsection{The Gaussian KDE}

\begin{equation}
  \hat{f}(x) = \frac{1}{Nh}\sum_{i=1}^{N}\frac{1}{\sqrt{2\pi}}
               \exp\!\left(-0.5\left[\frac{x_i - x_0}{h}\right]^2\right)
  \label{eq_GaussianKDE}
\end{equation}

\subsection{Bandwidth Selection}

Choice of bandwidth is more important in practice than choice of kernel;
it involves a fundamental trade-off:
\begin{itemize}
  \item \textbf{Small $h$} (narrow bandwidth): low bias (correct shape)
    but high variance (jagged, imprecise estimate)
  \item \textbf{Large $h$} (wide bandwidth): low variance (smooth)
    but high bias (oversmooths the distribution)
\end{itemize}

The optimal bandwidth minimises asymptotic mean integrated squared error
(MISE). Silverman's rule-of-thumb provides a data-driven starting point:
$h = 1.06\hat\sigma n^{-1/5}$.


\section{Simulation and Resampling}
\label{sec_Simulation}

\subsection{The Bootstrap}

The bootstrap \citep{Efron1979} estimates the sampling distribution of a
statistic by resampling from the observed data with replacement. It provides
standard errors and confidence intervals without distributional assumptions,
making it particularly valuable when asymptotic approximations are unreliable.

For spatial data, the standard bootstrap must be modified to account for
spatial dependence---naive resampling destroys the spatial correlation
structure. The \textbf{spatial block bootstrap} resamples spatially
contiguous blocks of observations.

\subsection{Monte Carlo Methods for Spatial Inference}

Moran's $I$ significance is typically assessed by comparing the observed
value to a reference distribution generated by random permutation of the
data across spatial units:
\begin{enumerate}
  \item Compute observed $I^*$
  \item Randomly permute values across locations; compute $I$ for each
    permutation (e.g.\ 999 times)
  \item $p$-value = proportion of permuted $I$ values exceeding $I^*$
\end{enumerate}

This permutation approach makes no distributional assumptions and correctly
accounts for the specific topology of the weight matrix.


\section{Regional Agent-Based Models}
\label{sec_ABMs}

\subsection{What ABMs Add}

Agent-based models\index{agent-based models} (ABMs) simulate the behaviour
of individual agents (households, firms, municipalities) and allow macro-level
patterns to emerge from micro-level rules. Three phenomena are particularly
suited to ABM:

\begin{itemize}
  \item \textbf{Emergence:} Spatial patterns (segregation, clustering)
    arising from simple individual rules
  \item \textbf{Heterogeneity:} Distributional effects that aggregate models
    miss by construction
  \item \textbf{Path dependence:} Outcomes that depend on the sequence of
    historical events, not just current conditions
\end{itemize}

\subsection{The Schelling Segregation Model as a Spatial Equilibrium Problem}

\citet{Schelling1971} showed that mild individual preferences for neighbours
of the same type generate extreme aggregate segregation---a result that
anticipates the spatial equilibrium framework of Chapter~\ref{chap_SpatialEquil}
by showing how individual optimising behaviour produces spatial patterns that
no individual intended.

% TODO: Add NetLogo/nlrx worked example with calibration to ACS data


\section{Machine Learning and Causal Inference}
\label{sec_ML}

\subsection{What ML Offers: Prediction, Not Identification}

Machine learning methods (random forests, gradient boosting, LASSO, neural
networks) optimise predictive accuracy---they minimise out-of-sample
prediction error. This is valuable for several tasks in regional analysis:
variable selection, nonparametric approximation of complex response surfaces,
and out-of-sample forecasting.

What ML methods cannot do, without additional structure, is identify causal
effects. A predictive model that achieves high out-of-sample $R^2$ may do so
by exploiting spurious correlations in the training data. The fundamental
identification challenges of Chapter~\ref{chap_Causal}---selection bias,
geographic confounding, omitted variables---are not resolved by adding
flexibility to the model.

\subsection{Regularisation: LASSO and Ridge}

For high-dimensional settings with many candidate regressors, penalised
regression methods impose shrinkage on coefficients:
\begin{equation}
  \hat\beta_\text{LASSO} = \arg\min_\beta \Bigl[
    \sum(y_i - X_i'\beta)^2 + \lambda\|\beta\|_1
  \Bigr]
\end{equation}
\begin{equation}
  \hat\beta_\text{Ridge} = \arg\min_\beta \Bigl[
    \sum(y_i - X_i'\beta)^2 + \lambda\|\beta\|_2^2
  \Bigr]
\end{equation}

LASSO ($L_1$ penalty) produces sparse solutions (many coefficients set
exactly to zero) and serves as a variable selection tool. Ridge ($L_2$
penalty) shrinks all coefficients continuously toward zero without setting
any to exactly zero.

\subsection{Causal Forests and Heterogeneous Treatment Effects}

\citet{Wager2018} combine the random forest algorithm with the
potential-outcomes framework to estimate
\textbf{heterogeneous treatment effects}\index{heterogeneous treatment
effects}---the individual-level causal effect $\tau_i = Y_{1i} - Y_{0i}$---
as a function of covariates. The causal forest modifies the splitting rule of
a random forest to maximise treatment effect heterogeneity rather than
predictive accuracy, and uses orthogonalisation (the Robinson transformation)
to partial out confounders. This connects directly to the GWR framework:
where GWR asks whether the \emph{association} between $X$ and $Y$ varies
across space, a spatial causal forest asks whether the \emph{causal effect}
of a treatment varies across space.

\renewcommand\bibname{Further reading}
\begin{subappendices}
\bibliographystyle{econometrica}
\bibliography{Literature/OverLord}
\IfFileExists{Literature/CompMethods.bbl}{\chapter{Computational Methods}\label{chap_CompMethods}

\section{Principles of coding}
\subsection{Debugging code}
Most questions about code are made more answerable by the addition of a minimal working example (MWE)\index{MWE|see{minimal working example}}\index{minimal working example} or a short, self-contained, correct example (SSCCE).

Your debugging question may have an obvious answer, or it may not. When enlisting the help of others in debugging\index{code!debugging} your code, think about it from the point of view of the person whose help you're seeking. To ensure that the answer she provides is sound, she will probably want to create a sample document, make the changes she recommends, and check that it works. If she thinks she knows the answer but doesn't have the time to create a sample document, she may in the interest of quality control refrain from answering at this time. Thus you should always provide a block of code that can be copied, pasted directly into a file, compiled, tweaked, corrected, and pasted back.

In this context, ``minimal'' means that the problem is isolated and there is nothing in the sample that distracts from the error you have or the effect you desire. ``Working'' doesn't mean the problem is solved, it means the code sample can be copied-and-pasted and directly compiled. The site \href{http://www.sscce.org}{sscce.org} is dedicated to the topic. There are also several articles on the topic as it applies specifically to LaTeX:

\begin{itemize}
\item \href{http://theoval.cmp.uea.ac.uk/~nlct/latex/minexample/minexample.pdf}{``Creating a LaTeX Minimal Example''} (PDF) from Dr. Nicola Talbot's LaTeX resources,
\item \href{http://www.tex.ac.uk/cgi-bin/texfaq2html?label=minxampl}{``How to make a `minimum example'"} from the UK \TeX{} FAQ,
\item \href{http://www.minimalbeispiel.de/mini-en.html}{``What is a minimal working example?"} from the de.comp.text.tex newsgroup.
\end{itemize}


\section{High-performance computing}
This section introduces the concept of high-performance computing\index{high-performance computing} and illustrates an application using the University of Michigan's high-performance Linux cluster Flux\index{Flux|see{software}}\index{Flux|seealso{high-performance computing}}\index{software!Flux}.


\subsection{Parallelised statistical computations}
Many modern statistical techniques rely on computations that are parallel. Resampling methods (such as cross-validation and the bootstrap) and many simulations can be separated into independent sets of computations. 

As an example, we use simulation to produce a power curve for an independent samples $t$-test. After demonstrating the simulation using the \texttt{lapply} function in \textsf{R}, we parallelise the simulation using the \texttt{mcapply} function from the \texttt{multicore} package.



\subsection{University of Michigan high-performance cluster}
\subsubsection{Logging into the cluster and copying files}
\subsubsection{Software on Flux}
\subsubsection{Description of key commands}
\subsubsection{Running a batch job}

\section{Nonparametric econometrics}
This section with provide a basic introduction to nonparametric techniques in econometrics.

\linenumbers
\footnotesize
\begin{verbatim}
R> y <- plot(density(rnorm))
R> lines(x, add=TRUE)
\end{verbatim}
\normalsize
\nolinenumbers
}{\typeout{Need bibtex on CompMethods}}
\end{subappendices}


% ── Appendices ────────────────────────────────────────────────────────
\appendix
\part*{Appendices}
% =====================================================================
%  Appendix A: Getting Started with R and Python
%  Source: IntroR.tex (Overleaf) — adapted
% =====================================================================

\chapter{Getting Started with R and Python}
\label{appx_Software}

\section{R and RStudio}

\textsf{R}\index{software!R} is an open-source implementation of the
object-oriented mathematical programming language \textsf{S}, developed by
statisticians worldwide. In contrast to commercial packages such as
\textsf{SPSS}\index{software!SPSS} or \textsf{Stata}\index{software!Stata},
\textsf{R} is a full mathematical programming language with capabilities
comparable to \textsf{MATLAB}\index{software!MATLAB}, not merely a program
that performs canned regression routines. It is free, has no licensing fees,
and is increasingly used by businesses, government agencies, and non-profit
organisations.

Run \textsf{R} via the integrated development environment (IDE)
\href{http://rstudio.org/}{RStudio}\index{software!RStudio} or via
\href{https://code.visualstudio.com/}{VS Code} with the R extension.

\subsection{Core Spatial Libraries for R}

\begin{lstlisting}[language=R, caption={Installing core spatial libraries}]
install.packages(c(
  "sf",          # vector spatial data
  "terra",       # raster spatial data
  "tidycensus",  # Census / ACS data API
  "tigris",      # TIGER/Line geographies
  "spdep",       # spatial weights and Moran's I
  "spatialreg",  # SAR, SEM, IV estimation
  "GWmodel",     # GWR
  "leontief",    # input-output modelling
  "fredr",       # FRED time series API
  "osmdata",     # OpenStreetMap
  "ggplot2",     # graphics
  "tidyverse"    # data manipulation
))
\end{lstlisting}

\subsection{A Minimal Working Example: Spatial Regression in R}

\begin{lstlisting}[language=R,
  caption={Spatial lag model with spdep and spatialreg}]
library(sf)
library(spdep)
library(spatialreg)

# Load data (any sf object with a response and covariates)
data(columbus)                     # bundled example dataset
col_sf <- st_as_sf(columbus)

# Construct spatial weight matrix (queen contiguity)
W_nb  <- poly2nb(col_sf)           # neighbour list
W_lw  <- nb2listw(W_nb)            # list-weight object

# OLS (ignores spatial dependence -- baseline)
ols <- lm(CRIME ~ INC + HOVAL, data = col_sf)

# Test for spatial autocorrelation in OLS residuals
moran.test(residuals(ols), W_lw)

# SAR model (spatial lag)
sar <- lagsarlm(CRIME ~ INC + HOVAL,
                data = col_sf, listw = W_lw)
summary(sar)

# SEM model (spatial error)
sem <- errorsarlm(CRIME ~ INC + HOVAL,
                  data = col_sf, listw = W_lw)
summary(sem)
\end{lstlisting}

\section{Python and the Spatial Stack}

Python provides an equally capable spatial analytics environment via the
PySAL ecosystem \citep{Rey2010}.

\begin{lstlisting}[language=Python,
  caption={Installing core spatial libraries (Python)}]
pip install geopandas libpysal esda spreg mgwr
\end{lstlisting}

\begin{lstlisting}[language=Python,
  caption={Spatial lag model with spreg (Python equivalent)}]
import geopandas as gpd
import libpysal
from spreg import GM_Lag

# Load data
data = gpd.read_file(libpysal.examples.get_path("columbus.shp"))

# Construct weights
W = libpysal.weights.Queen.from_dataframe(data)
W.transform = "r"  # row-standardise

# SAR model
y  = data[["CRIME"]].values
X  = data[["INC","HOVAL"]].values
sar = GM_Lag(y, X, w=W, name_y="CRIME",
             name_x=["INC","HOVAL"])
print(sar.summary)
\end{lstlisting}

\section{Quarto: The Reproducible Workflow}

Quarto unifies prose, code, and mathematics in a single document:

\begin{verbatim}
---
title: "Spatial Regression Analysis"
format: html
---

## Data and Weights

```{r}
#| echo: true
library(spdep)
data(columbus)
W_nb <- poly2nb(st_as_sf(columbus))
W_lw <- nb2listw(W_nb)
```

The Moran's $I$ test for the OLS residuals yields ...
\end{verbatim}

Render with \texttt{quarto render document.qmd}. Toggle between PDF and HTML
output by changing \texttt{format: pdf} or \texttt{format: html}.

\section{Getting Started with QGIS}

\textsf{QGIS}\index{software!QGIS} (\href{http://www.qgis.org}{qgis.org})
is a free, open-source desktop GIS that integrates with R and Python. For
most analytical tasks in this book, R and Python are preferable because
they produce reproducible, scriptable workflows. QGIS is most useful for:
(a) initial exploratory visual analysis of spatial data; (b) digitising
and editing geometries; (c) presentation-quality cartographic output.

\renewcommand\bibname{Further reading}
\begin{subappendices}
\bibliographystyle{econometrica}
\bibliography{Literature/OverLord}
\IfFileExists{Literature/AppA.bbl}{\input{Literature/AppA.bbl}}{\typeout{Need bibtex on AppA}}
\end{subappendices}


% ── Back matter ───────────────────────────────────────────────────────
\backmatter
\printindex

\end{document}
